\documentclass[11pt, letterpaper]{article}
\usepackage[margin=1in]{geometry}
\usepackage{booktabs}
\usepackage{array}
\usepackage{tabularx}
\usepackage[table]{xcolor}
\usepackage{multirow}
\usepackage{enumitem}
\usepackage{titlesec}
\usepackage[hidelinks]{hyperref}
\usepackage{amsmath}
\usepackage{amssymb}
\usepackage[expansion=false]{microtype}
\usepackage{tcolorbox}
\usepackage{multicol}
\usepackage{titletoc}
\usepackage{graphicx}

\tcbuselibrary{skins, breakable}

% ── Colors ──────────────────────────────────────────────────────────────────
\definecolor{sectionblue}{RGB}{30, 90, 160}
\definecolor{sectiongreen}{RGB}{30, 130, 80}
\definecolor{sectionpurple}{RGB}{100, 40, 140}
\definecolor{sectionorange}{RGB}{180, 90, 20}
\definecolor{sectionteal}{RGB}{20, 120, 120}
\definecolor{tableheader}{RGB}{220, 230, 245}
\definecolor{tableheadergreen}{RGB}{210, 240, 220}
\definecolor{tableheaderorange}{RGB}{255, 235, 210}
\definecolor{tableheaderteal}{RGB}{210, 240, 240}
\definecolor{rowA}{RGB}{252, 252, 252}
\definecolor{rowB}{RGB}{237, 244, 255}
\definecolor{rowBg}{RGB}{237, 248, 241}
\definecolor{rowOr}{RGB}{255, 245, 232}
\definecolor{rowTe}{RGB}{232, 248, 248}
\definecolor{partbg}{RGB}{240, 244, 255}
\definecolor{part2bg}{RGB}{240, 252, 244}
\definecolor{part3bg}{RGB}{248, 240, 255}
\definecolor{part4bg}{RGB}{255, 245, 232}
\definecolor{part5bg}{RGB}{232, 248, 248}
\definecolor{craapbg}{RGB}{255, 252, 228}
\definecolor{craapborder}{RGB}{180, 148, 30}
\definecolor{tableheaderpurple}{RGB}{235, 220, 245}
\definecolor{rowPu}{RGB}{245, 235, 255}

% ── Section format switchers ─────────────────────────────────────────────────
\newcommand{\bluesections}{%
  \titleformat{\section}{\large\bfseries\color{sectionblue}}{}{0em}{}[\color{sectionblue}\titlerule]%
  \titleformat{\subsection}{\normalsize\bfseries\color{sectionblue!80!black}}{}{0em}{}%
}
\newcommand{\greensections}{%
  \titleformat{\section}{\large\bfseries\color{sectiongreen}}{}{0em}{}[\color{sectiongreen}\titlerule]%
  \titleformat{\subsection}{\normalsize\bfseries\color{sectiongreen!80!black}}{}{0em}{}%
}
\newcommand{\purplesections}{%
  \titleformat{\section}{\large\bfseries\color{sectionpurple}}{}{0em}{}[\color{sectionpurple}\titlerule]%
  \titleformat{\subsection}{\normalsize\bfseries\color{sectionpurple!80!black}}{}{0em}{}%
}
\newcommand{\orangesections}{%
  \titleformat{\section}{\large\bfseries\color{sectionorange}}{}{0em}{}[\color{sectionorange}\titlerule]%
  \titleformat{\subsection}{\normalsize\bfseries\color{sectionorange!80!black}}{}{0em}{}%
}
\newcommand{\tealsections}{%
  \titleformat{\section}{\large\bfseries\color{sectionteal}}{}{0em}{}[\color{sectionteal}\titlerule]%
  \titleformat{\subsection}{\normalsize\bfseries\color{sectionteal!80!black}}{}{0em}{}%
}

% ── Part banner ──────────────────────────────────────────────────────────────
\newcommand{\partbanner}[3]{%
  \begin{tcolorbox}[colback=#1, colframe=#2, boxrule=1.2pt, arc=4pt,
    left=8pt, right=8pt, top=6pt, bottom=6pt]
    \begin{center}{\LARGE\bfseries #3}\end{center}
  \end{tcolorbox}%
}

\setlength{\parskip}{4pt}
\setlength{\parindent}{0pt}
\newcommand{\yes}{\checkmark}
\newcommand{\divider}{\par\noindent\rule{\textwidth}{0.4pt}\par\vspace{2pt}}
\hbadness=10000
\hfuzz=5pt

% ── TOC styling ──────────────────────────────────────────────────────────────
\titlecontents{section}[1.5em]{\smallskip}
  {\contentslabel{1.5em}}{\hspace*{-1.5em}}
  {\titlerule*[6pt]{.}\contentspage}
\titlecontents{subsection}[3.5em]{}
  {\contentslabel{2em}}{\hspace*{-2em}}
  {\titlerule*[6pt]{.}\contentspage}

%% ===========================================================================
\begin{document}

% ── Title ────────────────────────────────────────────────────────────────────
\begin{center}
  {\Huge\bfseries Biology Notes}\\[6pt]
  {\large Introductory Course Reference}\\[4pt]
  \rule{\textwidth}{0.8pt}
\end{center}

\vspace{8pt}
\tableofcontents
\newpage

%% ===========================================================================
%%  PART I – Early Class / Exam 1
%% ===========================================================================
\greensections
\addcontentsline{toc}{section}{\textcolor{sectiongreen}{PART I \textemdash\ Early Class Notes (Exam 1)}}
\addcontentsline{toc}{subsection}{The Scientific Method}
\addcontentsline{toc}{subsection}{The Four Biological Macromolecules}
\partbanner{part2bg}{sectiongreen}{Part I \quad Early Class Notes \textnormal{\small(Exam 1)}}

\vspace{6pt}
\begin{center}\small\itshape
  Introductory concepts from the start of the course.
  Lower priority for later exams but essential foundational knowledge.
\end{center}
\vspace{4pt}

% ── Scientific Method ────────────────────────────────────────────────────────
\section*{The Scientific Method}

\subsection*{Observation}
The starting point of all science. You notice something in the natural world and form a question. Good observations are objective, specific, and repeatable.

\subsection*{Hypothesis}
A hypothesis is a testable, observation-based explanation stated as a \textbf{declarative statement}. A strong hypothesis:
\begin{enumerate}[noitemsep]
  \item Is \textbf{testable} -- can be supported or refuted through experimentation.
  \item Is \textbf{based on prior observations} -- grounded in existing knowledge.
  \item \textbf{Leads to predictions} -- allows you to say ``if \ldots then \ldots''
  \item Is phrased as a \textbf{declarative statement} (not a question).
\end{enumerate}

\subsection*{Prediction}
A prediction is a specific, measurable outcome that follows logically from the hypothesis. It bridges hypothesis and experiment.

\textbf{Format:} \emph{``If [hypothesis is true], then [observable result] when [experimental condition].''}

Predictions must be \textbf{concrete and falsifiable} -- they are what you actually measure or observe in an experiment. One hypothesis can generate many different predictions depending on what you choose to test.

\textit{Example --- Hypothesis: ``Plants grow taller with more sunlight.''\\
Prediction: ``If plants receive 12 hours of light rather than 6, then the 12-hour group will be measurably taller after 4 weeks under otherwise identical conditions.''}

\subsection*{Theory}
A theory has broader scope than a hypothesis and has survived rigorous, repeated testing:
\begin{enumerate}[noitemsep]
  \item \textbf{Rigorously and repeatedly tested} across many independent studies.
  \item Has \textbf{high predictive value} -- reliably forecasts new outcomes.
  \item \textbf{Can still be contradicted} -- not absolute, but very strongly supported.
\end{enumerate}
\textit{Note: In everyday language ``theory'' suggests uncertainty. In science, a theory is the strongest form of explanation we have.}

\subsection*{Scientific Law}
A law describes a consistent, observed pattern -- often with a \textbf{mathematical component}. It states \emph{what} happens, but does \textbf{not} explain \emph{why}. That is the role of a theory.

\vspace{4pt}
\textbf{Progression of confidence:}
\[
  \text{Hypothesis} \;\longrightarrow\; \text{Predictions Tested} \;\longrightarrow\; \text{Theory} \;\longrightarrow\; \text{Law}
\]

\vspace{6pt}

% ── CRAAP ────────────────────────────────────────────────────────────────────
\begin{tcolorbox}[colback=craapbg, colframe=craapborder,
  title={\bfseries The CRAAP Test -- Evaluating Sources},
  fonttitle=\normalsize, boxrule=1pt, arc=4pt]
Use this checklist to decide whether a source is credible:
\begin{itemize}[noitemsep, leftmargin=*]
  \item \textbf{C}urrency   -- How recent is the information? Is it current enough for your topic?
  \item \textbf{R}elevance  -- Does it address your question? Is the depth appropriate?
  \item \textbf{A}uthority  -- Who wrote it? What are their credentials? Is the publisher reputable?
  \item \textbf{A}ccuracy   -- Is it supported by evidence? Has it been peer-reviewed?
  \item \textbf{P}urpose    -- Why was it written? To inform, sell, or persuade?
\end{itemize}
\end{tcolorbox}

\vspace{8pt}

% ── Macromolecules ───────────────────────────────────────────────────────────
\section*{The Four Biological Macromolecules}

\noindent
{\small
\renewcommand{\arraystretch}{1.5}
\begin{tabular}{>{\bfseries\raggedright\arraybackslash}p{2.7cm}
                >{\raggedright\arraybackslash}p{3.3cm}
                >{\raggedright\arraybackslash}p{2.3cm}
                >{\raggedright\arraybackslash}p{6.1cm}}
\toprule
\rowcolor{tableheadergreen}
\multicolumn{1}{>{\bfseries\raggedright\arraybackslash}p{2.7cm}}{\textbf{Molecule}} &
\multicolumn{1}{>{\raggedright\arraybackslash}p{3.3cm}}{\textbf{Monomer}} &
\multicolumn{1}{>{\raggedright\arraybackslash}p{2.3cm}}{\textbf{Solubility}} &
\multicolumn{1}{>{\raggedright\arraybackslash}p{6.1cm}}{\textbf{Key Functions}} \\
\midrule
\rowcolor{rowA}
Carbohydrates
  & Monosaccharides (e.g.\ glucose)
  & Hydrophilic
  & Quick energy (ATP); cell-surface ID tags for signaling \\
\rowcolor{rowBg}
Lipids
  & Fatty acids \& glycerol (no true monomer)
  & Hydrophobic
  & Long-term energy storage; steroid hormones; phospholipid bilayer \\
\rowcolor{rowA}
Proteins
  & Amino acids (20 types)
  & Mostly hydrophilic
  & Enzymes; antibodies; motor proteins; receptors; peptide hormones \\
\rowcolor{rowBg}
Nucleic Acids
  & Nucleotides
  & Hydrophilic
  & DNA (genetic storage); RNA (transcription \& translation) \\
\bottomrule
\end{tabular}
}

\vspace{4pt}
\textit{Memory tip: \textbf{CLPN} -- Carbs, Lipids, Proteins, Nucleic Acids. All are polymers built from monomers except lipids, which have no true single monomer.}

\vspace{8pt}

% ── Properties of Water ───────────────────────────────────────────────────────
\section*{Properties of Water}

Water is the \textbf{medium of life} --- every cellular reaction occurs in an aqueous
environment. Its unique properties arise from its \textbf{polarity}: the oxygen atom pulls
electrons toward itself, creating a partial negative charge ($\delta^-$) on oxygen and
partial positive charges ($\delta^+$) on each hydrogen. This allows water molecules to form
\textbf{hydrogen bonds} with each other and with other polar molecules.

\noindent
{\renewcommand{\arraystretch}{1.6}
\begin{tabularx}{\textwidth}{>{\bfseries\raggedright\arraybackslash}p{4cm}
                              >{\raggedright\arraybackslash}X
                              >{\raggedright\arraybackslash}p{3.5cm}}
\toprule
\rowcolor{tableheadergreen}
\textbf{Property} & \textbf{Explanation} & \textbf{Biological Importance} \\
\midrule
\rowcolor{rowA} Cohesion & Water molecules hydrogen-bond to each other, creating surface tension. & Allows water to travel up plant xylem (transpiration pull); surface tension supports small organisms. \\
\rowcolor{rowBg} Adhesion & Water clings to other polar surfaces (e.g., glass, cellulose). & Capillary action moves water up narrow plant vessels. \\
\rowcolor{rowA} High Specific Heat & Takes a lot of energy to raise water temperature (4.18 J/g\textdegree C). & Buffers temperature changes; oceans and lakes moderate climate; maintains stable body temperature. \\
\rowcolor{rowBg} High Heat of Vaporization & Requires much energy for liquid water to become gas. & Evaporative cooling (sweating, transpiration) efficiently removes heat. \\
\rowcolor{rowA} Universal Solvent & Polar and ionic substances (``like dissolves like'') dissolve readily in water. & Nutrients, gases, and wastes dissolve for transport in blood and across membranes. \\
\rowcolor{rowBg} Ice Floats (Low Density as Solid) & Hydrogen bonds in ice form a crystalline lattice, making it less dense than liquid water. & Ice insulates lakes from above; aquatic life survives beneath frozen surface in winter. \\
\bottomrule
\end{tabularx}}

\vspace{4pt}
\begin{tcolorbox}[colback=part2bg, colframe=sectiongreen, arc=3pt, boxrule=1pt]
\textbf{pH Scale:} Water dissociates slightly: H$_2$O $\rightleftharpoons$ H$^+$ + OH$^-$.
The pH scale measures the concentration of H$^+$ ions on a \emph{logarithmic} scale from 0--14.
\begin{itemize}[noitemsep, leftmargin=*]
  \item \textbf{pH 7 = Neutral} (equal H$^+$ and OH$^-$, like pure water).
  \item \textbf{pH $<$ 7 = Acidic} (more H$^+$ than OH$^-$). Each unit = 10$\times$ more acidic.
  \item \textbf{pH $>$ 7 = Basic / Alkaline} (more OH$^-$ than H$^+$).
  \item \textbf{Acids} donate H$^+$ (protons). \textbf{Bases} accept H$^+$ or release OH$^-$.
  \item \textbf{Buffers} resist pH changes by absorbing or releasing H$^+$. Blood is buffered near pH 7.4.
\end{itemize}
\end{tcolorbox}

\newpage

%% ===========================================================================
%%  PART II – Chapter 3: Cells + Chemistry of Life
%% ===========================================================================
\bluesections
\addcontentsline{toc}{section}{\textcolor{sectionblue}{PART II \textemdash\ Cells \& Chemistry of Life}}
\addcontentsline{toc}{subsection}{Classification by Cell Type}
\addcontentsline{toc}{subsection}{Eukaryotic Organelles}
\addcontentsline{toc}{subsection}{Cell Features by Type}
\addcontentsline{toc}{subsection}{Serial Endosymbiosis Theory (SET)}
\addcontentsline{toc}{subsection}{Three Tenets of Cell Theory}
\addcontentsline{toc}{subsection}{Additional Key Points}
\addcontentsline{toc}{subsection}{CHNOPS -- Key Biological Elements}
\addcontentsline{toc}{subsection}{DNA vs. RNA \& Central Dogma}
\addcontentsline{toc}{subsection}{Transcription \& Translation}
\addcontentsline{toc}{subsection}{The Genetic Code}
\addcontentsline{toc}{subsection}{DNA Mutations}
\partbanner{partbg}{sectionblue}{Part II \quad Cells \& Chemistry of Life}

\vspace{8pt}

% ── Cell Classification ──────────────────────────────────────────────────────
\section*{Classification by Cell Type}

\noindent
{\setlength{\tabcolsep}{5pt}
\renewcommand{\arraystretch}{1.4}
\begin{tabularx}{\textwidth}{>{\raggedright\arraybackslash}Xcc}
\toprule
\rowcolor{tableheader}
\textbf{Statement} & \textbf{Prokaryotes} & \textbf{Eukaryotes} \\
\midrule
\rowcolor{rowA} Small (1--10 \textmu m)            & \yes &  \\
\rowcolor{rowB} Large (10--100 \textmu m)           &      & \yes \\
\rowcolor{rowA} Unicellular                         & \yes & \yes\ (some) \\
\rowcolor{rowB} Multicellular                       &      & \yes \\
\rowcolor{rowA} Reproduces via binary fission       & \yes &  \\
\rowcolor{rowB} Reproduces via mitosis / meiosis    &      & \yes \\
\rowcolor{rowA} No membrane-bound organelles        & \yes &  \\
\rowcolor{rowB} Membrane-bound organelles           &      & \yes \\
\rowcolor{rowA} Plasma membrane                     & \yes & \yes \\
\rowcolor{rowB} Ribosomes present                   & \yes & \yes \\
\rowcolor{rowA} Bacterial cell                      & \yes &  \\
\rowcolor{rowB} Plant cells                         &      & \yes \\
\rowcolor{rowA} Can have cell wall                  & \yes & \yes\ (plants, fungi) \\
\rowcolor{rowB} Linear DNA in nucleus               &      & \yes \\
\rowcolor{rowA} Nucleoid region                     & \yes &  \\
\rowcolor{rowB} Cytoplasm                           & \yes & \yes \\
\rowcolor{rowA} Older cell type                     & \yes &  \\
\rowcolor{rowB} DNA present                         & \yes & \yes \\
\bottomrule
\end{tabularx}}

% ── Organelles ───────────────────────────────────────────────────────────────
\section*{Eukaryotic Organelles}

\noindent
{\setlength{\tabcolsep}{5pt}
\renewcommand{\arraystretch}{1.4}
\begin{tabularx}{\textwidth}{>{\bfseries\raggedright\arraybackslash}p{3cm}>{\raggedright\arraybackslash}X}
\toprule
\rowcolor{tableheader}
\textbf{Organelle} & \textbf{Function(s)} \\
\midrule
\rowcolor{rowA} Nucleus         & Contains genetic material (DNA); master control of cell activity \\
\rowcolor{rowB} Nuclear Lamina  & Protein network that supports the nuclear envelope; maintains nucleus shape; helps organize DNA \\
\rowcolor{rowA} Nuclear Pore    & Channel in the nuclear envelope; allows substances (e.g., ribosomes) to exit the nucleus after assembly \\
\rowcolor{rowB} Rough ER        & Synthesizes and processes proteins destined for secretion or the membrane (ribosomes attached) \\
\rowcolor{rowA} Smooth ER       & Synthesizes lipids; detoxifies chemicals (abundant in liver cells); stores Ca$^{2+}$ ions \\
\rowcolor{rowB} Golgi Apparatus & Modifies, sorts, and packages proteins \& lipids for storage or transport; the ``shipping/receiving center'' of the cell \\
\rowcolor{rowA} Vesicle         & Transports materials within the cell \\
\rowcolor{rowB} Lysosome        & Breaks down waste materials and cellular debris \\
\rowcolor{rowA} Peroxisome      & Breaks down fatty acid wastes into hydrogen peroxide (H$_2$O$_2$), which is then converted to water \& CO$_2$; found in abundance in liver cells \\
\rowcolor{rowB} Mitochondria    & Produces ATP through cellular respiration \\
\rowcolor{rowA} Chloroplast     & Site of photosynthesis; solar energy $\to$ chemical energy (plants only) \\
\rowcolor{rowB} Nucleolus       & Produces ribosomes \\
\bottomrule
\end{tabularx}}

% ── Cell Features ────────────────────────────────────────────────────────────
\section*{Cell Features by Type}

\noindent
\begin{tabular}{@{} p{0.31\textwidth} p{0.31\textwidth} p{0.31\textwidth} @{}}
\textbf{Prokaryotic} & \textbf{Animal Cell} & \textbf{Plant Cell} \\ \addlinespace[0.5pt]
\vspace{-0.7em}
\begin{itemize}[noitemsep, topsep=0pt, leftmargin=*]
  \item Nucleoid region
  \item Ribosomes
  \item Cytoplasm
  \item Plasma membrane
  \item Cell wall
\end{itemize} &
\vspace{-0.7em}
\begin{itemize}[noitemsep, topsep=0pt, leftmargin=*]
  \item Nucleus \& Mitochondria
  \item Lysosomes, peroxisomes
  \item Plasma membrane
  \item \textbf{No} central vacuole
  \item \textbf{No} cell wall
\end{itemize} &
\vspace{-0.7em}
\begin{itemize}[noitemsep, topsep=0pt, leftmargin=*]
  \item Nucleus \& Mitochondria
  \item Chloroplasts
  \item Central vacuole
  \item Cell wall
  \item Plasmodesmata
\end{itemize} \\
\end{tabular}

\vspace{6pt}

\begin{tcolorbox}[colback=rowB, colframe=sectionblue, title={\bfseries The ``Big Three'' Differentiators}, arc=3pt, boxrule=1pt]
\begin{enumerate}[noitemsep]
  \item \textbf{Nucleus:} Prokaryotes have a \textbf{Nucleoid} region; Eukaryotes have a true \textbf{Nucleus} with a membrane.
  \item \textbf{Organelles:} Prokaryotes have \textbf{no} membrane-bound organelles; Eukaryotes have many (Mitochondria, ER, Golgi, etc.).
  \item \textbf{Size:} Prokaryotes are \textbf{small} (1--10 \textmu m); Eukaryotes are \textbf{large} (10--100 \textmu m).
\end{enumerate}
\divider
\textbf{What EVERY cell has:} Plasma membrane, cytoplasm, ribosomes, and DNA.
\end{tcolorbox}

% ── SET ─────────────────────────────────────────────────────────────────────
\section*{Serial Endosymbiosis Theory (SET)}

\textbf{Claim:} Mitochondria and chloroplasts were once free-living prokaryotes engulfed by a larger host cell. Over time they became permanent, interdependent organelles.

\vspace{4pt}\textbf{Evidence:}
\begin{itemize}[noitemsep]
  \item \textbf{Own DNA} -- Both carry circular DNA similar to bacteria.
  \item \textbf{Double membranes} -- Two membranes reflect the original engulfment event.
  \item \textbf{Ribosomes} -- Resemble prokaryotic ribosomes, not eukaryotic ones.
  \item \textbf{Reproduction} -- Replicate independently via binary fission inside the host.
  \item \textbf{Genetic similarity} -- DNA closely matches certain bacterial lineages.
\end{itemize}

% ── Cell Theory ──────────────────────────────────────────────────────────────
\section*{Three Tenets of Cell Theory}

\begin{enumerate}[noitemsep]
  \item All living things are composed of cells.
  \item Cells are the basic unit of structure and function in living organisms.
  \item All cells come from pre-existing cells.
\end{enumerate}

% ── Additional Key Points ────────────────────────────────────────────────────
\section*{Additional Key Points}

\subsection*{Criteria for Life --- The 8 Characteristics}

All living things share \textbf{eight fundamental characteristics}. An entity must satisfy
\emph{all} of them to be considered alive. A virus, for example, fails several criteria
(it cannot reproduce independently or carry out metabolism on its own).

\begin{tcolorbox}[colback=part2bg, colframe=sectiongreen, arc=3pt, boxrule=1pt]
\begin{enumerate}[noitemsep, leftmargin=*]
  \item \textbf{Organization} -- Living things have ordered, complex structure built from cells.
        Complexity exists at every level: atoms $\to$ molecules $\to$ organelles $\to$ cells $\to$
        tissues $\to$ organs $\to$ organ systems $\to$ organism.
  \item \textbf{Metabolism} -- Life carries out chemical reactions to acquire and use energy.
        Includes \emph{catabolism} (breaking down) and \emph{anabolism} (building up).
  \item \textbf{Growth and Development} -- Organisms increase in size and complexity according
        to a genetic program; cells grow, divide, and differentiate into specialized types.
  \item \textbf{Response to Stimuli} -- Living things detect and respond to changes in their
        environment (e.g., plants grow toward light; humans pull away from heat).
  \item \textbf{Homeostasis} -- Maintaining a stable internal environment despite external
        fluctuations (e.g., human body temperature regulated near 37\textdegree C).
  \item \textbf{Reproduction} -- Organisms produce offspring, passing genetic information to
        the next generation. Can be \emph{asexual} (binary fission, budding, cloning) or
        \emph{sexual} (gamete fusion).
  \item \textbf{Heredity} -- Traits are encoded in DNA and transmitted from parent to offspring,
        ensuring continuity across generations.
  \item \textbf{Evolutionary Adaptation} -- Populations change over time via natural selection;
        heritable traits that improve fitness become more common.
\end{enumerate}
\end{tcolorbox}
\textit{Memory: \textbf{OM-GR-RHE} -- Organization, Metabolism, Growth, Response,
Reproduction, Homeostasis, Heredity, Evolution. Or simply remember the list as the things
a rock cannot do.}

\subsection*{Surface Area to Volume Ratio}
\[
  \text{SA:V} = \frac{(\text{Surface Area})^2}{(\text{Volume})^3}
\]
\begin{itemize}[noitemsep]
  \item \textbf{Small cell} (high SA:V): efficient exchange with the environment.
  \item \textbf{Large cell} (low SA:V): surface area becomes a bottleneck; requires more energy and thermoregulation.
\end{itemize}

\newpage

% ── CHNOPS ──────────────────────────────────────────────────────────────────
\bluesections
\section*{CHNOPS -- Key Biological Elements}

\begin{tcolorbox}[colback=craapbg, colframe=craapborder,
  title={\bfseries The Octet Rule \& Stability}, arc=4pt, boxrule=1pt]
Atoms are most stable when their valence (outer) shell is full -- usually 8 electrons (the octet rule).
\begin{itemize}[noitemsep, leftmargin=*]
  \item \textbf{Reactive elements (C, H, N, O, P, S):} Have incomplete valence shells and form bonds to reach stability. These are the building blocks of all biological molecules.
  \item \textbf{Noble gases} (e.g.\ Neon, Helium): Have naturally full outer shells -- they are chemically inert and do not form biological molecules.
\end{itemize}
\end{tcolorbox}

\vspace{6pt}

\noindent
{\small
\renewcommand{\arraystretch}{1.5}
\begin{tabularx}{\textwidth}{>{\bfseries\color{sectionpurple}\raggedright\arraybackslash}p{2.8cm} >{\raggedright\arraybackslash}X}
\toprule
\rowcolor{tableheader} \textbf{Element} & \textbf{Biological Role} \\
\midrule
\rowcolor{rowA} Hydrogen (H) & Found in water and all organic molecules; powers ATP production via proton gradients in cellular respiration. \\
\rowcolor{rowBg} Carbon (C) & The backbone of life -- forms 4 covalent bonds, enabling the complex skeletons of carbs, lipids, proteins, and nucleic acids. \\
\rowcolor{rowA} Nitrogen (N) & Essential component of the amino group ($-$NH$_2$) in \textbf{amino acids/proteins} and the nitrogenous bases in \textbf{DNA/RNA}. \\
\rowcolor{rowBg} Oxygen (O) & Vital for aerobic cellular respiration; major component of water; found in most biological molecules. \\
\rowcolor{rowA} Phosphorus (P) & Forms the sugar-phosphate backbone of nucleic acids (DNA/RNA) and the hydrophilic head of phospholipids in cell membranes. \\
\rowcolor{rowBg} Sulfur (S) & Found in specific amino acids (cysteine, methionine); stabilizes protein 3D structure via disulfide bridges ($-$S$-$S$-$). \\
\bottomrule
\end{tabularx}}

\vspace{6pt}
\textit{Mnemonic: \textbf{CHNOPS} -- the six most abundant elements in living organisms, in order of biological importance.}

\vspace{6pt}
\textbf{Which macromolecules contain which elements?}
\begin{itemize}[noitemsep]
  \item \textbf{Carbohydrates \& Lipids:} Primarily C, H, and O.
  \item \textbf{Proteins:} C, H, N, O, and sometimes \textbf{S} (in cysteine/methionine).
  \item \textbf{Nucleic Acids (DNA \& RNA):} C, H, N, O, and \textbf{P} (phosphate backbone).
\end{itemize}

\newpage

% -- Microbiology section -----------------------------------------------------
\orangesections
\addcontentsline{toc}{subsection}{Viruses}
\addcontentsline{toc}{subsection}{Bacteria \& Gram Staining}
\addcontentsline{toc}{subsection}{Mutations}

\section*{Microbiology: Viruses, Bacteria \& Mutations}

\subsection*{Viruses}

A \textbf{virus} is a non-living infectious particle. It cannot reproduce or carry out
metabolism on its own --- it must hijack a host cell. Every virus consists of:
\begin{itemize}[noitemsep]
  \item A \textbf{nucleic acid genome} (DNA or RNA, single- or double-stranded)
  \item A \textbf{protein coat} called the \textbf{capsid}
  \item Sometimes an outer \textbf{lipid envelope} (stolen from a previous host's membrane)
  \item \textbf{Surface glycoproteins} (spikes) that recognize and bind to specific host receptors
\end{itemize}

\textbf{Why viruses aren't ``alive'':} no metabolism, no growth, no homeostasis, cannot
reproduce without a host. They sit at the boundary of life --- biologically active only
when inside a cell.

\textbf{Replication strategy (general):}
\begin{enumerate}[noitemsep]
  \item \textbf{Attachment} -- spike protein binds to a specific host cell receptor.
  \item \textbf{Entry} -- viral genome enters the cell (membrane fusion, endocytosis, or injection).
  \item \textbf{Replication} -- the host's machinery copies the viral genome.
  \item \textbf{Assembly} -- new viral particles are built from synthesized parts.
  \item \textbf{Release} -- new viruses exit (lysis kills the cell, or budding releases them gradually).
\end{enumerate}

\subsection*{Segmented Genomes \& Influenza}

Some viruses (notably \textbf{influenza A}) have a \textbf{segmented genome} --- the genetic
material is split across multiple separate RNA strands. Influenza A has \textbf{8 segments}.

\textbf{Why segmentation matters:} when two different flu strains infect the same cell at
the same time, their segments can mix and match during assembly. This is called
\textbf{reassortment} and produces brand-new combinations the immune system has never
seen --- a major source of pandemic flu strains.

\subsection*{Antigenic Shift vs.\ Antigenic Drift}

These are the two ways flu viruses (and others) evade the immune system. Both rely on
changing the \textbf{surface proteins} (antigens) the body uses to recognize the virus.

\noindent
{\renewcommand{\arraystretch}{1.6}
\begin{tabularx}{\textwidth}{>{\bfseries\raggedright\arraybackslash}p{3.5cm}
                              >{\raggedright\arraybackslash}X}
\toprule
\rowcolor{tableheaderorange}
\textbf{Mechanism} & \textbf{What Happens} \\
\midrule
\rowcolor{rowOr} \textbf{Antigenic Drift} & \textbf{Gradual} change in surface proteins (hemagglutinin and neuraminidase, the H and N in ``H1N1'') from accumulated point mutations during normal replication. Small changes year to year. The reason \textbf{seasonal flu shots} need updating annually --- last year's antibodies recognize last year's virus, but not the slightly drifted current version. \\
\rowcolor{rowA}  \textbf{Antigenic Shift} & \textbf{Sudden, dramatic} change when two strains co-infect a host (often a pig or bird) and \textbf{reassort their segments} into a novel combination. The new virus has surface proteins the human immune system has never seen. This is what causes \textbf{flu pandemics} (1918 H1N1, 2009 H1N1, etc.). Only segmented viruses can shift. \\
\bottomrule
\end{tabularx}}

\textit{Quick rule: \textbf{drift} = small, slow, every year, requires only mutation. \textbf{Shift} = big, sudden, rare, requires reassortment of segments from different strains.}

\subsection*{Why We Need a New Flu Shot Every Year}

The flu virus continually changes its surface proteins through antigenic drift. Antibodies
from last year's vaccine recognize last year's H and N proteins, but not the drifted current
versions. Each year, public health agencies (WHO, CDC) predict which strains will dominate
and create a new vaccine targeting those strains. If their prediction is wrong, vaccine
effectiveness drops --- but it's still usually partially protective due to cross-reactivity.

\subsection*{Bacteria}

\textbf{Bacteria} are single-celled \textbf{prokaryotes} (no nucleus, no membrane-bound
organelles). They are alive, ubiquitous, and metabolically incredibly diverse.

\textbf{Bacterial cell components:}
\begin{itemize}[noitemsep]
  \item \textbf{Cell wall} (made of peptidoglycan in bacteria; rigid; protects against bursting)
  \item \textbf{Plasma membrane} (phospholipid bilayer)
  \item \textbf{Cytoplasm} containing ribosomes (70S, smaller than eukaryotic 80S)
  \item \textbf{Nucleoid} -- a region with the single circular chromosome (no nuclear membrane)
  \item \textbf{Plasmids} -- small extra circles of DNA, often carrying antibiotic resistance genes
  \item \textbf{Optional structures}: capsule, fimbriae, pili, flagella, endospores (see below)
\end{itemize}

\textbf{Reproduction:} \textbf{binary fission} --- a single cell copies its DNA, grows, and
splits into two identical daughter cells. Fast: \textit{E.\ coli} can divide every 20 minutes
in good conditions.

\textbf{How bacteria cause disease:}
\begin{itemize}[noitemsep]
  \item Producing \textbf{toxins} (e.g., botulinum toxin from \textit{Clostridium botulinum}; tetanus toxin)
  \item \textbf{Direct tissue damage} from massive replication (e.g., strep throat)
  \item Triggering an \textbf{immune overreaction} (sepsis, toxic shock syndrome)
  \item \textbf{Biofilm} formation (sticky communities resistant to antibiotics, e.g., on catheters)
\end{itemize}

\subsection*{Gram Staining}

\textbf{Gram staining} is a lab technique that classifies bacteria into two groups based
on their cell wall structure. Developed by Hans Christian Gram in 1884; still used today
because it tells you immediately what kind of antibiotic might work.

\noindent
{\renewcommand{\arraystretch}{1.6}
\begin{tabularx}{\textwidth}{>{\bfseries\raggedright\arraybackslash}p{3cm}
                              >{\raggedright\arraybackslash}X
                              >{\raggedright\arraybackslash}X}
\toprule
\rowcolor{tableheaderorange}
\textbf{Feature} & \textbf{Gram-Positive} & \textbf{Gram-Negative} \\
\midrule
\rowcolor{rowOr} Stain color    & \textbf{Purple/violet} (retains crystal violet) & \textbf{Pink/red} (loses crystal violet, takes up safranin counter-stain) \\
\rowcolor{rowA}  Cell wall      & \textbf{Thick} peptidoglycan layer (the dye gets trapped) & \textbf{Thin} peptidoglycan + outer membrane (the dye washes out) \\
\rowcolor{rowOr} Outer membrane & \textbf{None}                                    & \textbf{Yes} -- contains LPS (lipopolysaccharide), an endotoxin \\
\rowcolor{rowA}  Antibiotic susceptibility & Often more susceptible to penicillin (targets peptidoglycan) & The outer membrane blocks many antibiotics; harder to treat \\
\rowcolor{rowOr} Examples       & \textit{Staphylococcus, Streptococcus, Bacillus, Clostridium} & \textit{E.\ coli, Salmonella, Pseudomonas, Neisseria, Helicobacter pylori} \\
\bottomrule
\end{tabularx}}

\textbf{Why it matters clinically:} If a doctor takes a sample, sees Gram-positive cocci
in clusters, they can immediately suspect Staphylococcus and prescribe accordingly --- without
waiting days for a full culture.

\subsection*{Capsule, Fimbriae \& Endospore}

These are optional bacterial structures that increase survival or virulence.

\noindent
{\renewcommand{\arraystretch}{1.6}
\begin{tabularx}{\textwidth}{>{\bfseries\raggedright\arraybackslash}p{2.8cm}
                              >{\raggedright\arraybackslash}X}
\toprule
\rowcolor{tableheaderorange}
\textbf{Structure} & \textbf{Function} \\
\midrule
\rowcolor{rowOr} \textbf{Capsule} & A slippery polysaccharide layer outside the cell wall. \textbf{Protects from immune cells (phagocytes can't grip)}, prevents drying out. Major virulence factor in \textit{Streptococcus pneumoniae} and \textit{Klebsiella}. \\
\rowcolor{rowA}  \textbf{Fimbriae (pili)} & Short, hair-like surface projections used to \textbf{adhere to surfaces or host cells}. Without them, many pathogens can't establish infection. \textbf{Sex pili} (a special longer type) transfer DNA between bacteria during conjugation --- a major route for spreading antibiotic resistance. \\
\rowcolor{rowOr} \textbf{Flagella} & Long whip-like structures for motility. Driven by a rotating molecular motor at the base. \\
\rowcolor{rowA}  \textbf{Endospore} & A dormant, highly resistant survival form some bacteria produce when conditions are bad. The cell shrinks its DNA inside a tough shell. \textbf{Endospores can survive boiling, drying, radiation, and centuries of dormancy}, then reactivate when conditions improve. Made by \textit{Clostridium} (botulism, tetanus, \textit{C.\ difficile}) and \textit{Bacillus} (anthrax). This is why simple boiling does not sterilize --- you need an autoclave (121\textdegree C under pressure). \\
\bottomrule
\end{tabularx}}

\subsection*{Nitrogen Fixation}

Atmospheric nitrogen (\textbf{N\textsubscript{2}}, $\sim$78\% of air) is incredibly stable
because of its triple bond. Plants and animals \emph{cannot} use N\textsubscript{2} directly
--- they need nitrogen in usable forms like ammonia (NH\textsubscript{3}) or nitrate
(NO\textsubscript{3}\textsuperscript{$-$}). \textbf{Only certain bacteria can break the
N\textsubscript{2} triple bond} and convert it into usable forms. This process is called
\textbf{nitrogen fixation}.

\textbf{Who does it:}
\begin{itemize}[noitemsep]
  \item \textbf{Free-living soil bacteria} like \textit{Azotobacter}.
  \item \textbf{Symbiotic bacteria} like \textit{Rhizobium}, which live in nodules on the roots of legumes (peas, beans, clover, alfalfa). The plant feeds the bacteria sugars; the bacteria feed the plant fixed nitrogen.
  \item \textbf{Cyanobacteria} in aquatic environments.
\end{itemize}

\textbf{Why it matters:} all proteins, all DNA, all amino acids contain nitrogen. Without
nitrogen-fixing bacteria, life as we know it would not exist. Industrial nitrogen fixation
(the \textbf{Haber-Bosch process}, used to make synthetic fertilizer) is one of the most
important inventions of the 20th century --- it currently feeds about half the world's
population.

\subsection*{Bacterial Oxygen Tolerance Categories}

Different bacteria have different relationships with oxygen. This determines where they
can live and how they generate ATP.

\noindent
{\renewcommand{\arraystretch}{1.6}
\begin{tabularx}{\textwidth}{>{\bfseries\raggedright\arraybackslash}p{4cm}
                              >{\raggedright\arraybackslash}X
                              >{\raggedright\arraybackslash}p{3cm}}
\toprule
\rowcolor{tableheaderorange}
\textbf{Type} & \textbf{Relationship with O\textsubscript{2}} & \textbf{Examples} \\
\midrule
\rowcolor{rowOr} \textbf{Obligate Aerobe} & \textbf{Requires} oxygen to survive. Uses oxygen as the final electron acceptor for cellular respiration. Cannot live in oxygen-free environments. & \textit{Mycobacterium tuberculosis} \\
\rowcolor{rowA}  \textbf{Obligate Anaerobe} & \textbf{Cannot tolerate} oxygen --- it is toxic to them (they lack catalase or superoxide dismutase). Lives in deep soil, intestines, deep wounds. & \textit{Clostridium tetani, Clostridium botulinum} \\
\rowcolor{rowOr} \textbf{Facultative Anaerobe} & \textbf{Prefers oxygen} (more ATP), but can switch to fermentation or anaerobic respiration when oxygen is absent. Most flexible. & \textit{E.\ coli, Staphylococcus, Salmonella} \\
\rowcolor{rowA}  \textbf{Aerotolerant Anaerobe} & Doesn't \emph{use} oxygen, but can tolerate it. Always uses fermentation. & \textit{Lactobacillus} (yogurt cultures) \\
\rowcolor{rowOr} \textbf{Microaerophile} & Needs \emph{some} oxygen, but only at low concentrations (full atmospheric levels are toxic). & \textit{Helicobacter pylori} (causes ulcers) \\
\bottomrule
\end{tabularx}}

\textbf{Memory aid:} \textbf{Obligate} = ``stuck with one option,'' \textbf{Facultative} = ``has a choice.''

\subsection*{Mutations}

A \textbf{mutation} is any change in the DNA sequence of a cell. Mutations are the
\textbf{ultimate source of all genetic variation} --- without them, evolution could not occur
because there would be no new alleles for natural selection to act on.

\textbf{What causes mutations:}
\begin{itemize}[noitemsep]
  \item \textbf{Replication errors} -- DNA polymerase makes occasional mistakes ($\sim$1 in 10\textsuperscript{9} bases after proofreading)
  \item \textbf{Mutagens} -- chemicals (e.g., benzene, certain food preservatives), UV light, ionizing radiation
  \item \textbf{Transposable elements} -- ``jumping genes'' that insert into other genes
  \item \textbf{Spontaneous chemical changes} -- bases can rearrange (tautomerize) or be deaminated
\end{itemize}

\textbf{Types of point mutations} (single-base changes):

\noindent
{\renewcommand{\arraystretch}{1.6}
\begin{tabularx}{\textwidth}{>{\bfseries\raggedright\arraybackslash}p{3cm}
                              >{\raggedright\arraybackslash}X
                              >{\raggedright\arraybackslash}p{3cm}}
\toprule
\rowcolor{tableheaderorange}
\textbf{Type} & \textbf{Effect on Protein} & \textbf{Example} \\
\midrule
\rowcolor{rowOr} \textbf{Silent}    & \textbf{No change} in amino acid (genetic code redundancy --- multiple codons can code for the same amino acid). Protein is unchanged. & GCU $\to$ GCC, both = Alanine \\
\rowcolor{rowA}  \textbf{Missense}  & A different amino acid is inserted. Effect varies from negligible (similar amino acid) to catastrophic. & Sickle cell anemia (1 amino acid change, dramatic effect) \\
\rowcolor{rowOr} \textbf{Nonsense}  & Premature \textbf{stop codon} created --- protein is truncated, usually nonfunctional. & Many genetic diseases (e.g., Duchenne muscular dystrophy) \\
\rowcolor{rowA}  \textbf{Frameshift}& Insertion or deletion of a base (not a multiple of 3) shifts the entire reading frame downstream. Almost always disastrous --- nearly every codon afterward is wrong. & Cystic fibrosis ($\Delta$F508 deletion) \\
\bottomrule
\end{tabularx}}

\textbf{Effects of mutations:}
\begin{itemize}[noitemsep]
  \item \textbf{Most are neutral} -- in non-coding regions or silent.
  \item \textbf{Some are harmful} -- cause genetic diseases, cancer (when in growth-control genes).
  \item \textbf{Rarely beneficial} -- the raw material of evolution. A beneficial mutation in a population can spread by natural selection (e.g., lactose tolerance in adults, antibiotic resistance in bacteria).
  \item \textbf{Germline vs.\ somatic} -- germline mutations (in sperm/egg) are passed to offspring; somatic mutations (in body cells) are not heritable but can cause cancer.
\end{itemize}




% ── DNA vs RNA / Central Dogma ───────────────────────────────────────────────
\orangesections
\section*{DNA vs.\ RNA \& the Central Dogma}

\subsection*{Why DNA for Long-Term Storage?}
DNA is used instead of RNA for long-term genetic storage because it is significantly more chemically stable:
\begin{itemize}[noitemsep]
  \item DNA uses \textbf{deoxyribose} sugar, which lacks the reactive \textbf{2'-hydroxyl group} (--OH) present in RNA's ribose. This makes DNA far less susceptible to hydrolysis.
  \item DNA is \textbf{double-stranded}, providing a built-in backup copy and allowing damage repair mechanisms to use the complementary strand as a template.
  \item RNA is intentionally short-lived -- this allows the cell to rapidly change which proteins are being made in response to the environment.
\end{itemize}

\subsection*{DNA Strand Directionality \& Chargaff's Rules}
DNA is always read and written in a specific direction. The two strands run \textbf{antiparallel} -- the 5'$\to$3' direction of one strand runs counter to the 5'$\to$3' direction of the other:

\begin{tcolorbox}[colback=part4bg, colframe=sectionorange, arc=3pt, boxrule=1pt]
\ttfamily\small
\begin{tabbing}
Template strand: \quad \= 3\textquoteleft\ \ AAt ACA CAT TAC ATT\ \ 5\textquoteleft \\
Complementary:         \> 5\textquoteleft\ \ TTa TGT GTA ATG TAA\ \ 3\textquoteleft
\end{tabbing}
\normalfont
\vspace{2pt}
\footnotesize Note: A pairs with T (2 H-bonds); C pairs with G (3 H-bonds). Synthesis always proceeds \textbf{5'$\to$3'}.
\end{tcolorbox}

\begin{tcolorbox}[colback=craapbg, colframe=craapborder,
  title={\bfseries Chargaff's Rules -- Base Pairing Percentages}, arc=4pt, boxrule=1pt]
Because A always pairs with T, and C always pairs with G:
\begin{itemize}[noitemsep, leftmargin=*]
  \item \textbf{\%A = \%T} and \textbf{\%C = \%G} in any double-stranded DNA molecule.
  \item \textbf{\%A + \%T + \%C + \%G = 100\%}
  \item If you know one base's percentage, you can calculate all the others.
\end{itemize}
\textit{Example: If C = 34\%, then G = 34\%, leaving A + T = 32\%, so A = T = 16\%.}
\end{tcolorbox}

\subsection*{The Central Dogma of Molecular Biology}
\[
  \text{DNA} \;\longrightarrow\; \underbrace{\text{Pre-mRNA} \;\longrightarrow\; \text{mRNA}}_{\text{Transcription + Processing}} \;\longrightarrow\; \underbrace{\text{Protein}}_{\text{Translation}}
\]

\begin{itemize}[noitemsep]
  \item \textbf{Transcription} = RNA synthesis from a DNA template (in the nucleus).
  \item \textbf{Translation} = protein synthesis from an mRNA template (at ribosomes).
  \item DNA is never directly translated into protein -- RNA is the essential intermediary.
\end{itemize}

% ── Transcription & Translation ──────────────────────────────────────────────
\section*{Transcription \& mRNA Processing}

\subsection*{Identifying the Template Strand}
The \textbf{template strand} is the strand of DNA that RNA polymerase reads to synthesize mRNA. The resulting mRNA sequence is \textbf{complementary and antiparallel} to the template strand (and identical in sequence to the non-template/coding strand, but with U replacing T).

\begin{tcolorbox}[colback=part4bg, colframe=sectionorange, arc=3pt, boxrule=1pt]
\textbf{Worked example:} Given the coding strand 5'-GTCTATGCATTA-3', find the mRNA.
\begin{enumerate}[noitemsep]
  \item \textbf{Template strand} (antiparallel complement): 3'-CAGATACGTAAT-5'
  \item \textbf{mRNA} (same as coding strand, U replaces T): 5'-GUCUAUGCAUUA-3'
\end{enumerate}
\textit{The mRNA is read 5'$\to$3' starting from the AUG start codon.}
\end{tcolorbox}

\subsection*{Stages of Transcription}
\begin{enumerate}[noitemsep]
  \item \textbf{Initiation} -- RNA polymerase binds to the promoter region (e.g., TATAAA box) and unwinds the DNA. Transcription begins at the Transcription Start Site (TSS), \emph{after} the promoter.
  \item \textbf{Elongation} -- RNA polymerase moves along the template strand, synthesizing pre-mRNA (5'$\to$3').
  \item \textbf{Termination} -- Transcription stops at a terminator sequence; the pre-mRNA is released.
\end{enumerate}

\subsection*{mRNA Processing (Eukaryotes only)}
Before the pre-mRNA can leave the nucleus, it must be processed into mature mRNA:

\begin{tcolorbox}[colback=part2bg, colframe=sectiongreen, arc=3pt, boxrule=1pt]
\begin{itemize}[noitemsep, leftmargin=*]
  \item \textbf{5' Cap} -- A modified guanine nucleotide is added to the 5' end. Protects mRNA from degradation and helps ribosomes recognize it for translation.
  \item \textbf{3' Poly-A Tail} -- A string of $\sim$200 adenine nucleotides added to the 3' end. Increases stability and aids export from the nucleus.
  \item \textbf{RNA Splicing} -- \textbf{Introns} (non-coding intervening sequences) are removed; \textbf{exons} (expressed sequences) are joined together to form the final coding sequence.
\end{itemize}
\end{tcolorbox}

\subsection*{The Spliceosome}
RNA splicing is carried out by the \textbf{spliceosome}, a large complex of:
\begin{itemize}[noitemsep]
  \item \textbf{snRNPs} (small nuclear ribonucleoproteins, pronounced ``snurps'') -- recognize splice sites on the pre-mRNA.
  \item Additional proteins that catalyze the cutting and rejoining reactions.
\end{itemize}
The spliceosome removes introns and joins the remaining exons, producing the mature mRNA that is exported to the cytoplasm for translation.

\vspace{4pt}
\textit{Memory: \textbf{Introns} are \textbf{in}tervening sequences that stay \textbf{in} the nucleus and are discarded. \textbf{Exons} are \textbf{ex}pressed -- they exit and get translated.}

\subsection*{Translation Overview}
\begin{enumerate}[noitemsep]
  \item \textbf{Initiation} -- Ribosome assembles at the start codon (\textbf{AUG}, codes for methionine).
  \item \textbf{Elongation} -- tRNA molecules bring amino acids; ribosome reads codons and links amino acids into a polypeptide chain.
  \item \textbf{Termination} -- Ribosome reaches a stop codon (UAA, UAG, or UGA); polypeptide is released.
\end{enumerate}

\textbf{tRNA} is ``charged'' (loaded with its amino acid) by enzymes called \textbf{aminoacyl-tRNA synthetases}. Each tRNA has an \textbf{anticodon} that base-pairs with the complementary mRNA codon.

\begin{tcolorbox}[colback=part4bg, colframe=sectionorange, arc=3pt, boxrule=1pt]
\textbf{tRNA wobble:} tRNAs do \emph{not} always correctly read the 3rd base of a codon (the wobble position). A single tRNA can recognize multiple codons that differ only in the 3rd position. This is consistent with codon redundancy in the genetic code.
\end{tcolorbox}

% ── Genetic Code ─────────────────────────────────────────────────────────────
\section*{The Genetic Code}

The genetic code uses \textbf{codons} -- triplets of mRNA bases -- to specify amino acids. With 4 bases (A, U/T, C, G) in triplets:
\[
  4^2 = 16 \text{ (pairs -- not enough for 20 amino acids)} \qquad 4^3 = 64 \text{ (triplets -- more than enough)}
\]

\subsection*{Three Key Properties of the Genetic Code}

\begin{tcolorbox}[colback=part4bg, colframe=sectionorange, arc=3pt, boxrule=1pt]
\begin{itemize}[noitemsep, leftmargin=*]
  \item \textbf{Redundant (Degenerate):} Multiple codons can specify the \emph{same} amino acid. For example, both UUU and UUC code for phenylalanine. Changes in the \textbf{third base} of a codon often do not change the amino acid (called a ``wobble'' position). This provides tolerance against certain mutations.
  \item \textbf{Non-Ambiguous:} Each codon specifies \emph{only one} amino acid -- never two. The codon AUG always means methionine. This ensures precise, reliable protein synthesis.
  \item \textbf{Universal:} The same codons specify the same amino acids in \emph{virtually all} organisms -- from bacteria to humans. For example, UUU codes for phenylalanine in humans, bacteria, and plants. This universality is strong evidence that all life shares a common ancestor.
\end{itemize}
\end{tcolorbox}

\vspace{4pt}
\textbf{Special codons:}
\begin{itemize}[noitemsep]
  \item \textbf{AUG} -- Start codon; codes for methionine; signals where translation begins.
  \item \textbf{UAA, UAG, UGA} -- Stop codons; do not code for any amino acid; signal the end of translation.
\end{itemize}

% ── DNA Mutations ─────────────────────────────────────────────────────────────
\section*{DNA Mutations}

A \textbf{mutation} is any change in the DNA sequence. Point mutations affect a single nucleotide.

\noindent
{\setlength{\tabcolsep}{5pt}
\renewcommand{\arraystretch}{1.5}
\begin{tabularx}{\textwidth}{>{\bfseries\raggedright\arraybackslash}p{3.2cm}>{\raggedright\arraybackslash}X>{\raggedright\arraybackslash}p{4.5cm}}
\toprule
\rowcolor{tableheader}
\textbf{Mutation Type} & \textbf{Description} & \textbf{Example} \\
\midrule
\rowcolor{rowA}
Silent (Synonymous)
  & The base change does \emph{not} alter the amino acid produced, due to codon redundancy. Protein is unchanged.
  & UGU $\to$ UGC (both code for cysteine) \\
\rowcolor{rowB}
Missense
  & The base change results in a \emph{different} amino acid being incorporated. Protein function may be altered.
  & UGU (Cys) $\to$ UGG (Trp) \\
\rowcolor{rowA}
Nonsense
  & The base change creates a premature \emph{stop codon}. Translation terminates early, producing a truncated (usually nonfunctional) protein.
  & UGG (Trp) $\to$ UGA (Stop) \\
\bottomrule
\end{tabularx}}

\vspace{6pt}
\begin{tcolorbox}[colback=rowB, colframe=sectionblue, arc=3pt, boxrule=1pt]
\textbf{Key insight:} Silent mutations are possible because the genetic code is \textbf{redundant} -- multiple codons encode the same amino acid, especially differing only at the wobble (3rd) position. A nonsense mutation is generally more damaging than a missense mutation because the protein is cut short entirely.
\end{tcolorbox}

\newpage

%% ===========================================================================
%%  PART IIb – Cell Membranes & Transport
%% ===========================================================================
\tealsections
\addcontentsline{toc}{section}{\textcolor{sectionteal}{PART IIb \textemdash\ Cell Membranes \& Transport}}
\addcontentsline{toc}{subsection}{Membrane Structure \& Fluidity}
\addcontentsline{toc}{subsection}{Osmosis \& Tonicity}
\addcontentsline{toc}{subsection}{Types of Transport}
\partbanner{part5bg}{sectionteal}{Part IIb \quad Cell Membranes \& Transport}

\vspace{8pt}

% ── Membrane Structure ───────────────────────────────────────────────────────
\section*{Membrane Structure \& Fluidity}

\subsection*{The Fluid Mosaic Model}
The plasma membrane is described by the \textbf{fluid mosaic model}: it is a \textbf{phospholipid bilayer} in which proteins, cholesterol, and carbohydrates are embedded and can move laterally.

\begin{itemize}[noitemsep]
  \item \textbf{Phospholipids} have a hydrophilic (water-loving) head and two hydrophobic (water-fearing) fatty acid tails. They spontaneously arrange into a bilayer with heads facing outward and tails facing inward.
  \item \textbf{Integral proteins} (transmembrane) span the bilayer and may form channels or carriers.
  \item \textbf{Peripheral proteins} are loosely attached to the surface.
  \item \textbf{Carbohydrates} attached to proteins (glycoproteins) or lipids (glycolipids) on the extracellular surface serve as cell-identity tags and signaling molecules.
\end{itemize}

\subsection*{Membrane Fluidity \& Temperature Adaptation}
Membrane fluidity is critical -- too rigid and transport stops; too fluid and the membrane loses integrity. Cells adapt to temperature changes by altering membrane composition:

\noindent
{\renewcommand{\arraystretch}{1.5}
\begin{tabularx}{\textwidth}{>{\bfseries\raggedright\arraybackslash}p{3.5cm}>{\raggedright\arraybackslash}X>{\raggedright\arraybackslash}p{3.8cm}}
\toprule
\rowcolor{tableheaderteal}
\textbf{Adaptation} & \textbf{Effect on Membrane} & \textbf{When Used} \\
\midrule
\rowcolor{rowTe}
More \textbf{unsaturated} phospholipid tails
  & Unsaturated tails have kinks (double bonds) that prevent tight packing $\to$ \textbf{increases fluidity}
  & Cold temperatures (e.g., penguin feet) \\
\rowcolor{rowA}
More \textbf{saturated} phospholipid tails
  & Straight, tightly packed tails $\to$ \textbf{decreases fluidity} (more rigid)
  & Warm temperatures \\
\rowcolor{rowTe}
\textbf{Cholesterol}
  & Acts as a \textbf{fluidity buffer}: at cold temps it prevents crystallization (increases fluidity); at warm temps it restrains excessive movement (decreases fluidity)
  & Both extremes; stabilizes membrane \\
\bottomrule
\end{tabularx}}

\vspace{4pt}
\begin{tcolorbox}[colback=part5bg, colframe=sectionteal, arc=3pt, boxrule=1pt]
\textbf{Exam tip:} A penguin's feet are exposed to near-freezing temperatures. To keep the membrane fluid enough for transport, cells would \textbf{increase unsaturated phospholipid tails} (primary adaptation) and \textbf{increase cholesterol} (secondary stabilizer). More saturated tails would make the membrane too rigid -- the wrong direction.
\end{tcolorbox}

% ── Osmosis & Tonicity ───────────────────────────────────────────────────────
\section*{Osmosis \& Tonicity}

\textbf{Osmosis} is the passive diffusion of \emph{water} across a selectively permeable membrane, from a region of \textbf{lower solute concentration} (more water) to \textbf{higher solute concentration} (less water).

\vspace{4pt}
\textbf{Tonicity} describes the relative solute concentration of a solution compared to the cell's cytoplasm:

\noindent
{\renewcommand{\arraystretch}{1.5}
\begin{tabularx}{\textwidth}{>{\bfseries\raggedright\arraybackslash}p{3cm}>{\raggedright\arraybackslash}X>{\raggedright\arraybackslash}p{3.5cm}}
\toprule
\rowcolor{tableheaderteal}
\textbf{Term} & \textbf{Definition} & \textbf{Effect on Cell} \\
\midrule
\rowcolor{rowTe}
Hypertonic
  & External solution has \textbf{higher} solute concentration than the cell
  & Water leaves cell $\to$ cell shrivels \\
\rowcolor{rowA}
Hypotonic
  & External solution has \textbf{lower} solute concentration than the cell
  & Water enters cell $\to$ cell swells \\
\rowcolor{rowTe}
Isotonic
  & External solution has \textbf{equal} solute concentration to the cell
  & No net water movement \\
\bottomrule
\end{tabularx}}

\vspace{4pt}
\textit{Example: A cell with 0.2 M glucose placed in a 0.4 M glucose solution -- the external solution is \textbf{hypertonic} (higher solute outside), so water will leave the cell.}

% ── Types of Transport ───────────────────────────────────────────────────────
\section*{Types of Transport}

\subsection*{Passive Transport (no energy required)}
\begin{itemize}[noitemsep]
  \item \textbf{Simple Diffusion:} Small, nonpolar molecules (O$_2$, CO$_2$, H$_2$O) slip directly through the phospholipid bilayer from high $\to$ low concentration. \textit{No proteins needed.}
  \item \textbf{Facilitated Diffusion:} Larger or polar molecules (e.g., glucose) move from high $\to$ low concentration \emph{through specific transport proteins} (channels or carriers). The proteins are \textbf{specific} to the substances they help move. No energy required.
  \item \textbf{Osmosis:} The specific case of water diffusing across a membrane via aquaporins or directly through the bilayer.
\end{itemize}

\begin{tcolorbox}[colback=part5bg, colframe=sectionteal, arc=3pt, boxrule=1pt]
\textbf{Key distinction:} Diffusion always moves molecules from \textbf{high to low} concentration (down the gradient). A common misconception is that diffusion goes from low to high -- this is \textbf{false}; that would be active transport.
\end{tcolorbox}

\subsection*{Active Transport (energy required)}
\begin{itemize}[noitemsep]
  \item Moves molecules \textbf{against} their concentration gradient (low $\to$ high) -- requires energy (\textbf{ATP}).
  \item Uses \textbf{protein pumps} (e.g., sodium-potassium pump).
  \item \textbf{Co-transport (coupled transport):} One protein pumps one substance against its gradient; a second protein uses the gradient established to move another substance.
\end{itemize}

\subsection*{Bulk Transport (vesicle-mediated)}

\noindent
{\renewcommand{\arraystretch}{1.5}
\begin{tabularx}{\textwidth}{>{\bfseries\raggedright\arraybackslash}p{4cm}>{\raggedright\arraybackslash}X}
\toprule
\rowcolor{tableheaderteal}
\textbf{Process} & \textbf{Description} \\
\midrule
\rowcolor{rowTe}
Endocytosis (general)
  & Cell engulfs material by wrapping membrane around it to form a vesicle \\
\rowcolor{rowA}
Phagocytosis (``cell eating'')
  & Cell engulfs large solid particles (e.g., bacteria); forms a \textbf{phagosome} \\
\rowcolor{rowTe}
Pinocytosis (``cell drinking'')
  & Cell takes in extracellular fluid and small dissolved molecules in tiny vesicles \\
\rowcolor{rowA}
Receptor-mediated endocytosis
  & \textbf{Specific} ligands bind to receptors on the membrane before being brought in; \textbf{not} non-specific \\
\rowcolor{rowTe}
Exocytosis
  & Cell secretes materials by fusing vesicles with the plasma membrane; exports substances out of the cell \\
\bottomrule
\end{tabularx}}

\vspace{4pt}
\begin{tcolorbox}[colback=part5bg, colframe=sectionteal, arc=3pt, boxrule=1pt]
\textbf{Exam tip:} Receptor-mediated endocytosis is \textbf{specific} -- only substances that bind to the specific receptor are brought in. This is often tested as a true/false: ``receptor-mediated endocytosis is non-specific'' -- this is \textbf{FALSE}.
\end{tcolorbox}

\newpage

%% ===========================================================================
%%  PART IIc – Energy, Enzymes & Metabolism
%% ===========================================================================
\orangesections
\addcontentsline{toc}{section}{\textcolor{sectionorange}{PART IIc \textemdash\ Energy, Enzymes \& Metabolism}}
\addcontentsline{toc}{subsection}{Gibbs Free Energy}
\addcontentsline{toc}{subsection}{Enzymes}
\addcontentsline{toc}{subsection}{ATP -- The Energy Currency of the Cell}
\addcontentsline{toc}{subsection}{Photosynthesis vs. Cellular Respiration}
\addcontentsline{toc}{subsection}{Redox Reactions}
\addcontentsline{toc}{subsection}{Glycolysis}
\addcontentsline{toc}{subsection}{Cellular Respiration Summary}
\partbanner{part4bg}{sectionorange}{Part IIc \quad Energy, Enzymes \& Metabolism}

\vspace{8pt}

% ── Gibbs Free Energy ────────────────────────────────────────────────────────
\section*{Gibbs Free Energy ($\Delta G$)}

The \textbf{Gibbs free energy} equation determines whether a reaction is spontaneous:
\[
  G = H - TS \qquad \Longrightarrow \qquad \Delta G = \Delta H - T\Delta S
\]

where $H$ = enthalpy (heat content), $T$ = temperature (Kelvin), $S$ = entropy (disorder).

\vspace{4pt}
\noindent
{\renewcommand{\arraystretch}{1.5}
\begin{tabularx}{\textwidth}{>{\bfseries\raggedright\arraybackslash}p{3cm}>{\raggedright\arraybackslash}X>{\raggedright\arraybackslash}p{4cm}}
\toprule
\rowcolor{tableheaderorange}
\textbf{$\Delta G$ sign} & \textbf{Meaning} & \textbf{Reaction type} \\
\midrule
\rowcolor{rowOr}
$\Delta G < 0$ (negative)
  & Products have less free energy than reactants; energy is \textbf{released}; reaction is \textbf{spontaneous}
  & \textbf{Exergonic / Catabolic} \\
\rowcolor{rowA}
$\Delta G > 0$ (positive)
  & Products have more free energy than reactants; energy must be \textbf{input}; reaction is \textbf{non-spontaneous}
  & \textbf{Endergonic / Anabolic} \\
\rowcolor{rowOr}
$\Delta G = 0$
  & System is at chemical \textbf{equilibrium}; no net change in free energy
  & Equilibrium \\
\bottomrule
\end{tabularx}}

\vspace{6pt}
\textbf{Entropy and energy:}
\begin{itemize}[noitemsep]
  \item \textbf{Less disorder} (low entropy, high order) = \textbf{more stored chemical energy}
  \item \textbf{More disorder} (high entropy) = \textbf{less stored chemical energy}
  \item Higher $G$ = less stable, greater capacity to do work.
\end{itemize}

\vspace{4pt}
\begin{tcolorbox}[colback=part4bg, colframe=sectionorange, arc=3pt, boxrule=1pt]
\textbf{Memory:} \textbf{Catabolism} = breaking down (e.g., cellular respiration) = \textbf{exergonic} ($\Delta G < 0$, releases energy). \textbf{Anabolism} = building up (e.g., protein synthesis) = \textbf{endergonic} ($\Delta G > 0$, consumes energy). ``\textit{Catabolism is to anabolism as exergonic is to endergonic.}''
\end{tcolorbox}

% ── Enzymes ──────────────────────────────────────────────────────────────────
\section*{Enzymes}

Enzymes are biological \textbf{catalysts} -- they \textbf{increase the rate} of a reaction without being consumed or permanently altered by it.

\subsection*{How Enzymes Work}
\begin{itemize}[noitemsep]
  \item Enzymes \textbf{lower the activation energy} (E$_a$) barrier of a reaction, allowing it to proceed faster. They do \emph{not} change the overall $\Delta G$ of the reaction or whether it is spontaneous.
  \item They stabilize the \textbf{transition state} and provide an alternative reaction pathway.
  \item The \textbf{active site} is the region of the enzyme that binds the substrate; it is highly specific to the substrate's shape and chemistry.
\end{itemize}

\begin{tcolorbox}[colback=part4bg, colframe=sectionorange, arc=3pt, boxrule=1pt]
\textbf{What enzymes do NOT do:}
\begin{itemize}[noitemsep, leftmargin=*]
  \item Do \textbf{not} greatly reduce the free energy given off -- they only lower $E_a$.
  \item Do \textbf{not} change the direction a reaction goes.
  \item Do \textbf{not} become permanently altered by the reactions they catalyze.
  \item Do \textbf{not} prevent changes in substrate/product concentration.
\end{itemize}
\end{tcolorbox}

\subsection*{Activation Energy Diagram}
\[
\underbrace{\text{Reactants}}_{\text{start}} \;\xrightarrow{\quad E_a \text{ (lowered by enzyme)}\quad}\; \underbrace{\text{Transition State}}_{\text{peak}} \;\longrightarrow\; \underbrace{\text{Products}}_{\text{end}}
\]
A \textbf{starch solution at room temperature} does not readily break down to sugars because the \textbf{activation energy barrier cannot be overcome} without an enzyme (amylase) or added heat -- even though the hydrolysis of starch to sugar is thermodynamically favorable (exergonic).

% ── ATP ──────────────────────────────────────────────────────────────────────
\section*{ATP -- The Energy Currency of the Cell}

\textbf{ATP} (adenosine triphosphate) is the primary energy carrier in cells. It consists of a nitrogenous base (adenine), a ribose sugar, and 3 phosphate groups.

\subsection*{Why ATP is Important in Metabolism}
\begin{itemize}[noitemsep]
  \item When ATP is \textbf{hydrolyzed} (terminal phosphate bond broken), energy is released and can be used to drive \textbf{endergonic} (non-spontaneous) reactions.
  \item ATP \textbf{couples} exergonic and endergonic reactions: energy released from an exergonic reaction is captured in ATP, which then powers an endergonic reaction. This is called \textbf{energy coupling}.
  \item A \textbf{phosphorylated} molecule (one that has received a phosphate group from ATP) has \textbf{increased chemical reactivity} -- it can perform cellular work.
\end{itemize}

\begin{tcolorbox}[colback=part4bg, colframe=sectionorange, arc=3pt, boxrule=1pt]
\textbf{Common misconception:} The terminal phosphate bond in ATP is \textbf{not} strongly covalent -- it is actually relatively \textbf{unstable}, which is precisely why breaking it releases energy. The instability of the bond stores potential energy.
\end{tcolorbox}

\begin{tcolorbox}[colback=part2bg, colframe=sectiongreen,
  title={\bfseries How Do Bacteria Make ATP Without Mitochondria?}, arc=3pt, boxrule=1pt]
Bacteria use their \textbf{plasma membrane} in place of the inner mitochondrial membrane. They pump protons (H$^+$) out of the cell to create a \textbf{proton gradient} across the membrane. ATP synthase, embedded in the plasma membrane, harnesses this gradient to produce ATP as protons flow back into the cell -- the same mechanism used in mitochondria, just at a different membrane.
\end{tcolorbox}

% ── Photosynthesis vs Respiration ────────────────────────────────────────────
\section*{Photosynthesis vs.\ Cellular Respiration}

These two processes are essentially \textbf{opposites} -- one stores energy, the other releases it.

\subsection*{Overall Equations}
\[
  \textbf{Photosynthesis:} \quad 6\text{CO}_2 + 6\text{H}_2\text{O} + \text{light energy} \;\longrightarrow\; \text{C}_6\text{H}_{12}\text{O}_6 + 6\text{O}_2
\]
\[
  \textbf{Cellular Respiration:} \quad \text{C}_6\text{H}_{12}\text{O}_6 + 6\text{O}_2 \;\longrightarrow\; 6\text{CO}_2 + 6\text{H}_2\text{O} + \text{ATP (energy)}
\]

\vspace{4pt}
\noindent
{\renewcommand{\arraystretch}{1.5}
\begin{tabularx}{\textwidth}{>{\bfseries\raggedright\arraybackslash}p{3.5cm}
                             >{\raggedright\arraybackslash}X
                             >{\raggedright\arraybackslash}X}
\toprule
\rowcolor{tableheaderorange}
\textbf{Feature} & \textbf{Photosynthesis} & \textbf{Cellular Respiration} \\
\midrule
\rowcolor{rowOr} Energy change & Stores energy (endergonic/anabolic) & Releases energy (exergonic/catabolic) \\
\rowcolor{rowA}  Location & Chloroplasts (plants only) & Mitochondria (all eukaryotes) \\
\rowcolor{rowOr} Who does it & Photoautotrophs (plants, algae, some bacteria) & Virtually all living organisms \\
\rowcolor{rowA}  Inputs & CO$_2$, H$_2$O, light & Glucose, O$_2$ \\
\rowcolor{rowOr} Outputs & Glucose, O$_2$ & CO$_2$, H$_2$O, ATP \\
\rowcolor{rowA}  Relationship & \multicolumn{2}{>{\raggedright\arraybackslash}X}{Photosynthesis \textbf{stores} energy in glucose; respiration \textbf{releases} it. They are complementary, not reverses of each other's exact biochemical pathways.} \\
\bottomrule
\end{tabularx}}

\vspace{4pt}
\begin{tcolorbox}[colback=part4bg, colframe=sectionorange, arc=3pt, boxrule=1pt]
\textbf{Trophic types:}
\begin{itemize}[noitemsep, leftmargin=*]
  \item \textbf{Photoautotrophs} -- use light as energy source; produce own organic molecules (e.g., plants, algae).
  \item \textbf{Photoheterotrophs} -- use light for energy but obtain organic carbon from the environment.
  \item \textbf{Chemoautotrophs} -- use inorganic chemicals as energy source; synthesize own organic molecules.
  \item \textbf{Chemoheterotrophs} -- obtain both energy and carbon from organic molecules (most animals, fungi).
\end{itemize}
\end{tcolorbox}

\vspace{4pt}
\textbf{Potential energy examples:} A carbohydrate made of glucose monomers (chemical potential energy); a compressed spring; a charged battery. Note: \textit{a boy mid-dive, water falling over a waterfall, or electrons being passed along the electron transport chain} are examples of \textbf{kinetic} energy in action.


% ── Redox Reactions ──────────────────────────────────────────────────────────
\section*{Redox Reactions}

Cellular respiration is fundamentally a series of \textbf{redox (reduction-oxidation) reactions} in which electrons are transferred between molecules. This electron transfer is what drives ATP production.

\begin{tcolorbox}[colback=craapbg, colframe=craapborder,
  title={\bfseries OIL RIG -- Remembering Redox}, arc=3pt, boxrule=1pt]
\begin{multicols}{2}
\centering
\textbf{OIL}\\
\textbf{O}xidation \textbf{I}s \textbf{L}oss\\
(of electrons)\\[6pt]
A molecule that loses electrons is \textbf{oxidized}.\\
It is the \textbf{reducing agent}.
\columnbreak
\centering
\textbf{RIG}\\
\textbf{R}eduction \textbf{I}s \textbf{G}ain\\
(of electrons)\\[6pt]
A molecule that gains electrons is \textbf{reduced}.\\
It is the \textbf{oxidizing agent}.
\end{multicols}
\end{tcolorbox}

\vspace{6pt}

\subsection*{Key Electron Carriers}

The cell uses \textbf{electron carriers} to shuttle electrons from glucose oxidation to the electron transport chain. They exist in oxidized and reduced forms:

\vspace{4pt}
\noindent
{\small\renewcommand{\arraystretch}{1.5}
\begin{tabularx}{\textwidth}{>{\bfseries\raggedright\arraybackslash}p{2.8cm}
                              >{\centering\arraybackslash}p{2.4cm}
                              >{\centering\arraybackslash}p{2.4cm}
                              >{\raggedright\arraybackslash}X}
\toprule
\rowcolor{tableheader}
\textbf{Carrier} & \textbf{Oxidized Form} & \textbf{Reduced Form} & \textbf{What Happens} \\
\midrule
\rowcolor{rowA} NAD$^+$ / NADH & NAD$^+$ & NADH & NAD$^+$ gains 2e$^-$ + H$^+$ $\to$ NADH (reduced; carries energy) \\
\rowcolor{rowBg} FAD / FADH$_2$ & FAD & FADH$_2$ & FAD gains 2e$^-$ + 2H$^+$ $\to$ FADH$_2$ (reduced; carries energy) \\
\bottomrule
\end{tabularx}}

\vspace{4pt}
\textit{Remember: the \textbf{reduced} forms (NADH, FADH$_2$) are the ``loaded'' carriers -- they carry electrons to the electron transport chain to make ATP. The oxidized forms (NAD$^+$, FAD) are the empty carriers, ready to pick up more electrons.}

\subsection*{Overall Equation for Cellular Respiration}
\[
  \text{C}_6\text{H}_{12}\text{O}_6 + 6\,\text{O}_2 \;\longrightarrow\; 6\,\text{CO}_2 + 6\,\text{H}_2\text{O} + \text{ATP (energy)}
\]
Glucose is \textbf{oxidized} (loses electrons); oxygen is \textbf{reduced} (gains electrons). This reaction is the reverse of photosynthesis.

% ── Glycolysis ───────────────────────────────────────────────────────────────
\section*{Glycolysis}

Glycolysis (``splitting glucose'') is the \textbf{first stage of cellular respiration}. It occurs in the \textbf{cytoplasm} and does \textbf{not require oxygen}, making it universal to nearly all life.

\textbf{Net summary:}
\[
  1 \text{ Glucose (6C)} \;\longrightarrow\; 2 \text{ Pyruvate (3C)} + 2\,\text{ATP (net)} + 2\,\text{NADH}
\]

\subsection*{Three Phases of Glycolysis}

\begin{tcolorbox}[colback=partbg, colframe=sectionblue, arc=3pt, boxrule=1pt]
\begin{enumerate}[leftmargin=*, noitemsep]
  \item \textbf{Energy Investment Phase} (costs 2 ATP)\\
    Glucose is phosphorylated \textbf{twice} using 2 ATP, making it more reactive and trapping it inside the cell. This produces a 6-carbon molecule stabilized by two phosphate groups.

  \item \textbf{Cleavage Phase}\\
    An enzyme cleaves the 6-carbon molecule into \textbf{two 3-carbon molecules} called \textbf{G3P} (glyceraldehyde-3-phosphate).

  \item \textbf{Energy Payoff Phase} (yields 4 ATP + 2 NADH)\\
    Each G3P is \textbf{oxidized} and phosphorylated: 2 NAD$^+$ are reduced to \textbf{2 NADH} (electron carriers). \textbf{Substrate-level phosphorylation} then directly transfers phosphate groups from the substrate to ADP, producing \textbf{4 ATP} total.
\end{enumerate}
\end{tcolorbox}

\vspace{4pt}

\begin{center}
\renewcommand{\arraystretch}{1.3}
\begin{tabular}{lcc}
\toprule
\rowcolor{tableheader}
\textbf{Phase} & \textbf{ATP Used} & \textbf{ATP / NADH Produced} \\
\midrule
\rowcolor{rowA} Investment (phosphorylation) & $-$2 ATP & --- \\
\rowcolor{rowBg} Cleavage                    & ---      & --- \\
\rowcolor{rowA} Payoff (substrate-level phosphorylation) & --- & $+$4 ATP, $+$2 NADH \\
\midrule
\rowcolor{tableheader} \textbf{Net Yield} & & \textbf{+2 ATP, +2 NADH} \\
\bottomrule
\end{tabular}
\end{center}

\vspace{4pt}
\textbf{Substrate-level phosphorylation} = a phosphate group is transferred \emph{directly} from a substrate molecule to ADP, forming ATP. This does not require the electron transport chain or oxygen.

\vspace{4pt}
\textit{Note: Glycolysis is only the beginning. The 2 NADH and 2 pyruvate produced here continue into the Krebs cycle and electron transport chain (when oxygen is available) to generate the bulk of the cell's ATP.}

\vspace{6pt}

% ── Cellular Respiration Summary ─────────────────────────────────────────────
\section*{Cellular Respiration Summary}

\noindent
{\small\renewcommand{\arraystretch}{1.5}
\begin{tabularx}{\textwidth}{>{\bfseries\raggedright\arraybackslash}p{3.8cm}
                              >{\raggedright\arraybackslash}p{4cm}
                              >{\centering\arraybackslash}p{1.3cm}
                              >{\centering\arraybackslash}p{1.3cm}
                              >{\centering\arraybackslash}p{1.5cm}
                              >{\centering\arraybackslash}p{1.3cm}}
\toprule
\rowcolor{tableheader}
\textbf{Stage} & \textbf{Location} & \textbf{ATP} & \textbf{NADH} & \textbf{FADH$_2$} & \textbf{CO$_2$} \\
\midrule
\rowcolor{rowA}   Glycolysis              & Cytoplasm                    & 2 (net) & 2 & 0 & 0 \\
\rowcolor{rowB}   Pyruvate Oxidation      & Mitochondrial matrix         & 0       & 2 & 0 & 2 \\
\rowcolor{rowA}   Krebs Cycle (CAC)       & Mitochondrial matrix         & 2       & 6 & 2 & 4 \\
\rowcolor{rowB}   ETC + ATP Synthase      & Inner mitochondrial membrane & $\sim$26--28 & --- & --- & 0 \\
\midrule
\rowcolor{tableheader} \textbf{Total (per glucose)} & & \textbf{$\sim$30--32} & \textbf{10} & \textbf{2} & \textbf{6} \\
\bottomrule
\end{tabularx}}

\vspace{4pt}
\begin{tcolorbox}[colback=partbg, colframe=sectionblue, arc=3pt, boxrule=1pt]
\textbf{Key ETC points:}
\begin{itemize}[noitemsep, leftmargin=*]
  \item \textbf{O$_2$ is the final electron acceptor} in the ETC. It accepts electrons (and H$^+$) at the end of the chain, forming \textbf{water} as a byproduct.
  \item Without O$_2$, the ETC backs up and ATP production halts -- this is why aerobic respiration requires oxygen while glycolysis does not.
  \item The \textbf{mitochondrial matrix} is the site of both pyruvate oxidation and the Krebs cycle; the \textbf{inner mitochondrial membrane} houses the ETC and ATP synthase.
\end{itemize}
\end{tcolorbox}

\newpage

%% ===========================================================================
%%  PART IIf – Photosynthesis
%% ===========================================================================
\greensections
\addcontentsline{toc}{section}{\textcolor{sectiongreen}{PART IIf \textemdash\ Photosynthesis}}
\addcontentsline{toc}{subsection}{Nutritional Classification of Organisms}
\addcontentsline{toc}{subsection}{Light \& the Electromagnetic Spectrum}
\addcontentsline{toc}{subsection}{Thermodynamics of Photosynthesis}
\addcontentsline{toc}{subsection}{Chloroplast Structure}
\addcontentsline{toc}{subsection}{Light Reactions (Thylakoid Membranes)}
\addcontentsline{toc}{subsection}{Calvin Cycle (Stroma)}
\addcontentsline{toc}{subsection}{Photosystem I vs.\ Photosystem II}
\partbanner{part2bg}{sectiongreen}{Part IIf \quad Photosynthesis}

\vspace{8pt}

% ── Trophic Classification ────────────────────────────────────────────────────
\section*{Nutritional Classification of Organisms}

Organisms differ in how they obtain \textbf{energy} and \textbf{carbon}:

\noindent
{\renewcommand{\arraystretch}{1.5}
\begin{tabularx}{\textwidth}{>{\bfseries\raggedright\arraybackslash}p{3.8cm}>{\raggedright\arraybackslash}X>{\raggedright\arraybackslash}p{4cm}}
\toprule
\rowcolor{tableheadergreen}
\textbf{Type} & \textbf{Energy \& Carbon Source} & \textbf{Examples} \\
\midrule
\rowcolor{rowBg} Photoautotroph   & Light energy; builds own organic molecules from CO$_2$ & Plants, algae, cyanobacteria \\
\rowcolor{rowA}  Photoheterotroph & Light energy; obtains carbon from organic environment & Some purple non-sulfur bacteria \\
\rowcolor{rowBg} Chemoautotroph   & Inorganic chemicals for energy; synthesizes own organic molecules & Sulfur bacteria, nitrifying bacteria \\
\rowcolor{rowA}  Chemoheterotroph & Organic molecules for both energy and carbon & Most animals, fungi, decomposers \\
\bottomrule
\end{tabularx}}

\vspace{4pt}
\textit{Memory: ``auto'' = makes its own food; ``hetero'' = gets carbon externally; ``photo'' = uses light; ``chemo'' = uses chemical energy.}

% ── Light & the EM Spectrum ───────────────────────────────────────────────────
\section*{Light \& the Electromagnetic Spectrum}

Light is \textbf{electromagnetic radiation} that behaves simultaneously as a wave and as discrete particles called \textbf{photons}.

\begin{itemize}[noitemsep]
  \item \textbf{Visible light range:} approximately \textbf{400--700 nm} -- the portion of the spectrum that photosynthesis uses.
  \item \textbf{Wavelength and energy are inversely related:}
    \begin{itemize}[noitemsep]
      \item Shorter wavelength $\to$ \textbf{higher energy} (e.g., violet/blue $\approx$ 400--450 nm)
      \item Longer wavelength $\to$ \textbf{lower energy} (e.g., red $\approx$ 650--700 nm)
    \end{itemize}
  \item \textbf{Chlorophyll a} (the primary photosynthetic pigment) absorbs mostly \textbf{blue and red} wavelengths. It \emph{reflects} green light -- which is why plants appear green.
  \item \textbf{Accessory pigments} (chlorophyll b, carotenoids, xanthophylls) absorb additional wavelengths, broadening the usable range of the solar spectrum and transferring that energy to chlorophyll a.
\end{itemize}

\begin{tcolorbox}[colback=part2bg, colframe=sectiongreen, arc=3pt, boxrule=1pt]
\textbf{Why does chlorophyll reflect green?} Chlorophyll's molecular structure allows it to absorb the specific photon energies corresponding to red and blue light, exciting electrons to higher energy levels. Green photons do not match these energy levels and cannot be absorbed -- they are reflected back to our eyes.
\end{tcolorbox}

% ── Thermodynamics of Photosynthesis ─────────────────────────────────────────
\section*{Thermodynamics of Photosynthesis}

Photosynthesis is an \textbf{endergonic} (anabolic) process: it uses light energy to build energy-rich glucose.

\[
  6\,\text{CO}_2 + 6\,\text{H}_2\text{O} + \text{light energy} \;\longrightarrow\; \text{C}_6\text{H}_{12}\text{O}_6 + 6\,\text{O}_2 \quad (\Delta G > 0)
\]

It is also a \textbf{redox reaction}:
\begin{itemize}[noitemsep]
  \item \textbf{CO$_2$ is reduced} -- gains electrons (hydrogen atoms added); converted into glucose.
  \item \textbf{H$_2$O is oxidized} -- loses electrons; split to release \textbf{O$_2$} as a byproduct.
\end{itemize}

\begin{tcolorbox}[colback=part2bg, colframe=sectiongreen, arc=3pt, boxrule=1pt]
\textbf{Compare to cellular respiration:} The two processes are chemical opposites.
Photosynthesis is \textbf{endergonic} -- it \emph{stores} light energy as chemical energy in glucose ($\Delta G > 0$).
Cellular respiration is \textbf{exergonic} -- it \emph{releases} that stored energy as ATP ($\Delta G < 0$).
\end{tcolorbox}

% ── Chloroplast Structure ─────────────────────────────────────────────────────
\section*{Chloroplast Structure}

The chloroplast is a double-membrane organelle found in plant cells and algae. Its internal organization separates the two stages of photosynthesis:

\begin{itemize}[noitemsep]
  \item \textbf{Outer \& inner membranes} -- the double-membrane boundary (evidence for endosymbiotic origin).
  \item \textbf{Thylakoids} -- flattened, membrane-bound sacs stacked into \textbf{grana} (sing.\ granum). The \textbf{light reactions} occur here.
  \item \textbf{Thylakoid lumen} -- the interior space of the thylakoid membrane; H$^+$ ions accumulate here to build the proton gradient.
  \item \textbf{Stroma} -- the fluid-filled space surrounding the thylakoids. The \textbf{Calvin cycle} occurs here.
\end{itemize}

\textit{Memory: \textbf{Thylakoid} = light reactions (``thy-light-koid''); \textbf{Stroma} = Calvin cycle (happens in the ``stroma-ch'' of the chloroplast).}

% ── Light Reactions ───────────────────────────────────────────────────────────
\section*{Light Reactions (Thylakoid Membranes)}

The light reactions convert \textbf{solar energy} into \textbf{chemical energy} (ATP and NADPH) and split water to release O$_2$.

\subsection*{Photosystem II (PSII) -- Acts First}

\begin{tcolorbox}[colback=part2bg, colframe=sectiongreen, arc=3pt, boxrule=1pt]
\textbf{Naming note:} PSII was \emph{discovered after} PSI but acts \emph{first} in the electron transport pathway. The numbering reflects discovery order, not functional order. This is a common exam trap.
\end{tcolorbox}

\begin{enumerate}[noitemsep]
  \item \textbf{Pigment absorption} -- Antenna pigments (chlorophyll a, chlorophyll b, carotenoids) absorb light and funnel energy to the \textbf{P680 reaction center}.
  \item \textbf{Electron excitation} -- Light energy excites electrons in P680 to a higher energy level; they are transferred to an electron acceptor molecule.
  \item \textbf{Photolysis (water splitting)} -- To replace the lost electrons, PSII splits water:
  \[
    2\,\text{H}_2\text{O} \;\longrightarrow\; 4\,\text{H}^+ + 4\,e^- + \text{O}_2\uparrow
  \]
  The O$_2$ released is the \textbf{oxygen we breathe}.
  \item \textbf{Electron Transport Chain (ETC)} -- Excited electrons pass along the ETC (plastoquinone $\to$ cytochrome b6f $\to$ plastocyanin), releasing energy that \textbf{pumps H$^+$ into the thylakoid lumen}, building a proton gradient.
  \item \textbf{ATP synthesis (photophosphorylation)} -- H$^+$ ions flow back out through \textbf{ATP synthase} via chemiosmosis, driving the synthesis of \textbf{ATP}.
\end{enumerate}

\subsection*{Photosystem I (PSI) -- Produces NADPH}

\begin{enumerate}[noitemsep]
  \item Electrons arriving from the ETC (now lower energy) enter the \textbf{P700 reaction center} and are re-energized by another photon of light.
  \item The re-energized electrons are passed to \textbf{ferredoxin}, then used to reduce NADP$^+$ to \textbf{NADPH} (catalyzed by NADP$^+$ reductase).
\end{enumerate}

\begin{tcolorbox}[colback=part2bg, colframe=sectiongreen, arc=3pt, boxrule=1pt]
\textbf{Outputs of the Light Reactions:}
\begin{itemize}[noitemsep, leftmargin=*]
  \item \textbf{ATP} -- produced via photophosphorylation (chemiosmosis through ATP synthase).
  \item \textbf{NADPH} -- produced at PSI; carries electrons and reducing power to the Calvin cycle.
  \item \textbf{O$_2$} -- released as a byproduct of water splitting at PSII.
\end{itemize}
\end{tcolorbox}

\subsection*{Z-Scheme: Overall Electron Flow}

\[
  \underbrace{\text{H}_2\text{O}}_{\text{electron donor}} \;\longrightarrow\; \underbrace{\text{PSII (P680)}}_{\text{light}} \;\longrightarrow\; \text{ETC (ATP)} \;\longrightarrow\; \underbrace{\text{PSI (P700)}}_{\text{light}} \;\longrightarrow\; \underbrace{\text{NADPH}}_{\text{electron carrier}}
\]

Electrons travel from water (lowest energy) through two photosystems and end in NADPH (high chemical energy). The energy diagram of this pathway has a Z-shape, giving it the name ``Z-scheme.''

% ── Calvin Cycle ──────────────────────────────────────────────────────────────
\section*{Calvin Cycle (Stroma)}

The Calvin cycle uses \textbf{ATP and NADPH} from the light reactions to convert CO$_2$ into sugar. It does \emph{not} directly use light, but it depends entirely on the light reactions' products.

\textbf{Net inputs and output (to make 1 G3P):}
\[
  3\,\text{CO}_2 + 9\,\text{ATP} + 6\,\text{NADPH} \;\longrightarrow\; 1\,\text{G3P (net)} \quad (\times 2 \text{ for glucose})
\]

\begin{tcolorbox}[colback=part2bg, colframe=sectiongreen, arc=3pt, boxrule=1pt]
\textbf{Three stages of the Calvin Cycle:}
\begin{enumerate}[noitemsep, leftmargin=*]
  \item \textbf{Carbon Fixation} -- The enzyme \textbf{Rubisco} attaches CO$_2$ to the 5-carbon acceptor \textbf{RuBP} (ribulose-1,5-bisphosphate). The unstable 6-carbon intermediate immediately splits into \textbf{two 3-PGA} (3-phosphoglycerate) molecules.
  \item \textbf{Reduction} -- ATP and NADPH are used to reduce 3-PGA into \textbf{G3P} (glyceraldehyde-3-phosphate). G3P is the usable sugar product of the cycle; it can exit to form glucose or other organic molecules.
  \item \textbf{Regeneration of RuBP} -- Most G3P molecules (5 out of every 6) are used with ATP to regenerate \textbf{RuBP}, allowing the cycle to continue. For every 3 CO$_2$ fixed, only 1 net G3P exits.
\end{enumerate}
\end{tcolorbox}

\vspace{4pt}
\textit{Note: It takes \textbf{six turns} of the Calvin cycle (fixing 6 CO$_2$) to produce enough G3P to synthesize one glucose molecule.}

% ── Photosystem Comparison ────────────────────────────────────────────────────
\section*{Photosystem I vs.\ Photosystem II}

\noindent
{\renewcommand{\arraystretch}{1.5}
\begin{tabularx}{\textwidth}{>{\bfseries\raggedright\arraybackslash}p{3.5cm}>{\raggedright\arraybackslash}X>{\raggedright\arraybackslash}X}
\toprule
\rowcolor{tableheadergreen}
\textbf{Feature} & \textbf{Photosystem II (PSII)} & \textbf{Photosystem I (PSI)} \\
\midrule
\rowcolor{rowBg} Reaction center & P680 (absorbs 680 nm light) & P700 (absorbs 700 nm light) \\
\rowcolor{rowA}  Order in pathway & \textbf{First} in linear electron flow & \textbf{Second} in linear electron flow \\
\rowcolor{rowBg} Electron source & Water (H$_2$O -- photolysis) & Electrons arriving from ETC (from PSII) \\
\rowcolor{rowA}  Key product & O$_2$ (water splitting) + H$^+$ gradient for ATP & \textbf{NADPH} \\
\rowcolor{rowBg} Why important & Source of O$_2$ we breathe; drives ATP synthesis via chemiosmosis & Produces reducing power for the Calvin cycle \\
\bottomrule
\end{tabularx}}

\vspace{4pt}
\begin{tcolorbox}[colback=part2bg, colframe=sectiongreen, arc=3pt, boxrule=1pt]
\textbf{Exam-ready summary:} Light reactions occur in the \textbf{thylakoid membranes}; outputs are \textbf{ATP, NADPH, and O$_2$}. Calvin cycle occurs in the \textbf{stroma}; uses ATP and NADPH to fix CO$_2$ into \textbf{G3P} (precursor to glucose). Photosystem II acts \textbf{first} and produces O$_2$; PSI acts \textbf{second} and produces NADPH.
\end{tcolorbox}

\newpage



\newpage

%% ===========================================================================
%%  PART IIg – Cancer Biology
%% ===========================================================================
\bluesections
\addcontentsline{toc}{section}{\textcolor{sectionblue}{PART IIg \textemdash\ Cancer Biology}}
\addcontentsline{toc}{subsection}{What Is Cancer?}
\addcontentsline{toc}{subsection}{Why Cancer Matters}
\addcontentsline{toc}{subsection}{Classifications of Cancer}
\addcontentsline{toc}{subsection}{How Cancer Is Treated}
\addcontentsline{toc}{subsection}{Challenges in Treatment}
\addcontentsline{toc}{subsection}{Future Directions in Cancer Research}
\addcontentsline{toc}{subsection}{Connections to Other Fields}
\partbanner{partbg}{sectionblue}{Part IIg \quad Cancer Biology}

\vspace{8pt}

% ── What Is Cancer ────────────────────────────────────────────────────────────
\section*{What Is Cancer?}

Cancer is a disease of \textbf{uncontrolled cell division} caused by mutations in genes that regulate the cell cycle, DNA repair, or programmed cell death (\textbf{apoptosis}).

\begin{itemize}[noitemsep]
  \item Normal cells divide in a tightly regulated manner and undergo apoptosis when damaged.
  \item Cancer arises from mutations in two classes of regulatory genes:
    \begin{itemize}[noitemsep]
      \item \textbf{Proto-oncogenes} -- normally promote cell division. When mutated into \textbf{oncogenes}, they act like a stuck accelerator, continuously driving proliferation.
      \item \textbf{Tumor suppressor genes} (e.g., \textit{p53}, \textit{BRCA1/2}, \textit{RB}) -- normally inhibit cell division or trigger apoptosis. Mutations that silence them remove the ``brakes'' on division.
    \end{itemize}
  \item A \textbf{benign tumor} is a localized mass of abnormal cells that does not invade surrounding tissue.
  \item A \textbf{malignant tumor} (cancer) invades surrounding tissue and can \textbf{metastasize} -- spreading via blood or lymph to distant sites.
\end{itemize}

\begin{tcolorbox}[colback=partbg, colframe=sectionblue, arc=3pt, boxrule=1pt]
\textbf{Two defining hallmarks of cancer:}
\begin{enumerate}[noitemsep, leftmargin=*]
  \item \textbf{Uncontrolled proliferation} -- cells divide without the normal regulatory signals.
  \item \textbf{Metastasis} -- cancer cells detach from the original tumor, invade surrounding tissue, enter the bloodstream or lymph, and seed new tumors elsewhere.
\end{enumerate}
\end{tcolorbox}

% ── Why Cancer Matters ────────────────────────────────────────────────────────
\section*{Why Cancer Matters}

\begin{itemize}[noitemsep]
  \item Cancer is the \textbf{second leading cause of death} worldwide (after cardiovascular disease), responsible for nearly 10 million deaths per year globally.
  \item It affects all ages, sexes, and populations, though risk increases substantially with age.
  \item The \textbf{economic burden} is enormous -- treatment, lost productivity, and research costs total trillions of dollars globally each year.
  \item Studying cancer drives fundamental advances in \textbf{cell biology, molecular genetics, immunology, and pharmacology}.
\end{itemize}

% ── Classifications ───────────────────────────────────────────────────────────
\section*{Classifications of Cancer}

Cancer is named by the \textbf{tissue or cell type of origin}:

\noindent
{\renewcommand{\arraystretch}{1.5}
\begin{tabularx}{\textwidth}{>{\bfseries\raggedright\arraybackslash}p{3.2cm}>{\raggedright\arraybackslash}X>{\raggedright\arraybackslash}p{4.2cm}}
\toprule
\rowcolor{tableheader}
\textbf{Category} & \textbf{Cell / Tissue of Origin} & \textbf{Examples} \\
\midrule
\rowcolor{rowA} Carcinoma      & Epithelial cells (skin, organ linings) & Lung, breast, colon, prostate \\
\rowcolor{rowB} Sarcoma        & Connective tissue (bone, muscle, fat, cartilage) & Osteosarcoma, liposarcoma \\
\rowcolor{rowA} Leukemia       & Blood-forming (hematopoietic) cells in bone marrow & AML, CLL \\
\rowcolor{rowB} Lymphoma       & Lymphocytes in lymph nodes or spleen & Hodgkin's, Non-Hodgkin's \\
\rowcolor{rowA} Melanoma       & Melanocytes (pigment cells) in skin & Cutaneous melanoma \\
\rowcolor{rowB} CNS tumors     & Glial cells or neurons of brain/spinal cord & Glioblastoma, astrocytoma \\
\bottomrule
\end{tabularx}}

\vspace{4pt}
Cancers are also classified by \textbf{stage} (I--IV) -- how far the disease has spread -- and by \textbf{grade} -- how abnormal the cells appear under a microscope. Higher stage and grade generally indicate a worse prognosis.

% ── Treatments ────────────────────────────────────────────────────────────────
\section*{How Cancer Is Treated}

\noindent
{\renewcommand{\arraystretch}{1.5}
\begin{tabularx}{\textwidth}{>{\bfseries\raggedright\arraybackslash}p{3.5cm}>{\raggedright\arraybackslash}X}
\toprule
\rowcolor{tableheader}
\textbf{Treatment Type} & \textbf{Description} \\
\midrule
\rowcolor{rowA} Surgery            & Physical removal of the tumor; most effective when cancer is localized and has not metastasized. \\
\rowcolor{rowB} Radiation therapy  & High-energy radiation (X-rays or proton beams) damages DNA in cancer cells, halting division. Nearby normal tissue is also affected. \\
\rowcolor{rowA} Chemotherapy       & Drugs that kill rapidly dividing cells. Non-specific -- also damages fast-dividing healthy cells (hair follicles, gut lining), causing common side effects. \\
\rowcolor{rowB} Immunotherapy      & Boosts or redirects the immune system to destroy cancer cells. Includes \textbf{checkpoint inhibitors} (releasing immune brakes) and \textbf{CAR-T cell therapy} (engineering T cells to recognize tumor antigens). \\
\rowcolor{rowA} Targeted therapy   & Drugs that block specific molecular targets unique to cancer cells (e.g., mutated proteins, overactive receptors). More precise than chemotherapy; resistance can develop. \\
\rowcolor{rowB} Hormone therapy    & Blocks hormones that fuel hormone-sensitive cancers (e.g., estrogen in breast cancer; testosterone in prostate cancer). \\
\rowcolor{rowA} New: Antisense RNA, Angiogenesis inhibitors, nanotechnology    & New approaches in cancer treatment. \\
\bottomrule
\end{tabularx}}

% ── Challenges ────────────────────────────────────────────────────────────────
\section*{Challenges in Treating Cancer}

\begin{itemize}[noitemsep]
  \item \textbf{Tumor heterogeneity:} Even within a single tumor, different cells carry different mutations. A drug may eliminate most cells while leaving resistant variants that repopulate the tumor.
  \item \textbf{Metastasis:} Once cancer has spread to distant sites, surgery cannot remove all foci, and systemic treatments have limited reach.
  \item \textbf{Drug resistance:} Cancer cells evolve resistance to chemotherapy and targeted therapies through additional mutations or compensatory pathways -- a direct consequence of natural selection acting within the body.
  \item \textbf{Side effects:} Many treatments are toxic to healthy, rapidly dividing cells, limiting therapeutic doses.
  \item \textbf{Late detection:} Many cancers are asymptomatic until advanced stages; early detection dramatically improves survival rates but requires accessible screening programs.
  \item \textbf{Cost and access:} Novel therapies (e.g., CAR-T cell therapy can cost \$400,000+) remain inaccessible to many patients worldwide.
\end{itemize}

% ── Future Directions ─────────────────────────────────────────────────────────
\section*{Future Directions in Cancer Research}

\begin{itemize}[noitemsep]
  \item \textbf{Liquid biopsies:} Detecting circulating tumor DNA (ctDNA) in blood for early, non-invasive diagnosis and monitoring of treatment response.
  \item \textbf{Precision / personalized medicine:} Sequencing a patient's tumor genome to identify specific mutations and tailor treatment accordingly.
  \item \textbf{mRNA therapeutics:} Using mRNA vaccine technology (similar to COVID-19 mRNA vaccines) to train the immune system against patient-specific tumor antigens (neoantigen vaccines).
  \item \textbf{CRISPR gene editing:} Correcting cancer-driving mutations or engineering immune cells (e.g., T cells) to more effectively target tumors.
  \item \textbf{Artificial intelligence:} Machine learning models trained on genomic and pathology data to improve diagnostic accuracy, predict treatment outcomes, and accelerate drug discovery.
  \item \textbf{Combination therapies:} Simultaneously attacking tumors via multiple mechanisms (e.g., targeted therapy + immunotherapy) to reduce the chance of resistance evolving.
\end{itemize}

% ── Connections to Other Fields ───────────────────────────────────────────────
\section*{Connections to Other Fields}

\begin{tcolorbox}[colback=partbg, colframe=sectionblue, arc=3pt, boxrule=1pt]
\begin{itemize}[noitemsep, leftmargin=*]
  \item \textbf{Cell Biology:} Oncogenes, tumor suppressors, the cell cycle, and apoptosis are core cell biology concepts that are directly dysregulated in cancer.
  \item \textbf{Molecular Biology / Central Dogma:} Cancer-driving mutations occur in DNA and alter transcription and translation, changing the proteins a cell produces and how it behaves.
  \item \textbf{Immunology:} The immune system naturally surveils and eliminates abnormal cells; cancer evolves mechanisms to evade this surveillance. Immunotherapy exploits this relationship.
  \item \textbf{Evolutionary Biology:} Tumor progression follows Darwinian natural selection -- mutations that increase proliferation or survival are selected for within the body's microenvironment.
  \item \textbf{Pharmacology \& Chemistry:} Designing drugs that target cancer-specific molecules requires detailed knowledge of protein structure, enzyme kinetics, and receptor binding.
  \item \textbf{Epidemiology \& Public Health:} Identifying carcinogens (tobacco, UV radiation, HPV, asbestos) and behavioral risk factors enables prevention campaigns and evidence-based screening programs.
\end{itemize}
\end{tcolorbox}


%% ===========================================================================
%%  PART IIh – Protein Structure
%% ===========================================================================
\bluesections
\addcontentsline{toc}{section}{\textcolor{sectionblue}{PART IIh \textemdash\ Protein Structure}}
\addcontentsline{toc}{subsection}{The Four Levels of Protein Structure}
\addcontentsline{toc}{subsection}{Denaturation \& Renaturation}
\addcontentsline{toc}{subsection}{Structure Determines Function}
\partbanner{partbg}{sectionblue}{Part IIh \quad Protein Structure}

\vspace{8pt}

\section*{The Four Levels of Protein Structure}

Proteins are the most structurally and functionally diverse macromolecules in the cell.
Their function depends entirely on their \textbf{three-dimensional shape}, which arises
from four hierarchical levels of organization.

\noindent
{\renewcommand{\arraystretch}{1.6}
\begin{tabularx}{\textwidth}{>{\bfseries\centering\arraybackslash}p{1cm}
                              >{\bfseries\raggedright\arraybackslash}p{3cm}
                              >{\raggedright\arraybackslash}X
                              >{\raggedright\arraybackslash}p{3.2cm}}
\toprule
\rowcolor{tableheader}
\textbf{Level} & \textbf{Name} & \textbf{Description \& Bonds} & \textbf{Examples} \\
\midrule
\rowcolor{rowA} 1 & Primary & The linear sequence of amino acids. Determined directly by the gene (DNA). Held together by \textbf{peptide bonds} (covalent, very strong). & Any specific protein sequence \\
\rowcolor{rowB} 2 & Secondary & Local, repeating folding patterns of the polypeptide backbone driven by \textbf{hydrogen bonds} between backbone NH and C=O groups. Does \emph{not} involve R groups. & $\alpha$-helix; $\beta$-pleated sheet \\
\rowcolor{rowA} 3 & Tertiary & The overall \textbf{3D shape} of a single polypeptide chain. Driven by \textbf{R-group (side chain) interactions}: disulfide bridges (--S--S--), ionic bonds, hydrophobic interactions, and H-bonds. & Myoglobin; most enzymes \\
\rowcolor{rowB} 4 & Quaternary & Two or more polypeptide subunits assembled together. Held by the same forces as tertiary structure. \textbf{Not all proteins have this level.} & Hemoglobin (4 subunits); collagen (3 subunits) \\
\bottomrule
\end{tabularx}}

\vspace{4pt}
\begin{tcolorbox}[colback=craapbg, colframe=craapborder, arc=3pt, boxrule=1pt]
\textbf{Key concept --- peptide bonds:} When two amino acids join, the carboxyl group
($-$COOH) of one reacts with the amino group ($-$NH$_2$) of the next, releasing water
(\textbf{dehydration synthesis / condensation reaction}). The resulting covalent bond is a
\textbf{peptide bond}. Hydrolysis (addition of water) breaks it. The chain of amino acids
linked by peptide bonds is called a \textbf{polypeptide}.
\end{tcolorbox}

\section*{Denaturation \& Renaturation}

\textbf{Denaturation} is the loss of a protein's 3D structure (secondary, tertiary, and
quaternary levels) due to disruption of the bonds holding it together. The \textbf{primary
structure (amino acid sequence) is NOT broken} --- peptide bonds remain intact.

\textbf{Causes of denaturation:}
\begin{itemize}[noitemsep]
  \item \textbf{Heat:} Increased thermal energy disrupts hydrogen bonds and hydrophobic
        interactions (e.g., cooking an egg white denatures albumin --- irreversible).
  \item \textbf{Extreme pH:} Excess H$^+$ or OH$^-$ alters the charge on R groups, disrupting
        ionic bonds and hydrogen bonds (e.g., acid in the stomach denatures ingested proteins).
  \item \textbf{Chemical agents:} Urea, detergents, or heavy metals break non-covalent bonds.
\end{itemize}

\textbf{Renaturation:} Some proteins can refold correctly after the denaturing agent is
removed if conditions are mild --- demonstrating that the sequence alone encodes the shape.
\textit{Chaperone proteins} (heat shock proteins) assist in correct folding in the cell and
prevent misfolding aggregates.

\section*{Structure Determines Function}

\begin{tcolorbox}[colback=partbg, colframe=sectionblue, arc=3pt, boxrule=1pt]
The relationship between protein structure and function is absolute:
\begin{itemize}[noitemsep, leftmargin=*]
  \item An enzyme's \textbf{active site} is shaped to fit only specific substrates (induced
        fit model: active site adjusts slightly when substrate binds).
  \item A \textbf{single amino acid substitution} can destroy function (e.g., sickle cell
        disease: glutamic acid $\to$ valine at position 6 of hemoglobin's $\beta$-chain
        causes the protein to clump under low oxygen, deforming red blood cells).
  \item \textbf{Antibodies} recognize pathogens through precise structural complementarity.
  \item \textbf{Structural proteins} (collagen, keratin) derive their strength from specific
        quaternary arrangements (triple helices, coiled coils).
\end{itemize}
\end{tcolorbox}

\newpage

%% ===========================================================================
%%  PART IIi – DNA Replication
%% ===========================================================================
\orangesections
\addcontentsline{toc}{section}{\textcolor{sectionorange}{PART IIi \textemdash\ DNA Replication}}
\addcontentsline{toc}{subsection}{Semi-Conservative Replication}
\addcontentsline{toc}{subsection}{Enzymes of DNA Replication}
\addcontentsline{toc}{subsection}{Leading vs.\ Lagging Strand}
\addcontentsline{toc}{subsection}{Replication vs.\ Transcription Comparison}
\partbanner{part4bg}{sectionorange}{Part IIi \quad DNA Replication}

\vspace{8pt}

\section*{Semi-Conservative Replication}

Before a cell can divide, it must copy its entire genome. DNA replication is
\textbf{semi-conservative}: each new double helix consists of \textbf{one original (parental)
strand and one newly synthesized strand}.

\begin{tcolorbox}[colback=part4bg, colframe=sectionorange, arc=3pt, boxrule=1pt]
\textbf{Meselson--Stahl Experiment (1958):} Proved semi-conservative replication by growing
bacteria in heavy $^{15}$N medium, then switching to light $^{14}$N. After one generation,
all DNA had intermediate density (one heavy, one light strand per helix). After two generations,
half was intermediate and half was light. This ruled out conservative and dispersive models.
\end{tcolorbox}

Replication begins at specific sequences called \textbf{origins of replication}. In bacteria,
there is one origin; in eukaryotes, thousands of origins fire simultaneously, allowing the
entire genome to be copied quickly. Each origin creates two \textbf{replication forks} moving
outward in opposite directions (\textbf{bidirectional replication}).

\section*{Enzymes of DNA Replication}

\noindent
{\renewcommand{\arraystretch}{1.6}
\begin{tabularx}{\textwidth}{>{\bfseries\raggedright\arraybackslash}p{3.5cm}
                              >{\raggedright\arraybackslash}X}
\toprule
\rowcolor{tableheaderorange}
\textbf{Enzyme / Protein} & \textbf{Function} \\
\midrule
\rowcolor{rowOr} Helicase & Unwinds and separates the double helix by breaking hydrogen bonds between base pairs at the replication fork. \\
\rowcolor{rowA}  Topoisomerase & Relieves the torsional stress (positive supercoils) that builds up ahead of helicase as the DNA unwinds. Cuts and rejoins the DNA. \\
\rowcolor{rowOr} Single-Strand Binding Proteins (SSBPs) & Stabilize the separated single strands and prevent them from re-annealing. \\
\rowcolor{rowA}  Primase & Synthesizes a short RNA \textbf{primer} ($\sim$10 nucleotides) complementary to the template. Required because DNA polymerase cannot start a new chain from scratch --- it can only extend an existing strand. \\
\rowcolor{rowOr} DNA Polymerase III & The main replicating enzyme. Reads the template 3'$\to$5' and synthesizes the new strand 5'$\to$3' by adding deoxyribonucleotides. Also \textbf{proofreads}: removes and replaces mismatched bases. \\
\rowcolor{rowA}  DNA Polymerase I & Removes the RNA primers and replaces them with DNA nucleotides. \\
\rowcolor{rowOr} DNA Ligase & Seals the nick between adjacent Okazaki fragments on the lagging strand by forming a phosphodiester bond. Joins DNA fragments into a continuous strand. \\
\bottomrule
\end{tabularx}}

\section*{Leading vs.\ Lagging Strand}

DNA polymerase can only synthesize \textbf{5'$\to$3'}. Since the two template strands run
antiparallel, the two new strands must be made differently.

\begin{tcolorbox}[colback=part4bg, colframe=sectionorange, arc=3pt, boxrule=1pt]
\begin{multicols}{2}
\textbf{Leading strand (continuous):}
\begin{itemize}[noitemsep, leftmargin=*]
  \item Template runs 3'$\to$5' toward the fork.
  \item One RNA primer, then DNA pol III synthesizes continuously in the same direction the fork moves.
  \item Simple and fast.
\end{itemize}
\columnbreak
\textbf{Lagging strand (discontinuous):}
\begin{itemize}[noitemsep, leftmargin=*]
  \item Template runs 5'$\to$3' toward the fork.
  \item Synthesized in short fragments called \textbf{Okazaki fragments} (100--200 nt each in eukaryotes), each needing its own primer.
  \item Primers removed by DNA pol I; gaps filled; ligase joins fragments.
\end{itemize}
\end{multicols}
\end{tcolorbox}

\section*{Replication vs.\ Transcription --- Side-by-Side}

\noindent
{\renewcommand{\arraystretch}{1.5}
\begin{tabularx}{\textwidth}{>{\bfseries\raggedright\arraybackslash}p{3.5cm}
                              >{\raggedright\arraybackslash}X
                              >{\raggedright\arraybackslash}X}
\toprule
\rowcolor{tableheaderorange}
\textbf{Feature} & \textbf{DNA Replication} & \textbf{Transcription} \\
\midrule
\rowcolor{rowOr} Purpose & Copy the entire genome for cell division & Produce mRNA (or rRNA/tRNA) for protein synthesis \\
\rowcolor{rowA}  Template & Both strands (entire genome) & One strand (one gene at a time) \\
\rowcolor{rowOr} Enzyme & DNA Polymerase III & RNA Polymerase \\
\rowcolor{rowA}  Product & DNA (double-stranded) & RNA (single-stranded) \\
\rowcolor{rowOr} Primer needed? & Yes (RNA primer by primase) & No \\
\rowcolor{rowA}  Direction & 5'$\to$3' & 5'$\to$3' \\
\rowcolor{rowOr} When & S phase (before cell division) & Continuously, as needed \\
\bottomrule
\end{tabularx}}

\newpage

%% ===========================================================================
%%  PART IIj – Cell Cycle, Mitosis & Meiosis
%% ===========================================================================
\tealsections
\addcontentsline{toc}{section}{\textcolor{sectionteal}{PART IIj \textemdash\ Cell Cycle, Mitosis \& Meiosis}}
\addcontentsline{toc}{subsection}{The Cell Cycle \& Checkpoints}
\addcontentsline{toc}{subsection}{Mitosis --- PMAT}
\addcontentsline{toc}{subsection}{Meiosis --- Two Divisions}
\addcontentsline{toc}{subsection}{Mitosis vs.\ Meiosis Comparison}
\addcontentsline{toc}{subsection}{Key Vocabulary}
\partbanner{part5bg}{sectionteal}{Part IIj \quad Cell Cycle, Mitosis \& Meiosis}

\vspace{8pt}

\section*{The Cell Cycle \& Checkpoints}

The \textbf{cell cycle} is the ordered sequence of events a cell undergoes to grow and
divide. Most of the cycle is spent in \textbf{interphase} --- actively growing and preparing
--- not dividing.

\noindent
{\renewcommand{\arraystretch}{1.6}
\begin{tabularx}{\textwidth}{>{\bfseries\raggedright\arraybackslash}p{2.5cm}
                              >{\bfseries\raggedright\arraybackslash}p{2cm}
                              >{\raggedright\arraybackslash}X}
\toprule
\rowcolor{tableheaderteal}
\textbf{Phase} & \textbf{Sub-phase} & \textbf{What happens} \\
\midrule
\rowcolor{rowTe} Interphase & G$_1$ (Gap 1) & Cell grows; proteins synthesized; organelles duplicated; cell checks whether conditions are right to divide. \\
\rowcolor{rowA}  Interphase & S (Synthesis) & DNA replication --- each chromosome is duplicated into two identical sister chromatids joined at the centromere. \\
\rowcolor{rowTe} Interphase & G$_2$ (Gap 2) & Cell continues to grow; mitochondria and other organelles duplicated; cell prepares materials for division. \\
\rowcolor{rowA}  M Phase    & Mitosis + Cytokinesis & Nuclear division (PMAT) followed by cytoplasmic division. \\
\rowcolor{rowTe} --- & G$_0$ & Quiescent state; cells that are not actively dividing (e.g., fully differentiated neurons, muscle cells). \\
\bottomrule
\end{tabularx}}

\textbf{Cell Cycle Checkpoints} --- molecular ``stop signs'' that prevent a cell from
proceeding if something is wrong:
\begin{itemize}[noitemsep]
  \item \textbf{G$_1$/S checkpoint} (Restriction point): Is the cell large enough? Is DNA undamaged? Are growth factors present? The tumor suppressor \textbf{p53} monitors DNA damage here.
  \item \textbf{G$_2$/M checkpoint}: Did DNA replicate completely? Is any DNA damaged?
  \item \textbf{Spindle Assembly Checkpoint} (Metaphase): Are all chromosomes correctly attached to spindle fibers?
\end{itemize}

\begin{tcolorbox}[colback=part5bg, colframe=sectionteal, arc=3pt, boxrule=1pt]
\textbf{Cancer = checkpoint failure.} If a mutation disables a checkpoint gene (e.g., p53,
RB), the cell loses its ``brakes'' and divides uncontrollably. p53 is mutated in over 50\%
of all human cancers. Properly functioning checkpoints = tumor suppression.
\end{tcolorbox}

\section*{Mitosis --- PMAT}

Mitosis produces \textbf{two genetically identical daughter cells} from one parent cell,
each with the \textbf{same chromosome number} as the parent (diploid $\to$ diploid). Used
for: growth, tissue repair, asexual reproduction.

\begin{tcolorbox}[colback=part5bg, colframe=sectionteal,
  title={\bfseries PMAT: The Four Stages of Mitosis}, arc=3pt, boxrule=1pt]
\begin{enumerate}[leftmargin=*, noitemsep]
  \item \textbf{Prophase} --- Chromatin condenses into visible chromosomes (each consisting of
        two sister chromatids joined at the centromere). The mitotic spindle (made of
        microtubules) begins to form from the centrosomes. The nuclear envelope breaks down.
  \item \textbf{Metaphase} --- All chromosomes align along the cell's equator (the
        \textbf{metaphase plate}). Spindle fibers attach to the centromere of each chromosome
        via protein complexes called \textbf{kinetochores}. This is the spindle assembly
        checkpoint.
  \item \textbf{Anaphase} --- Sister chromatids \textbf{separate}: the centromere splits and
        the spindle fibers pull each chromatid toward opposite poles of the cell. Each pole
        now has a complete set of chromosomes.
  \item \textbf{Telophase} --- A nuclear envelope reforms around each set of chromosomes.
        Chromosomes begin to decondense back into chromatin. The spindle breaks down.
\end{enumerate}
\textbf{Cytokinesis} (occurs alongside telophase): the cytoplasm divides.
\begin{itemize}[noitemsep, leftmargin=*]
  \item \textbf{Animal cells:} A \textbf{cleavage furrow} pinches the cell in two (actin ring contracts).
  \item \textbf{Plant cells:} A \textbf{cell plate} forms down the middle, becoming a new cell wall.
\end{itemize}
\end{tcolorbox}

\section*{Meiosis --- Two Divisions}

Meiosis produces \textbf{four haploid ($n$) gametes} (sperm or eggs) from one diploid ($2n$) cell.
The chromosome number is halved. Without meiosis, fertilization would double the chromosome
number each generation.

\textbf{Meiosis consists of two sequential divisions (Meiosis I and II) with only one round
of DNA replication beforehand.}

\subsection*{Meiosis I --- Reductional Division (homologs separate)}

\begin{tcolorbox}[colback=part5bg, colframe=sectionteal, arc=3pt, boxrule=1pt]
\textbf{Prophase I} (the most important phase for genetic diversity):
\begin{itemize}[noitemsep, leftmargin=*]
  \item Homologous chromosomes (homologs) pair up side-by-side: called \textbf{synapsis}.
  \item \textbf{Crossing over (recombination)} occurs: non-sister chromatids of homologs
        exchange segments at points called \textbf{chiasmata}. This \textbf{reshuffles alleles}
        between chromosomes, generating new gene combinations not found in either parent.
\end{itemize}
\textbf{Metaphase I:} Homologous pairs align at the equator (not individual chromosomes as in
mitosis). The orientation of each pair is \textbf{random} --- this is \textbf{independent
assortment} (Mendel's 2nd law). With 23 pairs, humans can produce $2^{23}$ $\approx$ 8 million
different gamete combinations from this step alone.\\[4pt]
\textbf{Anaphase I:} \textbf{Homologs} separate to opposite poles (sister chromatids stay joined).\\[2pt]
\textbf{Telophase I + Cytokinesis:} Two haploid cells form, each with one of each homolog (but
each chromosome still consists of two sister chromatids).
\end{tcolorbox}

\subsection*{Meiosis II --- Equational Division (like mitosis)}

No DNA replication between Meiosis I and II.
\begin{itemize}[noitemsep]
  \item \textbf{Prophase II:} Spindle re-forms in each cell.
  \item \textbf{Metaphase II:} Individual chromosomes (not pairs) align at equator.
  \item \textbf{Anaphase II:} Sister chromatids separate and move to opposite poles.
  \item \textbf{Telophase II + Cytokinesis:} Four haploid daughter cells result.
\end{itemize}

\textbf{Sources of genetic variation in meiosis:}
\begin{enumerate}[noitemsep]
  \item \textbf{Crossing over} (Prophase I) --- shuffles alleles along chromosomes.
  \item \textbf{Independent assortment} (Metaphase I) --- random orientation of homolog pairs.
  \item \textbf{Random fertilization} --- any sperm can fertilize any egg.
\end{enumerate}

\subsection*{Nondisjunction \& Aneuploidy}

\textbf{Nondisjunction} = failure of chromosomes or sister chromatids to separate properly.
Results in gametes with the wrong number of chromosomes (\textbf{aneuploidy}).
\begin{itemize}[noitemsep]
  \item \textbf{Trisomy 21 (Down syndrome):} Extra chromosome 21 ($2n+1 = 47$). Most common
        chromosomal disorder; risk increases with maternal age.
  \item \textbf{Monosomy X (Turner syndrome):} 45,X --- one X chromosome in a female.
  \item \textbf{XXY (Klinefelter syndrome):} Extra X in a male.
\end{itemize}

\section*{Mitosis vs.\ Meiosis Comparison}

\noindent
{\renewcommand{\arraystretch}{1.5}
\begin{tabularx}{\textwidth}{>{\bfseries\raggedright\arraybackslash}p{4cm}
                              >{\raggedright\arraybackslash}X
                              >{\raggedright\arraybackslash}X}
\toprule
\rowcolor{tableheaderteal}
\textbf{Feature} & \textbf{Mitosis} & \textbf{Meiosis} \\
\midrule
\rowcolor{rowTe} Purpose & Growth, repair, asexual reproduction & Sexual reproduction; produce gametes \\
\rowcolor{rowA}  Divisions & 1 & 2 \\
\rowcolor{rowTe} Daughter cells & 2 & 4 \\
\rowcolor{rowA}  Ploidy of products & Diploid ($2n$) -- same as parent & Haploid ($n$) -- half of parent \\
\rowcolor{rowTe} Genetic identity & Identical to parent (barring mutation) & Unique (crossing over + ind. assortment) \\
\rowcolor{rowA}  Synapsis / crossing over & No & Yes (Prophase I) \\
\rowcolor{rowTe} Where in humans & All somatic (body) cells & Gonads (testes, ovaries) \\
\rowcolor{rowA}  Chromosome alignment at equator & Individual chromosomes & Homologous pairs (Meiosis I) \\
\bottomrule
\end{tabularx}}

what is the resulting amounts and when the chromosome number is halved in meiosis? The resulting amount of gametes is 4 and the chromosome number is halved in meiosis I.

\section*{Key Vocabulary}

\begin{itemize}[noitemsep]
  \item \textbf{Chromatin:} Loosely coiled DNA + histone proteins (during interphase).
  \item \textbf{Chromosome:} Highly condensed chromatin (visible during division).
  \item \textbf{Sister chromatids:} Two identical copies of a chromosome joined at the
        centromere after DNA replication; separated in Anaphase (mitosis) or Anaphase II
        (meiosis).
  \item \textbf{Homologous chromosomes (homologs):} Chromosome pairs --- one inherited from
        each parent. Same genes at same loci, but may have different alleles.
  \item \textbf{Diploid ($2n$):} Having two sets of chromosomes. Humans: $2n = 46$ (23 pairs).
  \item \textbf{Haploid ($n$):} Having one set. Human gametes: $n = 23$.
  \item \textbf{Centromere:} The constricted region joining sister chromatids; where spindle
        fibers (kinetochores) attach.
  \item \textbf{Karyotype:} A photograph of all chromosomes in a cell, arranged in homologous
        pairs by size.
\end{itemize}

\newpage


%% ===========================================================================
%%  PART II extra - Genetics \& Probability Logic
%% ===========================================================================
\section*{Practice: Genetics \& Probability Logic}

\begin{tcolorbox}[colback=part5bg, colframe=sectionteal, title={\bfseries Part 1: Calculating Gamete Diversity}, arc=3pt, boxrule=1pt]
The number of genetically different gametes is determined by the formula $2^n$, where \textbf{$n$ = the number of heterozygous gene pairs}.

\begin{multicols}{2}
\begin{itemize}[noitemsep, leftmargin=*]
    \item \textbf{aa:} $n=0 \to 2^0 =$ \textbf{1}
    \item \textbf{AaBB:} $n=1 \to 2^1 =$ \textbf{2}
    \item \textbf{AaBbCc:} $n=3 \to 2^3 =$ \textbf{8}
    \item \textbf{AaBbCcDdEe:} $n=5 \to 2^5 =$ \textbf{32}
    \item \textbf{AABB:} $n=0 \to 2^0 =$ \textbf{1}
    \item \textbf{AaBb:} $n=2 \to 2^2 =$ \textbf{4}
    \item \textbf{AabbCcDd:} $n=3 \to 2^3 =$ \textbf{8}
    \item \textbf{AAbbCCddEeFf:} $n=2 \to 2^2 =$ \textbf{4}
\end{itemize}
\end{multicols}
\end{tcolorbox}



\begin{tcolorbox}[colback=white, colframe=sectionteal, title={\bfseries Part 2: Multihybrid Cross Probabilities}, arc=3pt, boxrule=1pt]
\textbf{Parents:} (1) $AaBBCcddEe \times$ (2) $AaBbCcddEe$

\begin{enumerate}[leftmargin=*]
    \item \textbf{Gamete Count:}
    \begin{itemize}
        \item Individual 1 ($n=3$): $2^3 = 8$ gametes.
        \item Individual 2 ($n=4$): $2^4 = 16$ gametes.
    \end{itemize}
    
    \item \textbf{Prob. of Offspring \textit{AABBCCddee}:} \\
    Calculate individual gene probabilities and multiply (\textit{Product Rule}):
    \begin{itemize}[noitemsep]
        \item $Aa \times Aa \to AA = 1/4$
        \item $BB \times Bb \to BB = 1/2$
        \item $Cc \times Cc \to CC = 1/4$
        \item $dd \times dd \to dd = 1$
        \item $Ee \times Ee \to ee = 1/4$
    \end{itemize}
    \textbf{Total Probability:} $(1/4) \times (1/2) \times (1/4) \times 1 \times (1/4) = $ \textbf{1/128}

    \item \textbf{Prob. of Offspring \textit{AaBbCcDdEe}:} \\
    Look at the \textbf{D} gene: Parents are $dd \times dd$. It is impossible for them to produce a $Dd$ offspring. \\
    \textbf{Total Probability: 0}
\end{enumerate}
\end{tcolorbox}

\begin{tcolorbox}[colback=part5bg, colframe=sectionteal, title={\bfseries Part 3: Lethal Alleles \& Dihybrid Phenotypes}, arc=3pt, boxrule=1pt]
\textbf{The Manx Cat (Lethal Inheritance):}
\begin{itemize}[leftmargin=*]
    \item $M$ (Manx) is dominant; $m$ (Long tail) is recessive. $MM$ = Death.
    \item \textbf{Cross:} $Mm \times Mm$. The possible genotypes are $MM$ (dies), $Mm$, $Mm$, and $mm$.
    \item \textbf{Living Ratio:} Since $MM$ dies, we only count the 3 living kittens.
    \item \textbf{Result:} Probability of a long tail ($mm$) is \textbf{1/3}.
\end{itemize}



\textbf{Dihybrid Cat Cross:}
\textbf{Male:} $MmPp$ (Manx, Extra toes) $\times$ \textbf{Female:} $mmPp$ (Long tail, Extra toes)
\begin{itemize}[leftmargin=*]
    \item \textbf{Tail ($Mm \times mm$):} $1/2$ Manx, $1/2$ Long tail.
    \item \textbf{Toes ($Pp \times Pp$):} $3/4$ Extra toes, $1/4$ Normal toes.
\end{itemize}
\textbf{Phenotype Ratios:}
\begin{itemize}[noitemsep]
    \item Manx, Extra Toes: $(1/2) \times (3/4) = $ \textbf{3/8}
    \item Manx, Normal Toes: $(1/2) \times (1/4) = $ \textbf{1/8}
    \item Long Tail, Extra Toes: $(1/2) \times (3/4) = $ \textbf{3/8}
    \item Long Tail, Normal Toes: $(1/2) \times (1/4) = $ \textbf{1/8}
\end{itemize}
\end{tcolorbox}

\newpage

%% ===========================================================================
%%  PART IIk – Mendelian & Non-Mendelian Genetics
%% ===========================================================================
\purplesections
\addcontentsline{toc}{section}{\textcolor{sectionpurple}{PART IIk \textemdash\ Mendelian \& Non-Mendelian Genetics}}
\addcontentsline{toc}{subsection}{Essential Vocabulary}
\addcontentsline{toc}{subsection}{Mendel's Laws}
\addcontentsline{toc}{subsection}{Punnett Squares \& Ratios}
\addcontentsline{toc}{subsection}{Non-Mendelian Patterns}
\addcontentsline{toc}{subsection}{ABO Blood Types \& Pedigrees}
\partbanner{part3bg}{sectionpurple}{Part IIk \quad Mendelian \& Non-Mendelian Genetics}

\vspace{8pt}

\section*{Essential Vocabulary}

\noindent
{\renewcommand{\arraystretch}{1.5}
\begin{tabularx}{\textwidth}{>{\bfseries\raggedright\arraybackslash}p{3.5cm}
                              >{\raggedright\arraybackslash}X}
\toprule
\rowcolor{tableheaderpurple}
\textbf{Term} & \textbf{Definition} \\
\midrule
\rowcolor{rowPu} Gene            & A segment of DNA on a chromosome that codes for a specific trait or protein. \\
\rowcolor{rowA}  Allele          & Alternative versions of a gene (e.g., \textit{B} = brown eyes; \textit{b} = blue eyes). \\
\rowcolor{rowPu} Locus (pl.\ loci)& The specific physical location of a gene on a chromosome. \\
\rowcolor{rowA}  Genotype        & The actual allele combination an organism carries (e.g., \textit{Bb}, \textit{BB}, \textit{bb}). \\
\rowcolor{rowPu} Phenotype       & The observable physical or biochemical expression of the genotype (e.g., brown eyes). \\
\rowcolor{rowA}  Homozygous      & Both alleles at a locus are identical (\textit{BB} or \textit{bb}). \\
\rowcolor{rowPu} Heterozygous    & The two alleles at a locus differ (\textit{Bb}). Also called a \emph{carrier} for recessive alleles. \\
\rowcolor{rowA}  Dominant        & The allele whose phenotype is expressed whenever at least one copy is present (written as capital letter). \\
\rowcolor{rowPu} Recessive       & The allele whose phenotype is expressed \emph{only} when the organism is homozygous recessive (\textit{bb}). \\
\rowcolor{rowA}  Wild type       & The most common allele or phenotype found in a natural population; used as a reference. \\
\rowcolor{rowPu} Testcross       & Crossing an unknown genotype with a homozygous recessive (\textit{aa}) to determine the unknown's genotype. \\
\bottomrule
\end{tabularx}}

\begin{tcolorbox}[colback=craapbg, colframe=craapborder, arc=3pt, boxrule=1pt]
\textbf{Genotype vs.\ phenotype is not 1:1.} A brown-eyed person (\textit{Bb}) and a brown-eyed
person (\textit{BB}) have the same phenotype but different genotypes. The dominant allele
``masks'' the recessive. Two blue-eyed (\textit{bb}) parents cannot have a brown-eyed child
--- but two \textit{Bb} parents can have a blue-eyed (\textit{bb}) child.
\end{tcolorbox}

\section*{Mendel's Laws}

Gregor Mendel's pea plant experiments (1860s) established the quantitative rules of inheritance.

\begin{tcolorbox}[colback=part3bg, colframe=sectionpurple, arc=3pt, boxrule=1pt]
\textbf{Law of Segregation:} Each organism has two alleles for each trait. The two alleles
separate (segregate) during gamete formation so each gamete carries \textbf{only one allele}.
The pair reunites randomly at fertilization. \\[4pt]
\textbf{Law of Independent Assortment:} The alleles of two different genes assort independently
during gamete formation --- the inheritance of one gene does \emph{not} affect another, provided
the genes are on \textbf{different chromosomes} (or far apart on the same chromosome). This
explains why traits are inherited independently.
\end{tcolorbox}

\section*{Punnett Squares \& Ratios}

\subsection*{Monohybrid Cross (one gene, two heterozygous parents: \textit{Aa} $\times$ \textit{Aa})}
\begin{center}
\renewcommand{\arraystretch}{1.4}
\begin{tabular}{c|cc}
         & \textit{A}  & \textit{a}  \\
\hline
\textit{A} & \textit{AA} & \textit{Aa} \\
\textit{a} & \textit{Aa} & \textit{aa} \\
\end{tabular}
\end{center}
Genotype ratio: 1 \textit{AA} : 2 \textit{Aa} : 1 \textit{aa}\quad
Phenotype ratio: \textbf{3 dominant : 1 recessive}

\subsection*{Dihybrid Cross (\textit{AaBb} $\times$ \textit{AaBb})}
With two independently assorting genes, the phenotype ratio of the offspring is
\textbf{9 : 3 : 3 : 1}:
\begin{center}
9 A\_ B\_ (both dominant) : 3 A\_ \textit{bb} : 3 \textit{aa} B\_ : 1 \textit{aa bb}
\end{center}

\subsection*{Testcross}
Cross the unknown phenotype with homozygous recessive (\textit{aa}):
\begin{itemize}[noitemsep]
  \item All offspring show dominant phenotype $\to$ unknown parent was \textit{AA}.
  \item Half dominant, half recessive $\to$ unknown parent was \textit{Aa}.
\end{itemize}

\section*{Non-Mendelian Patterns of Inheritance}

\noindent
{\renewcommand{\arraystretch}{1.6}
\begin{tabularx}{\textwidth}{>{\bfseries\raggedright\arraybackslash}p{4cm}
                              >{\raggedright\arraybackslash}X
                              >{\raggedright\arraybackslash}p{3.5cm}}
\toprule
\rowcolor{tableheaderpurple}
\textbf{Pattern} & \textbf{Description} & \textbf{Example} \\
\midrule
\rowcolor{rowPu} Incomplete Dominance & Heterozygote shows an \textbf{intermediate} phenotype --- neither allele is fully dominant. & Red $\times$ White snapdragons $\to$ Pink offspring \\
\rowcolor{rowA}  Codominance & \textbf{Both} alleles are fully and simultaneously expressed in the heterozygote. & Type AB blood (both A and B antigens present) \\
\rowcolor{rowPu} Multiple Alleles & More than two alleles exist for a gene in the population (though each individual still has only two). & ABO blood type: $I^A$, $I^B$, $i$ \\
\rowcolor{rowA}  Sex-Linked (X-linked) & Gene is on the X chromosome. Males (XY) have only one copy, so they express recessive X-linked traits more frequently. & Color blindness, hemophilia \\
\rowcolor{rowPu} Polygenic Inheritance & \textbf{Multiple genes} influence one trait, producing a continuous distribution (bell curve). & Height, skin color, intelligence \\
\rowcolor{rowA}  Pleiotropy & \textbf{One gene} affects \textbf{multiple} phenotypic traits. & Sickle cell anemia (one mutation $\to$ anemia, joint pain, organ damage, malaria resistance) \\
\rowcolor{rowPu} Epistasis & One gene masks or modifies the expression of another gene. & Coat color in Labrador retrievers (B/b controls color; E/e controls whether pigment is deposited) \\
\bottomrule
\end{tabularx}}

\section*{ABO Blood Types \& Pedigrees}

ABO blood type is the classic example of both \textbf{multiple alleles} and
\textbf{codominance}. Three alleles exist: $I^A$, $I^B$, and $i$ (recessive).

\noindent
{\renewcommand{\arraystretch}{1.4}
\footnotesize\begin{tabularx}{\textwidth}{>{\bfseries\centering\arraybackslash}p{2.8cm}
                              >{\centering\arraybackslash}p{1.5cm}
                              >{\centering\arraybackslash}p{2.5cm}
                              >{\centering\arraybackslash}p{2.8cm}
                              >{\centering\arraybackslash}p{2.8cm}}
\toprule
\rowcolor{tableheaderpurple}
\textbf{Genotype} & \textbf{Type} & \textbf{Antigen} & \textbf{Donate to} & \textbf{Receive from} \\
\midrule
\rowcolor{rowPu} $I^A I^A$ or $I^A i$ & A & A & A, AB & A, O \\
\rowcolor{rowA}  $I^B I^B$ or $I^B i$ & B & B & B, AB & B, O \\
\rowcolor{rowPu} $I^A I^B$             & AB & A and B & AB only & All (universal recipient) \\
\rowcolor{rowA}  $ii$                  & O & None & All (universal donor) & O only \\
\bottomrule
\end{tabularx}}

\vspace{4pt}
\textbf{Pedigree analysis rules:}
\begin{itemize}[noitemsep]
  \item Squares = males; circles = females; filled = affected; half-filled = carrier.
  \item \textbf{Autosomal dominant:} appears every generation; males and females equally affected; affected parent in every generation.
  \item \textbf{Autosomal recessive:} can skip generations; two unaffected parents can have an affected child (both carriers); more common when related parents mate.
  \item \textbf{X-linked recessive:} more common in males; affected males often have unaffected carrier mothers; no father-to-son transmission (fathers give Y to sons).
\end{itemize}

\newpage

%% ===========================================================================
%%  PART IIl – Evolution & Natural Selection
%% ===========================================================================
\greensections
\addcontentsline{toc}{section}{\textcolor{sectiongreen}{PART IIl \textemdash\ Evolution \& Natural Selection}}
\addcontentsline{toc}{subsection}{Darwin's Key Observations}
\addcontentsline{toc}{subsection}{Natural Selection \& Fitness}
\addcontentsline{toc}{subsection}{Types of Selection}
\addcontentsline{toc}{subsection}{Five Mechanisms of Evolution}
\addcontentsline{toc}{subsection}{Hardy-Weinberg Equilibrium}
\addcontentsline{toc}{subsection}{Speciation \& Evidence for Evolution}
\partbanner{part2bg}{sectiongreen}{Part IIl \quad Evolution \& Natural Selection}

\vspace{8pt}

\section*{Darwin's Key Observations}

Charles Darwin's theory of evolution by natural selection (1859) is the unifying framework of
all biology. It is built on four observations:
\begin{enumerate}[noitemsep]
  \item \textbf{Overproduction:} Organisms produce more offspring than can survive given limited resources.
  \item \textbf{Variation:} Individuals in a population differ in heritable traits.
  \item \textbf{Heritability:} Traits are passed from parents to offspring through genes.
  \item \textbf{Differential survival:} Individuals with traits better suited to the current
        environment survive and reproduce more (differential reproductive success).
\end{enumerate}

\begin{tcolorbox}[colback=part2bg, colframe=sectiongreen, arc=3pt, boxrule=1pt]
\textbf{Evolution is NOT:}
\begin{itemize}[noitemsep, leftmargin=*]
  \item Goal-directed or progressive --- it doesn't make organisms ``better'' in any absolute sense.
  \item Acting on individuals --- selection acts on \emph{phenotype} (individuals), but
        \emph{evolution} (change in allele frequency) occurs in \textbf{populations} over generations.
  \item Equivalent to ``survival of the fittest'' in the popular sense --- \textbf{fitness} means
        \textbf{reproductive success}, not physical strength.
\end{itemize}
\end{tcolorbox}

\section*{Natural Selection \& Fitness}

\textbf{Fitness} = an organism's \textbf{relative reproductive success} --- how many viable,
fertile offspring it contributes to the next generation compared to others in the population.

\textbf{Adaptation} = a heritable trait that \textbf{increases fitness} in a given environment.
Adaptations arise by natural selection acting on existing genetic variation; they are not
created on demand.

\section*{Types of Natural Selection}

Natural selection can shift the distribution of a trait in a population in three ways.
Imagine a bell curve of a quantitative trait (e.g., body size):

\noindent
{\renewcommand{\arraystretch}{1.6}
\begin{tabularx}{\textwidth}{>{\bfseries\raggedright\arraybackslash}p{3.5cm}
                              >{\raggedright\arraybackslash}X
                              >{\raggedright\arraybackslash}p{3.5cm}}
\toprule
\rowcolor{tableheadergreen}
\textbf{Type} & \textbf{Description} & \textbf{Example} \\
\midrule
\rowcolor{rowA}  Directional & Favors one extreme; \textbf{shifts} the mean of the distribution. Reduces variation over time in one direction. & Antibiotic resistance (bacteria with mutations survive); beak elongation in a drought \\
\rowcolor{rowBg} Stabilizing & Favors the \textbf{middle/average}; eliminates both extremes. \textbf{Reduces} variation; most common type in stable environments. & Human birth weight (very small or very large babies have higher mortality) \\
\rowcolor{rowA}  Disruptive  & Favors \textbf{both extremes} simultaneously; eliminates the middle. \textbf{Increases} variation; can lead to speciation. & Seed size in a population of birds with bimodal diet (small and large seeds dominate) \\
\bottomrule
\end{tabularx}}

\section*{Five Mechanisms of Evolution}

Natural selection is only one of five mechanisms that change allele frequencies in a population.

\begin{tcolorbox}[colback=part2bg, colframe=sectiongreen, arc=3pt, boxrule=1pt]
\begin{enumerate}[noitemsep, leftmargin=*]
  \item \textbf{Natural selection:} Non-random differential survival based on fitness. Adaptive evolution.
  \item \textbf{Genetic drift:} \emph{Random} changes in allele frequency due to chance sampling errors. Most powerful in \textbf{small populations}.
    \begin{itemize}[noitemsep]
      \item \textbf{Bottleneck effect:} Population crash drastically reduces gene pool. Survivors may not represent original diversity (e.g., cheetahs --- very low genetic diversity from ancient bottleneck).
      \item \textbf{Founder effect:} A few individuals colonize a new area. New population has less diversity and may have allele frequencies very different from the source population.
    \end{itemize}
  \item \textbf{Gene flow:} Movement of individuals (and their alleles) between populations. Reduces differences between populations; prevents speciation.
  \item \textbf{Mutation:} The \textbf{ultimate source of all new alleles}. Random; not directed. Most are neutral or harmful; rarely beneficial. Without mutation, there would be no variation for selection to act on.
  \item \textbf{Sexual selection:} A form of natural selection based on \textbf{mate choice}. Can lead to exaggerated traits (peacock tail, elk antlers) that increase mating success even at survival cost.
\end{enumerate}
\end{tcolorbox}

\section*{Hardy-Weinberg Equilibrium}

The \textbf{Hardy-Weinberg principle} states that allele and genotype frequencies in a
population remain \textbf{constant} from generation to generation if \emph{no} evolutionary
forces are acting. It describes a null model --- a population that is \emph{not} evolving.

\textbf{Five conditions required (all must be true):}
\begin{enumerate}[noitemsep]
  \item No natural selection
  \item No mutation
  \item No gene flow (no migration)
  \item Random mating
  \item Very large population (no genetic drift)
\end{enumerate}

\[
  p + q = 1 \quad \text{(allele frequencies sum to 1)}
\]
\[
  p^2 + 2pq + q^2 = 1 \quad \text{(genotype frequencies sum to 1)}
\]

where $p$ = frequency of dominant allele, $q$ = frequency of recessive allele; $p^2$ =
homozygous dominant; $2pq$ = heterozygous; $q^2$ = homozygous recessive.

\begin{tcolorbox}[colback=part2bg, colframe=sectiongreen, arc=3pt, boxrule=1pt]
\textbf{Exam problem type:} If 9 out of 100 people have cystic fibrosis (autosomal recessive):
$q^2 = 0.09 \Rightarrow q = 0.3 \Rightarrow p = 0.7 \Rightarrow 2pq = 2(0.7)(0.3) = 0.42$
So 42\% of the population are carriers.
\end{tcolorbox}

\section*{Speciation \& Evidence for Evolution}

\textbf{Species:} A group of organisms that can interbreed and produce \textbf{fertile offspring}
in nature. Reproductive isolation is what creates and maintains species.

\noindent
{\renewcommand{\arraystretch}{1.5}
\begin{tabularx}{\textwidth}{>{\bfseries\raggedright\arraybackslash}p{4cm}
                              >{\raggedright\arraybackslash}X}
\toprule
\rowcolor{tableheadergreen}
\textbf{Type of Speciation} & \textbf{Description \& Example} \\
\midrule
\rowcolor{rowA} Allopatric & Geographic barrier \textbf{physically separates} a population. Isolated populations diverge under different selective pressures. \textit{Example: Grand Canyon squirrels, Darwin's finches.} \\
\rowcolor{rowBg} Sympatric & Speciation without geographic isolation. Occurs through polyploidy (common in plants), habitat differentiation, or assortative mating. \textit{Example: polyploid plants, cichlid fish in African lakes.} \\
\bottomrule
\end{tabularx}}

\vspace{6pt}
\textbf{Evidence for Evolution:}
\begin{itemize}[noitemsep]
  \item \textbf{Fossil record:} Shows progressive change over geological time; transitional forms (e.g., Tiktaalik, fish-tetrapod link).
  \item \textbf{Homologous structures:} Same underlying anatomy, different functions in different species (e.g., human arm, whale flipper, bat wing) --- evidence of common descent.
  \item \textbf{Analogous structures:} Similar function but different evolutionary origin (e.g., insect wing vs.\ bird wing) --- evidence of \emph{convergent evolution}, not common ancestry.
  \item \textbf{Vestigial structures:} Reduced, non-functional remnants of structures that were functional in ancestors (e.g., human coccyx, whale pelvis, cave fish eyes).
  \item \textbf{Molecular evidence:} DNA and protein sequences are more similar in closely related species. All life uses the same genetic code (evidence of universal common ancestor).
  \item \textbf{Biogeography:} Species distributions make sense in light of continental drift and evolution (e.g., marsupials concentrated in Australia).
\end{itemize}

\newpage

%% ===========================================================================
%%  PART IIm – Classification, Taxonomy & Ecology
%% ===========================================================================
\bluesections
\addcontentsline{toc}{section}{\textcolor{sectionblue}{PART IIm \textemdash\ Classification, Taxonomy \& Ecology}}
\addcontentsline{toc}{subsection}{Taxonomic Hierarchy \& Binomial Nomenclature}
\addcontentsline{toc}{subsection}{The Three Domains}
\addcontentsline{toc}{subsection}{Levels of Ecological Organization}
\addcontentsline{toc}{subsection}{Population Ecology}
\addcontentsline{toc}{subsection}{Community Ecology \& Symbiosis}
\addcontentsline{toc}{subsection}{Energy Flow \& Trophic Levels}
\addcontentsline{toc}{subsection}{Biogeochemical Cycles}
\partbanner{partbg}{sectionblue}{Part IIm \quad Classification, Taxonomy \& Ecology}

\vspace{8pt}

\section*{Taxonomic Hierarchy \& Binomial Nomenclature}

\textbf{Taxonomy} is the science of naming and classifying organisms. Carl Linnaeus devised
the hierarchical system still used today.

\textbf{Hierarchy (broadest $\to$ most specific):}
\[
  \text{Domain} \to \text{Kingdom} \to \text{Phylum} \to \text{Class} \to
  \text{Order} \to \text{Family} \to \text{Genus} \to \text{Species}
\]
\textit{Memory: \textbf{D}ear \textbf{K}ing \textbf{P}hilip \textbf{C}ame \textbf{O}ver
\textbf{F}or \textbf{G}ood \textbf{S}oup.}

\textbf{Binomial nomenclature:} Each species has a two-part Latin name:
\textit{Genus species} (e.g., \textit{Homo sapiens}). Rules:
\begin{itemize}[noitemsep]
  \item Always \textbf{italicized} (or underlined when handwriting).
  \item Genus is \textbf{capitalized}; species is \textbf{lowercase}.
  \item Provides an unambiguous universal name understood by all biologists worldwide.
\end{itemize}

\textbf{Phylogenetic trees (cladograms):} Diagrams showing evolutionary relationships.
Nodes represent \emph{common ancestors}; branches represent lineages. A \textbf{clade}
includes an ancestor and \emph{all} of its descendants.

\section*{The Three Domains}

The \textbf{Three-Domain System} (Woese, 1990) divides all life based on rRNA sequences:

\noindent
{\renewcommand{\arraystretch}{1.6}
\begin{tabularx}{\textwidth}{>{\bfseries\centering\arraybackslash}p{2.5cm}
                              >{\centering\arraybackslash}p{2cm}
                              >{\raggedright\arraybackslash}X}
\toprule
\rowcolor{tableheader}
\textbf{Domain} & \textbf{Cell Type} & \textbf{Key Characteristics \& Examples} \\
\midrule
\rowcolor{rowA} Bacteria & Prokaryote & Peptidoglycan cell walls; most diverse group; includes pathogens and beneficial bacteria; first life forms on Earth. \\
\rowcolor{rowB} Archaea  & Prokaryote & \textbf{No} peptidoglycan; found in extreme environments (thermophiles, halophiles, methanogens); more closely related to Eukarya than to Bacteria. \\
\rowcolor{rowA} Eukarya  & Eukaryote  & Membrane-bound organelles; contains all four kingdoms below. \\
\bottomrule
\end{tabularx}}

\textbf{Four Kingdoms of Eukarya:}
\begin{itemize}[noitemsep]
  \item \textbf{Animalia:} Multicellular; heterotrophic (ingest food); no cell wall; motile.
  \item \textbf{Plantae:} Multicellular; autotrophic (photosynthesis); cellulose cell walls.
  \item \textbf{Fungi:} Mostly multicellular; heterotrophic (absorb nutrients by secreting
        digestive enzymes externally); \textbf{chitin} cell walls; decomposers.
  \item \textbf{Protista:} Mostly unicellular eukaryotes; catch-all kingdom for organisms
        that don't fit elsewhere (algae, amoeba, paramecium).
\end{itemize}

\section*{Levels of Ecological Organization}

\[
  \text{Organism} \to \text{Population} \to \text{Community} \to
  \text{Ecosystem} \to \text{Biome} \to \text{Biosphere}
\]
\begin{itemize}[noitemsep]
  \item \textbf{Population:} All individuals of the same species in a given area.
  \item \textbf{Community:} All populations of different species living and interacting in
        an area.
  \item \textbf{Ecosystem:} Community + the abiotic (non-living) environment (soil, water, climate).
  \item \textbf{Biome:} Large geographic region characterized by a specific climate and
        community type (tundra, rainforest, grassland, desert, ocean).
  \item \textbf{Biosphere:} All regions of Earth inhabited by life.
\end{itemize}

\section*{Population Ecology}

\subsection*{Growth Models}

\begin{tcolorbox}[colback=partbg, colframe=sectionblue, arc=3pt, boxrule=1pt]
\begin{multicols}{2}
\textbf{Exponential Growth (J-curve):}
\[\frac{dN}{dt} = rN\]
Unlimited resources; population doubles at regular intervals.
\textit{Real examples:} bacteria in ideal conditions, invasive species early in colonization.
\columnbreak
\textbf{Logistic Growth (S-curve):}
\[\frac{dN}{dt} = rN\left(\frac{K-N}{K}\right)\]
Growth slows as population approaches \textbf{carrying capacity $K$} (maximum sustainable population given available resources). Rate of growth is fastest at $N = K/2$.
\end{multicols}
\end{tcolorbox}

\begin{itemize}[noitemsep]
  \item \textbf{$r$-strategists:} High reproductive rate; many small offspring; little parental care; short lifespan. Thrive in unstable environments. \textit{Examples:} insects, mice, bacteria.
  \item \textbf{$K$-strategists:} Low reproductive rate; few large offspring; much parental care; long lifespan. Thrive in stable environments near $K$. \textit{Examples:} elephants, humans, eagles.
\end{itemize}

\section*{Community Ecology \& Symbiosis}

\noindent
{\renewcommand{\arraystretch}{1.6}
\begin{tabularx}{\textwidth}{>{\bfseries\raggedright\arraybackslash}p{3.5cm}
                              >{\centering\arraybackslash}p{2cm}
                              >{\raggedright\arraybackslash}X}
\toprule
\rowcolor{tableheader}
\textbf{Interaction Type} & \textbf{Effect (A / B)} & \textbf{Description \& Example} \\
\midrule
\rowcolor{rowA}  Mutualism (symbiosis)  & $+$ / $+$ & Both species benefit. \textit{Mycorrhizal fungi + plant roots (fungi get sugars; plant gets minerals); nitrogen-fixing Rhizobium in legume roots.} \\
\rowcolor{rowB}  Commensalism          & $+$ / 0   & One benefits; other is unaffected. \textit{Barnacles on whale skin; birds following cattle.} \\
\rowcolor{rowA}  Parasitism            & $+$ / $-$ & Parasite benefits; host harmed (but usually not killed immediately). \textit{Tapeworms, ticks, Plasmodium (malaria).} \\
\rowcolor{rowB}  Predation             & $+$ / $-$ & Predator kills and eats prey. Drives co-evolution of offense and defense adaptations. \\
\rowcolor{rowA}  Competition           & $-$ / $-$ & Both species harmed (resources depleted). Interspecific: between species. Intraspecific: within species. \\
\bottomrule
\end{tabularx}}

\textbf{Keystone species:} A species with disproportionately large ecological impact relative
to its abundance. Removing it causes cascading changes throughout the ecosystem
(e.g., sea otters control sea urchin populations, which protects kelp forests).

\textbf{Ecological succession:} Gradual change in community composition over time.
\begin{itemize}[noitemsep]
  \item \textbf{Primary succession:} Begins on bare rock with no soil (e.g., after a lava
        flow). Pioneers: \emph{lichens} (break down rock $\to$ soil), then mosses,
        shrubs, trees. Slow process (hundreds to thousands of years).
  \item \textbf{Secondary succession:} Begins on disturbed land that still has soil (e.g.,
        after a forest fire or abandoned field). Faster recovery because soil remains.
\end{itemize}

\section*{Energy Flow \& Trophic Levels}

\textbf{Trophic levels} describe feeding position in a food chain:

\[
  \underbrace{\text{Producers}}_{\text{Level 1}} \to
  \underbrace{\text{Primary Consumers}}_{\text{Level 2 (herbivores)}} \to
  \underbrace{\text{Secondary Consumers}}_{\text{Level 3}} \to
  \underbrace{\text{Tertiary Consumers}}_{\text{Level 4}}
\]

\begin{tcolorbox}[colback=partbg, colframe=sectionblue, arc=3pt, boxrule=1pt]
\textbf{The 10\% Rule:} Only approximately \textbf{10\%} of the energy at one trophic level
is transferred to the next. The other 90\% is lost as heat (metabolic processes, cellular
respiration, movement). This is why:
\begin{itemize}[noitemsep, leftmargin=*]
  \item Food chains rarely have more than 4--5 trophic levels (too little energy remains).
  \item Eating lower on the food chain is more energy-efficient (why plant-based diets have
        a smaller ecological footprint).
  \item Biomass pyramids narrow sharply: it takes $\sim$10{,}000 kg of grass to support
        $\sim$1{,}000 kg of grasshoppers to support $\sim$100 kg of frogs to support
        $\sim$10 kg of snakes.
\end{itemize}
\end{tcolorbox}

\section*{Biogeochemical Cycles}

\textbf{Carbon cycle:} CO$_2$ removed from atmosphere by \textbf{photosynthesis} (stored in
organic molecules); returned by \textbf{cellular respiration}, decomposition, and combustion
of fossil fuels. Burning fossil fuels has elevated atmospheric CO$_2$, driving climate change.

\textbf{Nitrogen cycle:} Nitrogen gas (N$_2$) makes up 78\% of atmosphere but is unusable
directly by most organisms. Key steps:
\begin{itemize}[noitemsep]
  \item \textbf{Nitrogen fixation:} Bacteria (\textit{Rhizobium} in legume root nodules;
        free-living \textit{Azotobacter}) convert N$_2$ $\to$ NH$_3$/NH$_4^+$.
        Also occurs via lightning.
  \item \textbf{Nitrification:} NH$_4^+$ $\to$ NO$_2^-$ $\to$ NO$_3^-$ (by nitrifying
        bacteria). Plants can absorb NO$_3^-$.
  \item \textbf{Assimilation:} Plants incorporate nitrogen into amino acids and nucleotides;
        moves up food chain.
  \item \textbf{Decomposition:} Decomposers return N from organic matter to NH$_4^+$.
  \item \textbf{Denitrification:} Bacteria convert NO$_3^-$ back to N$_2$, returning it
        to the atmosphere.
\end{itemize}



\newpage

%% =============================================================================
%%  PART IIn -- Biological Cycles & Pathways Reference
%% =============================================================================
\tealsections
\addcontentsline{toc}{section}{\textcolor{sectionteal}{PART IIn \textemdash\ Biological Cycles \& Pathways Reference}}
\addcontentsline{toc}{subsection}{Cell Cycle: G1, S, G2, M}
\addcontentsline{toc}{subsection}{Mitosis (PMAT)}
\addcontentsline{toc}{subsection}{Meiosis I \& II}
\addcontentsline{toc}{subsection}{Glycolysis (Cytoplasm)}
\addcontentsline{toc}{subsection}{Pyruvate Oxidation \& Krebs Cycle}
\addcontentsline{toc}{subsection}{Electron Transport Chain \& Chemiosmosis}
\addcontentsline{toc}{subsection}{Light-Dependent Reactions}
\addcontentsline{toc}{subsection}{Calvin Cycle}
\addcontentsline{toc}{subsection}{C3, C4 \& CAM Photosynthesis}
\addcontentsline{toc}{subsection}{ATP Production Across Organisms}
\addcontentsline{toc}{subsection}{Fermentation Pathways}
\addcontentsline{toc}{subsection}{Nitrogen \& Carbon Cycles}
\partbanner{part5bg}{sectionteal}{Part IIn \quad Biological Cycles \& Pathways Reference}

\vspace{6pt}
\begin{center}\small\itshape
  All the cycles, phases, and pathways from earlier sections in one place.
  Use this as a quick lookup when you need to remember the inputs, outputs,
  or location of any biological process.
\end{center}

\section*{Cell Cycle: G1, S, G2, M}

The \textbf{cell cycle} describes the life of a dividing cell --- how it grows, copies its
DNA, and divides into two daughter cells. Most of the cycle ($\sim$90\%) is spent in
\textbf{interphase}; the actual division (\textbf{M phase}) is brief.

\noindent
{\renewcommand{\arraystretch}{1.6}
\begin{tabularx}{\textwidth}{>{\bfseries\centering\arraybackslash}p{1.5cm}
                              >{\bfseries\raggedright\arraybackslash}p{3cm}
                              >{\raggedright\arraybackslash}X}
\toprule
\rowcolor{tableheaderteal}
\textbf{Phase} & \textbf{Name} & \textbf{What happens} \\
\midrule
\rowcolor{rowTe} G\textsubscript{1} & Gap 1 / Growth 1 & Cell grows in size, makes proteins and organelles, performs normal functions. The longest, most variable phase. The cell ``decides'' here whether to divide. \\
\rowcolor{rowA}  \textbf{S} & \textbf{Synthesis} & \textbf{DNA is replicated.} Every chromosome makes an identical copy --- now exists as two sister chromatids joined at the centromere. The cell now has \emph{double} its normal DNA content. \\
\rowcolor{rowTe} G\textsubscript{2} & Gap 2 / Growth 2 & More growth, more protein synthesis. Cell prepares the spindle apparatus and checks DNA for damage before dividing. \\
\rowcolor{rowA}  M & Mitosis + Cytokinesis & The actual division. Sister chromatids separate; the cytoplasm splits into two daughter cells. \\
\rowcolor{rowTe} G\textsubscript{0} & Resting / Quiescent & Cells that have exited the cycle (most adult neurons, muscle cells). Some can re-enter G\textsubscript{1} when needed (liver cells after damage); others cannot. \\
\bottomrule
\end{tabularx}}

\textbf{Cell cycle checkpoints} (where the cell verifies everything is OK before proceeding):
\begin{itemize}[noitemsep]
  \item \textbf{G\textsubscript{1}/S checkpoint (the ``restriction point''):} Is DNA undamaged? Is environment favorable? Are nutrients available? \textbf{p53} is the master regulator here --- mutated p53 is found in $\sim$50\% of human cancers.
  \item \textbf{G\textsubscript{2}/M checkpoint:} Was DNA replicated correctly and completely? Is it undamaged?
  \item \textbf{M (spindle) checkpoint:} Are all chromosomes attached to spindle fibers from both poles? Prevents nondisjunction.
\end{itemize}

\section*{Mitosis (PMAT)}

\textbf{Mitosis} is the division of a single nucleus into two genetically identical
daughter nuclei. Used for growth, tissue repair, and asexual reproduction. \textbf{1 parent
cell $\to$ 2 identical diploid daughter cells.}

\noindent
{\renewcommand{\arraystretch}{1.6}
\begin{tabularx}{\textwidth}{>{\bfseries\raggedright\arraybackslash}p{2.5cm}
                              >{\raggedright\arraybackslash}X}
\toprule
\rowcolor{tableheaderteal}
\textbf{Phase} & \textbf{What happens} \\
\midrule
\rowcolor{rowTe} \textbf{P}rophase & Chromatin \textbf{condenses} into visible chromosomes (each as two sister chromatids). Nuclear envelope breaks down. Spindle fibers form from centrosomes at opposite poles. \\
\rowcolor{rowA}  \textbf{M}etaphase & Chromosomes \textbf{align} at the cell's equator (the metaphase plate). Each sister chromatid is attached to a spindle fiber from opposite poles. \\
\rowcolor{rowTe} \textbf{A}naphase & Sister chromatids \textbf{separate} (centromeres split) and are pulled to opposite poles by shortening spindle fibers. \\
\rowcolor{rowA}  \textbf{T}elophase & Chromosomes arrive at poles and \textbf{decondense}. New nuclear envelopes form around each set. Spindle disassembles. \\
\rowcolor{rowTe} Cytokinesis      & The cytoplasm physically divides. Animal cells form a \textbf{cleavage furrow}; plant cells build a \textbf{cell plate} that becomes a new cell wall. \\
\bottomrule
\end{tabularx}}

\textit{Memory aid: \textbf{IPMAT} = Interphase, Prophase, Metaphase, Anaphase, Telophase, plus cytokinesis.}

\section*{Meiosis I \& II}

\textbf{Meiosis} is a specialized division producing \textbf{gametes} (sperm and eggs). It
goes through \emph{two} divisions and produces \textbf{4 genetically unique haploid cells}
from one diploid parent. The chromosome number is halved (2n $\to$ n).

\textbf{Meiosis I (the reductional division):}
\begin{itemize}[noitemsep]
  \item \textbf{Prophase I:} \textbf{Crossing over} occurs! Homologous chromosomes pair up (forming \textbf{tetrads}) and exchange segments at \textbf{chiasmata}. Major source of genetic variation.
  \item \textbf{Metaphase I:} Homologous pairs (tetrads) line up at the equator. \textbf{Independent assortment} happens here --- each pair orients randomly, giving 2\textsuperscript{23} possible combinations in humans.
  \item \textbf{Anaphase I:} Homologous chromosomes \textbf{separate} (sister chromatids stay together!). This is what halves the chromosome number.
  \item \textbf{Telophase I + Cytokinesis:} Two haploid cells form, each with chromosomes still as sister chromatid pairs.
\end{itemize}

\textbf{Meiosis II (looks just like mitosis):}
\begin{itemize}[noitemsep]
  \item Sister chromatids \textbf{separate} (like mitosis). 
  \item Result: \textbf{4 haploid genetically unique gametes}.
\end{itemize}

\textbf{Three sources of genetic variation in meiosis:}
\begin{enumerate}[noitemsep]
  \item \textbf{Crossing over} in Prophase I (recombines maternal and paternal alleles).
  \item \textbf{Independent assortment} in Metaphase I (which homolog goes to which pole is random).
  \item \textbf{Random fertilization} (any sperm can fertilize any egg).
\end{enumerate}

\section*{Glycolysis (Cytoplasm)}

\textbf{Glycolysis} = ``splitting glucose.'' The first stage of cellular respiration, occurring
in the \textbf{cytoplasm} of \emph{every} cell type --- it does \textbf{not} require oxygen
or mitochondria, making it universal to all life.

\textbf{Net equation:} \quad Glucose (6C) + 2 NAD\textsuperscript{+} + 2 ADP + 2 P\textsubscript{i}
$\to$ 2 Pyruvate (3C) + 2 NADH + \textbf{2 ATP (net)} + 2 H\textsubscript{2}O

\textbf{Three phases:}
\begin{itemize}[noitemsep]
  \item \textbf{Energy investment} -- 2 ATP are spent to phosphorylate and split glucose into 2 G3P molecules.
  \item \textbf{Cleavage} -- the 6-carbon sugar is split into two 3-carbon molecules (G3P).
  \item \textbf{Energy harvest} -- 4 ATP and 2 NADH are produced. Net gain: 2 ATP, 2 NADH, 2 pyruvate.
\end{itemize}

\textit{Glycolysis evolved before oxygen was abundant in Earth's atmosphere. That's why it's
in the cytoplasm of every cell --- it predates mitochondria.}

\section*{Pyruvate Oxidation \& Krebs Cycle}

If oxygen is available, pyruvate enters the \textbf{mitochondrial matrix} and is processed
through two more stages.

\subsection*{Pyruvate Oxidation (Link Reaction)}

Each pyruvate (3C) is converted to \textbf{Acetyl-CoA} (2C):
\begin{itemize}[noitemsep]
  \item Loses one carbon as CO\textsubscript{2}
  \item Reduces 1 NAD\textsuperscript{+} to NADH
  \item Attaches Coenzyme A
\end{itemize}
Per glucose (which yields 2 pyruvate): \textbf{2 CO\textsubscript{2}, 2 NADH, 2 Acetyl-CoA}.

\subsection*{Krebs Cycle (Citric Acid Cycle / TCA Cycle)}

Acetyl-CoA enters the cycle by combining with oxaloacetate (4C) to form \textbf{citrate} (6C).
The cycle then strips off carbons as CO\textsubscript{2}, generates electron carriers
(NADH and FADH\textsubscript{2}), and regenerates oxaloacetate to start again.

\textbf{Per turn of the cycle (1 Acetyl-CoA):} \quad 2 CO\textsubscript{2}, 3 NADH, 1 FADH\textsubscript{2}, 1 ATP (or GTP)

\textbf{Per glucose (2 turns):} \quad \textbf{4 CO\textsubscript{2}, 6 NADH, 2 FADH\textsubscript{2}, 2 ATP}

\textbf{Where it occurs:} the \textbf{mitochondrial matrix}. The cycle has 8 enzyme-catalyzed steps, but for an intro course you only need to know the inputs, outputs, and location.

\textit{The CO\textsubscript{2} you exhale comes from this cycle (and from pyruvate oxidation).
Every breath out is the carbon skeleton of yesterday's food.}

\section*{Electron Transport Chain \& Chemiosmosis}

The \textbf{ETC} is where the \textbf{vast majority of ATP} is produced. Located in the
\textbf{inner mitochondrial membrane} (or in the plasma membrane in bacteria).

\textbf{Step 1 -- Electrons enter the chain.} NADH and FADH\textsubscript{2} (made in
glycolysis, pyruvate oxidation, and Krebs) donate their high-energy electrons to a series
of protein complexes (Complexes I, II, III, IV).

\textbf{Step 2 -- Electrons cascade down the energy gradient.} As electrons pass from one
complex to the next, each transfer releases energy. This energy is used to \textbf{pump
H\textsuperscript{+} ions} from the matrix into the intermembrane space, building a steep
\textbf{electrochemical proton gradient}.

\textbf{Step 3 -- Oxygen accepts the spent electrons.} At Complex IV, electrons (now
low-energy) combine with \textbf{O\textsubscript{2}} and H\textsuperscript{+} to form
\textbf{water}. \emph{This is why we need to breathe oxygen --- it's the final electron
acceptor.} Without O\textsubscript{2}, the entire chain backs up and ATP production stops.

\textbf{Step 4 -- Chemiosmosis: ATP synthesis.} The H\textsuperscript{+} gradient pushes
protons back into the matrix through \textbf{ATP synthase}, a molecular turbine. The flow
of protons spins the rotor, mechanically generating ATP from ADP + P\textsubscript{i}.
This is called \textbf{oxidative phosphorylation}.

\textbf{ATP yield from one glucose (eukaryotes):}

\noindent
{\renewcommand{\arraystretch}{1.4}
\begin{tabularx}{\textwidth}{>{\bfseries\raggedright\arraybackslash}p{4cm}
                              >{\centering\arraybackslash}p{2cm}
                              >{\raggedright\arraybackslash}X}
\toprule
\rowcolor{tableheaderteal}
\textbf{Stage} & \textbf{ATP Yield} & \textbf{Notes} \\
\midrule
\rowcolor{rowTe} Glycolysis            & 2 (net)  & Substrate-level phosphorylation \\
\rowcolor{rowA}  Pyruvate oxidation    & 0        & But generates 2 NADH \\
\rowcolor{rowTe} Krebs cycle           & 2        & Substrate-level phosphorylation \\
\rowcolor{rowA}  ETC + chemiosmosis    & $\sim$26--28 & From all the NADH and FADH\textsubscript{2} \\
\rowcolor{rowTe} \textbf{Total}        & \textbf{$\sim$30--32 ATP} & Per glucose, with oxygen \\
\bottomrule
\end{tabularx}}

\textit{Older textbooks say 36-38 ATP. Modern values are lower because some energy is used
to transport NADH into the mitochondria.}

\section*{Light-Dependent Reactions (Photosynthesis Stage 1)}

\textbf{Where:} the \textbf{thylakoid membranes} of chloroplasts.\\
\textbf{Inputs:} light, H\textsubscript{2}O, NADP\textsuperscript{+}, ADP + P\textsubscript{i}.\\
\textbf{Outputs:} O\textsubscript{2}, NADPH, ATP.

\textbf{Steps:}
\begin{enumerate}[noitemsep]
  \item \textbf{Photosystem II} absorbs light; an electron is excited to a high energy level.
  \item To replace the lost electron, water is split (\textbf{photolysis}): 2 H\textsubscript{2}O $\to$ 4 H\textsuperscript{+} + 4 e\textsuperscript{$-$} + O\textsubscript{2}. \textbf{The O\textsubscript{2} you breathe came from this step.}
  \item Electrons travel down the chloroplast \textbf{ETC}, pumping H\textsuperscript{+} into the thylakoid lumen.
  \item \textbf{Photosystem I} absorbs more light, re-energizes the electrons.
  \item Electrons reduce NADP\textsuperscript{+} $\to$ NADPH.
  \item H\textsuperscript{+} flow through ATP synthase $\to$ ATP. (\textbf{Photophosphorylation}.)
\end{enumerate}

\textit{Notice the parallel: photosynthesis and respiration both use chemiosmosis with ATP synthase. The chloroplast and mitochondrion are evolutionarily related --- both descend from ancient bacteria.}

\section*{Calvin Cycle (Photosynthesis Stage 2)}

\textbf{Where:} the \textbf{stroma} of chloroplasts (the fluid surrounding the thylakoids).\\
\textbf{Inputs:} CO\textsubscript{2}, ATP and NADPH (from light reactions).\\
\textbf{Output:} G3P (a 3-carbon sugar; 2 G3P combine to make glucose).

\textbf{Three stages:}
\begin{enumerate}[noitemsep]
  \item \textbf{Carbon fixation:} The enzyme \textbf{RuBisCO} attaches CO\textsubscript{2} to a 5C molecule (RuBP), making an unstable 6C molecule that immediately splits into two 3C molecules (3-PGA).
  \item \textbf{Reduction:} ATP and NADPH are used to convert 3-PGA into G3P.
  \item \textbf{Regeneration:} Most G3P is recycled, using ATP, to regenerate RuBP so the cycle can continue.
\end{enumerate}

\textbf{It takes 6 turns of the Calvin cycle (fixing 6 CO\textsubscript{2}) to produce 12 G3P, of which 2 leave to make 1 glucose. The other 10 are recycled.}

\textit{Calvin cycle is sometimes called ``light-independent'' --- but that's misleading. It needs ATP and NADPH from the light reactions, so it stops at night.}

\section*{C3, C4 \& CAM Photosynthesis}

Plants in different climates use different strategies to fix carbon. Hot, dry climates pose
a problem: closing stomata to save water also blocks CO\textsubscript{2} entry. Different
plants solve this differently.

\noindent
{\renewcommand{\arraystretch}{1.6}
\begin{tabularx}{\textwidth}{>{\bfseries\raggedright\arraybackslash}p{1.5cm}
                              >{\raggedright\arraybackslash}X
                              >{\raggedright\arraybackslash}p{3cm}}
\toprule
\rowcolor{tableheaderteal}
\textbf{Type} & \textbf{Strategy} & \textbf{Examples} \\
\midrule
\rowcolor{rowTe} \textbf{C3} & Standard photosynthesis. CO\textsubscript{2} is fixed directly by RuBisCO into a 3-carbon compound (3-PGA). Suffers from \textbf{photorespiration} when hot/dry --- RuBisCO accidentally grabs O\textsubscript{2} instead of CO\textsubscript{2}, wasting energy. & Most plants (rice, wheat, soybeans, trees) \\
\rowcolor{rowA}  \textbf{C4} & Pre-fixes CO\textsubscript{2} into a 4-carbon compound in mesophyll cells, then shuttles it to specialized bundle sheath cells where the Calvin cycle runs at high CO\textsubscript{2}. Avoids photorespiration. Better in hot, sunny conditions. & Corn, sugarcane, sorghum, crabgrass \\
\rowcolor{rowTe} \textbf{CAM} & Crassulacean Acid Metabolism. \textbf{Opens stomata only at night} to take in CO\textsubscript{2} (stored as malate). During the day, stomata close to conserve water; the stored CO\textsubscript{2} feeds the Calvin cycle. & Cacti, succulents, pineapple, agave \\
\bottomrule
\end{tabularx}}

\textbf{Trade-offs:} C3 is most efficient in cool, wet climates. C4 is more efficient in hot,
sunny climates (~3\% of plant species, but $\sim$25\% of land photosynthesis). CAM is most
efficient in arid deserts (sacrifices growth speed for water conservation).

\section*{ATP Production Across Organisms}

Different organisms generate ATP using different terminal electron acceptors and pathways.

\noindent
{\small\renewcommand{\arraystretch}{1.6}
\begin{tabularx}{\textwidth}{>{\bfseries\raggedright\arraybackslash}p{2.8cm}
                              >{\raggedright\arraybackslash}X
                              >{\raggedright\arraybackslash}p{3.5cm}}
\toprule
\rowcolor{tableheaderteal}
\textbf{Organism / Type} & \textbf{ATP production strategy} & \textbf{Examples} \\
\midrule
\rowcolor{rowTe} \textbf{Aerobic eukaryotes (humans, animals, plants)} & Glycolysis $\to$ Krebs $\to$ ETC with O\textsubscript{2} as final electron acceptor. $\sim$30--32 ATP per glucose. & Humans, mice, oak trees \\
\rowcolor{rowA}  \textbf{Yeast (fermentation)} & Glycolysis (2 ATP) $\to$ ethanol fermentation (regenerates NAD\textsuperscript{+} but no extra ATP). 2 ATP per glucose. & Brewer's/baker's yeast (\textit{Saccharomyces}) \\
\rowcolor{rowTe} \textbf{Muscle cells (anaerobic)} & Glycolysis (2 ATP) $\to$ lactic acid fermentation when O\textsubscript{2} is limited (sprinting). 2 ATP per glucose. ``Burn'' is from lactic acid. & Human muscle during intense exercise \\
\rowcolor{rowA}  \textbf{Aerobic bacteria} & Same as eukaryotes but the ETC is in the plasma membrane (no mitochondria). & \textit{E.\ coli} with O\textsubscript{2}, \textit{Pseudomonas} \\
\rowcolor{rowTe} \textbf{Sulfate-reducing bacteria} & Anaerobic respiration using SO\textsubscript{4}\textsuperscript{2$-$} as final electron acceptor instead of O\textsubscript{2}. Produces H\textsubscript{2}S (rotten egg smell). & \textit{Desulfovibrio} (deep-sea sediments) \\
\rowcolor{rowA}  \textbf{Methanogens (Archaea)} & Anaerobic respiration; reduce CO\textsubscript{2} with H\textsubscript{2} to produce \textbf{methane (CH\textsubscript{4})}. & Cow rumen, swamps, deep-sea vents \\
\rowcolor{rowTe} \textbf{Nitrate-reducing bacteria} & Use NO\textsubscript{3}\textsuperscript{$-$} as final electron acceptor, producing N\textsubscript{2}, N\textsubscript{2}O, or NH\textsubscript{3}. Important in soil ecology. & Many soil bacteria \\
\rowcolor{rowA}  \textbf{Photosynthetic organisms} & Use light directly to make ATP \emph{and} fix carbon. Plants, algae, cyanobacteria. & Plants, algae, cyanobacteria \\
\rowcolor{rowTe} \textbf{Chemoautotrophs} & Get energy from oxidizing inorganic chemicals (H\textsubscript{2}S, NH\textsubscript{3}, Fe\textsuperscript{2+}) rather than light or food. & Hydrothermal vent bacteria \\
\bottomrule
\end{tabularx}}

\section*{Fermentation Pathways}

\textbf{Fermentation} = an anaerobic way of regenerating NAD\textsuperscript{+} so glycolysis
can continue. It does \textbf{not} produce additional ATP --- the only ATP is from glycolysis
itself (2 net per glucose). The point is to \emph{recycle} NADH back to NAD\textsuperscript{+}.

\noindent
{\renewcommand{\arraystretch}{1.6}
\begin{tabularx}{\textwidth}{>{\bfseries\raggedright\arraybackslash}p{3cm}
                              >{\raggedright\arraybackslash}X}
\toprule
\rowcolor{tableheaderteal}
\textbf{Type} & \textbf{Process \& Where} \\
\midrule
\rowcolor{rowTe} \textbf{Lactic Acid Fermentation} & Pyruvate $\to$ Lactate (lactic acid). Used by \textbf{human muscle cells} during intense exercise (when O\textsubscript{2} can't keep up), \textbf{lactobacillus} bacteria (yogurt, cheese, sourdough, sauerkraut, kimchi). \\
\rowcolor{rowA}  \textbf{Alcoholic Fermentation} & Pyruvate $\to$ Acetaldehyde $\to$ Ethanol + CO\textsubscript{2}. Used by \textbf{yeast} (beer, wine, bread). The CO\textsubscript{2} makes bread rise; the ethanol stays in alcoholic drinks but evaporates from bread. \\
\bottomrule
\end{tabularx}}

\textbf{Why is lactic acid produced when you sprint?} Your muscles can't get enough oxygen
fast enough during high-intensity activity. The ETC backs up. To keep glycolysis going (and
keep producing some ATP), pyruvate accepts electrons from NADH, becoming lactate. The lactate
is later transported to the liver, where (when oxygen is available again) it's converted
back to pyruvate. The ``burn'' you feel is acidification from H\textsuperscript{+}, not the
lactate itself.

\section*{Nitrogen \& Carbon Cycles}

\subsection*{Nitrogen Cycle}

Nitrogen is essential for proteins, DNA, and amino acids. The atmosphere is 78\%
N\textsubscript{2}, but plants and animals can't use it directly.

\begin{enumerate}[noitemsep]
  \item \textbf{Nitrogen fixation:} Bacteria convert N\textsubscript{2} $\to$ NH\textsubscript{3} (ammonia). Done by \textit{Rhizobium} (in legume nodules), \textit{Azotobacter}, and cyanobacteria.
  \item \textbf{Nitrification:} Soil bacteria convert NH\textsubscript{3} $\to$ NO\textsubscript{2}\textsuperscript{$-$} $\to$ NO\textsubscript{3}\textsuperscript{$-$} (nitrate). Plants absorb nitrate.
  \item \textbf{Assimilation:} Plants use nitrogen to build proteins and DNA. Animals get nitrogen by eating plants.
  \item \textbf{Ammonification:} Decomposers convert dead matter back to NH\textsubscript{3}.
  \item \textbf{Denitrification:} Anaerobic bacteria convert NO\textsubscript{3}\textsuperscript{$-$} back to N\textsubscript{2}, returning it to the atmosphere.
\end{enumerate}

\subsection*{Carbon Cycle}

Carbon moves between the atmosphere, oceans, living things, and rocks/fossil fuels.

\begin{enumerate}[noitemsep]
  \item \textbf{Photosynthesis} pulls CO\textsubscript{2} out of the atmosphere into plant tissue.
  \item \textbf{Cellular respiration} (by all organisms) returns CO\textsubscript{2} to the atmosphere.
  \item \textbf{Decomposition} releases CO\textsubscript{2} as decomposers break down dead matter.
  \item \textbf{Combustion} (burning fossil fuels, wildfires, deforestation) releases stored carbon as CO\textsubscript{2}.
  \item \textbf{Ocean absorption} dissolves CO\textsubscript{2} into seawater (forming carbonic acid; this drives ocean acidification).
  \item \textbf{Sedimentation} traps carbon long-term in shells, limestone, and fossil fuels.
\end{enumerate}

\textit{Burning fossil fuels releases carbon stored over millions of years in just centuries
--- this is the geological reason for climate change.}


\newpage

%% =============================================================================
%%  PART IIo -- Taxonomy & The Tree of Life
%% =============================================================================
\purplesections
\addcontentsline{toc}{section}{\textcolor{sectionpurple}{PART IIo \textemdash\ Taxonomy \& The Tree of Life}}
\addcontentsline{toc}{subsection}{The Hierarchy: Domain to Species}
\addcontentsline{toc}{subsection}{Binomial Nomenclature}
\addcontentsline{toc}{subsection}{The Three Domains}
\addcontentsline{toc}{subsection}{Kingdom Animalia}
\addcontentsline{toc}{subsection}{Major Animal Phyla}
\addcontentsline{toc}{subsection}{Vertebrate Classes}
\addcontentsline{toc}{subsection}{Kingdom Plantae}
\addcontentsline{toc}{subsection}{Kingdom Fungi}
\addcontentsline{toc}{subsection}{Kingdom Protista}
\addcontentsline{toc}{subsection}{Domains Bacteria \& Archaea}
\addcontentsline{toc}{subsection}{Cladistics \& Phylogenetic Trees}
\partbanner{part3bg}{sectionpurple}{Part IIo \quad Taxonomy \& The Tree of Life}

\vspace{6pt}
\begin{center}\small\itshape
  How biologists organize the millions of species on Earth into a unified system,
  starting from broad categories and narrowing all the way down to individual species.
\end{center}

\section*{The Hierarchy: Domain to Species}

\textbf{Taxonomy} is the science of naming, describing, and classifying organisms. The
modern system uses \textbf{8 hierarchical ranks}, each more specific than the last. From
broadest (most inclusive) to narrowest:

\noindent
{\renewcommand{\arraystretch}{1.6}
\begin{tabularx}{\textwidth}{>{\bfseries\centering\arraybackslash}p{1cm}
                              >{\bfseries\raggedright\arraybackslash}p{2.5cm}
                              >{\raggedright\arraybackslash}X
                              >{\itshape\raggedright\arraybackslash}p{4cm}}
\toprule
\rowcolor{tableheaderpurple}
\textbf{\#} & \textbf{Rank} & \textbf{What it groups together} & \textbf{Human example} \\
\midrule
\rowcolor{rowPu} 1 & Domain    & Largest, broadest grouping. Only 3 exist (Bacteria, Archaea, Eukarya). & Eukarya \\
\rowcolor{rowA}  2 & Kingdom   & The largest grouping below domain. Contains all related phyla. & Animalia \\
\rowcolor{rowPu} 3 & Phylum (pl.\ phyla) & Major body plan / structural design. & Chordata (have a notochord) \\
\rowcolor{rowA}  4 & Class     & Major subgroup with shared traits within a phylum. & Mammalia \\
\rowcolor{rowPu} 5 & Order     & Closely related families within a class. & Primates \\
\rowcolor{rowA}  6 & Family    & Closely related genera; ends in \emph{-idae} for animals, \emph{-aceae} for plants. & Hominidae (great apes \& humans) \\
\rowcolor{rowPu} 7 & Genus (pl.\ genera) & A small group of very closely related species. Italicized; capitalized first letter. & \textit{Homo} \\
\rowcolor{rowA}  8 & Species   & The most specific level. Italicized, lowercase. & \textit{sapiens} \\
\bottomrule
\end{tabularx}}

\textbf{Memory aids} (pick whichever is easiest):
\begin{itemize}[noitemsep]
  \item \textbf{D}ear \textbf{K}ing \textbf{P}hilip \textbf{C}ame \textbf{O}ver \textbf{F}or \textbf{G}ood \textbf{S}oup
  \item \textbf{D}id \textbf{K}ing \textbf{P}hilip \textbf{C}ome \textbf{O}ver \textbf{F}or \textbf{G}reat \textbf{S}pecies?
  \item \textbf{D}umb \textbf{K}ids \textbf{P}laying \textbf{C}atch \textbf{O}n \textbf{F}reeway \textbf{G}et \textbf{S}quashed
\end{itemize}

\textit{Note: There are sub-ranks too (subphylum, suborder, etc.), but these 8 are the major ranks every biology student must know.}

\section*{Binomial Nomenclature}

Devised by Carl Linnaeus in the 1750s. Every species is given a unique \textbf{two-part Latin name}: \textbf{Genus species}.

\textbf{Rules:}
\begin{itemize}[noitemsep]
  \item \textbf{Genus} is capitalized; \textbf{species} is lowercase.
  \item Both parts are \emph{italicized} (or underlined when handwritten).
  \item After first use, genus can be abbreviated: \textit{Homo sapiens} $\to$ \textit{H.\ sapiens}.
  \item Latin/Latinized to be universal --- the same name is used worldwide regardless of native language.
\end{itemize}

\textbf{Examples:}
\begin{itemize}[noitemsep]
  \item \textit{Homo sapiens} (humans)
  \item \textit{Canis lupus familiaris} (domestic dog --- a subspecies)
  \item \textit{Felis catus} (house cat)
  \item \textit{Escherichia coli} (a common gut bacterium)
  \item \textit{Tyrannosaurus rex} (extinct dinosaur)
  \item \textit{Quercus alba} (white oak tree)
\end{itemize}

\textbf{Why this matters:} ``Crow'' might mean different birds in different countries. \textit{Corvus brachyrhynchos} means exactly one species, everywhere on Earth.

\section*{The Three Domains}

The highest level of classification. All known life falls into one of three domains.

\noindent
{\renewcommand{\arraystretch}{1.6}
\begin{tabularx}{\textwidth}{>{\bfseries\raggedright\arraybackslash}p{2.5cm}
                              >{\raggedright\arraybackslash}X}
\toprule
\rowcolor{tableheaderpurple}
\textbf{Domain} & \textbf{Description} \\
\midrule
\rowcolor{rowPu} \textbf{Bacteria}  & Prokaryotic (no nucleus). Single-celled. Peptidoglycan cell walls. The most diverse, abundant organisms on Earth. Includes \textit{E.\ coli}, \textit{Streptococcus}, cyanobacteria. \\
\rowcolor{rowA}  \textbf{Archaea}   & Prokaryotic. Single-celled. Look like bacteria but biochemically distinct (different cell wall composition; different membrane lipids; some genes more like eukaryotes). Often live in \textbf{extreme environments}: hot springs, deep-sea vents, salt lakes, animal guts. \\
\rowcolor{rowPu} \textbf{Eukarya}   & \textbf{Eukaryotic}: cells have a nucleus and membrane-bound organelles. All multicellular life lives here, plus single-celled protists. Contains 4 kingdoms: Animalia, Plantae, Fungi, Protista. \\
\bottomrule
\end{tabularx}}

\textbf{Why three domains, not two?} Until 1977, scientists used a 5-kingdom system (Monera,
Protista, Plantae, Fungi, Animalia). Then Carl Woese compared rRNA gene sequences and
discovered Archaea were \emph{not} just weird bacteria --- they're as different from Bacteria
as we are. The 3-domain system is now standard.

\section*{Kingdom Animalia}

\textbf{Multicellular, eukaryotic, heterotrophic} (consume other organisms for energy),
\textbf{no cell walls}, capable of movement (at least at some life stage), develop from a
\textbf{blastula} embryo. About 1.5 million described species; estimates of total species
range from 7 million to over 10 million.

\textbf{Major shared traits:}
\begin{itemize}[noitemsep]
  \item \textbf{Tissues} differentiated from embryonic germ layers (ectoderm, mesoderm, endoderm)
  \item Most have \textbf{nervous systems} and \textbf{muscle tissue}
  \item Most reproduce \textbf{sexually}; some can also reproduce asexually
  \item Generally \textbf{diploid} as adults; only gametes are haploid
\end{itemize}

\section*{Major Animal Phyla}

There are about 35 animal phyla; here are the 9 you should know.

\noindent
{\small\renewcommand{\arraystretch}{1.6}
\begin{tabularx}{\textwidth}{>{\bfseries\raggedright\arraybackslash}p{2.8cm}
                              >{\raggedright\arraybackslash}X
                              >{\raggedright\arraybackslash}p{3.2cm}}
\toprule
\rowcolor{tableheaderpurple}
\textbf{Phylum} & \textbf{Key features} & \textbf{Examples} \\
\midrule
\rowcolor{rowPu} \textbf{Porifera}    & Sponges. \textbf{No true tissues}; filter feeders. Asymmetrical. Sessile (don't move). Considered the most ancient animal lineage. & Sea sponges \\
\rowcolor{rowA}  \textbf{Cnidaria}    & \textbf{Radial symmetry}, stinging cells (\emph{cnidocytes}). Two body forms: polyp (attached) and medusa (swimming). True tissues but no organs. & Jellyfish, coral, sea anemones, hydra \\
\rowcolor{rowPu} \textbf{Platyhelminthes} & Flatworms. \textbf{Bilateral symmetry}; first phylum with a recognizable head. No body cavity. & Tapeworms, planarians, flukes \\
\rowcolor{rowA}  \textbf{Nematoda}    & \textbf{Roundworms}. Cylindrical, unsegmented. Have a complete digestive tract (mouth + anus). \textit{C.\ elegans} is a key model organism. & Hookworm, pinworm, \textit{C.\ elegans} \\
\rowcolor{rowPu} \textbf{Mollusca}    & Soft body, often with a shell. Most have a muscular ``foot,'' a mantle, and a radula (rasping tongue). Second-largest phylum. & Snails, clams, octopus, squid \\
\rowcolor{rowA}  \textbf{Annelida}    & \textbf{Segmented worms}. Each segment has its own coelom and partial repetition of organs. & Earthworms, leeches, polychaetes \\
\rowcolor{rowPu} \textbf{Arthropoda}  & The \textbf{largest animal phylum} --- $\sim$80\% of known species! Exoskeleton (chitin), segmented body, jointed appendages. & Insects, spiders, crabs, lobsters, scorpions \\
\rowcolor{rowA}  \textbf{Echinodermata} & ``Spiny-skinned.'' Adults have \textbf{five-part radial symmetry} (but larvae are bilateral). Water vascular system with tube feet. & Starfish, sea urchins, sand dollars, sea cucumbers \\
\rowcolor{rowPu} \textbf{Chordata}    & Have a \textbf{notochord} (flexible rod) at some life stage, dorsal hollow nerve cord, pharyngeal slits, and a post-anal tail. Includes all vertebrates. & Humans, fish, birds, reptiles, mammals, plus tunicates and lancelets \\
\bottomrule
\end{tabularx}}

\begin{tcolorbox}[colback=part3bg, colframe=sectionpurple, arc=3pt, boxrule=1pt]
\textbf{Symmetry types in animals:}
\begin{itemize}[noitemsep, leftmargin=*]
  \item \textbf{Asymmetrical} -- no symmetry at all (sponges).
  \item \textbf{Radial} -- multiple planes of symmetry through the center; no front/back; common in sessile or floating animals (cnidarians, adult echinoderms).
  \item \textbf{Bilateral} -- one plane of symmetry; left/right mirror images; has a distinct head end (\textbf{cephalization}). Allows fast directional movement. Almost all complex animals are bilateral.
\end{itemize}
\end{tcolorbox}

\section*{Vertebrate Classes (Phylum Chordata)}

Within phylum Chordata, the subphylum \textbf{Vertebrata} contains all animals with backbones.
Six major classes:

\noindent
{\small\renewcommand{\arraystretch}{1.6}
\begin{tabularx}{\textwidth}{>{\bfseries\raggedright\arraybackslash}p{2.8cm}
                              >{\raggedright\arraybackslash}X
                              >{\raggedright\arraybackslash}p{3cm}}
\toprule
\rowcolor{tableheaderpurple}
\textbf{Class} & \textbf{Distinguishing features} & \textbf{Examples} \\
\midrule
\rowcolor{rowPu} \textbf{Agnatha} (or Cyclostomata) & \textbf{Jawless fish}. No paired fins. Cartilage skeleton. The most ancient vertebrate lineage. & Lampreys, hagfish \\
\rowcolor{rowA}  \textbf{Chondrichthyes} & \textbf{Cartilaginous fish.} Have jaws and paired fins, but skeletons are cartilage, not bone. & Sharks, rays, skates \\
\rowcolor{rowPu} \textbf{Osteichthyes} & \textbf{Bony fish}. Bone skeleton, scales, swim bladder, gills. Most numerous vertebrate class. & Salmon, tuna, goldfish, bass \\
\rowcolor{rowA}  \textbf{Amphibia}  & ``Double life.'' \textbf{Aquatic larvae, terrestrial adults.} Moist permeable skin. Eggs lack shells $\to$ must be laid in water. \textbf{Ectothermic.} & Frogs, toads, salamanders, newts \\
\rowcolor{rowPu} \textbf{Reptilia}  & \textbf{Scaly, dry skin.} \textbf{Amniotic eggs} with leathery shells (a major innovation --- allowed full terrestrial life). \textbf{Ectothermic.} Includes birds in modern phylogeny. & Snakes, lizards, turtles, crocodiles. (Birds branch from theropod dinosaurs.) \\
\rowcolor{rowA}  \textbf{Aves} (Birds) & Feathers, wings, hollow bones, beaks (no teeth), \textbf{endothermic} (warm-blooded). Lay hard-shelled amniotic eggs. Modern phylogeny groups them within reptiles, but they're often listed as their own class. & Eagles, sparrows, penguins, ostriches \\
\rowcolor{rowPu} \textbf{Mammalia} (Mammals) & \textbf{Hair/fur}, \textbf{mammary glands} (milk for young), 3 middle ear bones, \textbf{endothermic}. Most are viviparous (live birth). & Humans, dogs, whales, bats, kangaroos \\
\bottomrule
\end{tabularx}}

\textbf{Major mammal subgroups:}
\begin{itemize}[noitemsep]
  \item \textbf{Monotremes:} egg-laying mammals (platypus, echidna). Only 5 species.
  \item \textbf{Marsupials:} young born tiny, develop in a pouch (kangaroo, koala, opossum). Concentrated in Australia.
  \item \textbf{Placentals:} the vast majority --- young develop fully inside the uterus connected via placenta (humans, dogs, whales).
\end{itemize}

\textbf{Major mammal orders:}
\begin{itemize}[noitemsep]
  \item \textbf{Primates}: humans, monkeys, apes, lemurs.
  \item \textbf{Carnivora}: dogs, cats, bears, seals.
  \item \textbf{Cetacea}: whales, dolphins, porpoises (descended from land mammals).
  \item \textbf{Rodentia}: rodents --- the largest mammal order ($\sim$40\% of mammal species).
  \item \textbf{Chiroptera}: bats --- the only mammals capable of true flight.
  \item \textbf{Artiodactyla}: even-toed ungulates (cows, pigs, deer, giraffes; cetaceans are nested within!).
\end{itemize}

\section*{Kingdom Plantae}

\textbf{Multicellular, eukaryotic, photosynthetic} (autotrophs), with \textbf{cellulose cell walls}
and \textbf{chloroplasts}. About 350{,}000 species described.

\textbf{Four major plant groups} (in evolutionary order):

\noindent
{\renewcommand{\arraystretch}{1.6}
\begin{tabularx}{\textwidth}{>{\bfseries\raggedright\arraybackslash}p{3.5cm}
                              >{\raggedright\arraybackslash}X
                              >{\raggedright\arraybackslash}p{2.8cm}}
\toprule
\rowcolor{tableheaderpurple}
\textbf{Group} & \textbf{Key innovations \& features} & \textbf{Examples} \\
\midrule
\rowcolor{rowPu} \textbf{Bryophytes} (non-vascular) & Earliest land plants. \textbf{No vascular tissue} --- water and nutrients move by diffusion, so they stay small. Need water for reproduction (sperm swim). & Mosses, liverworts, hornworts \\
\rowcolor{rowA}  \textbf{Pteridophytes} (seedless vascular) & Have \textbf{xylem and phloem} (vascular tissue) --- can grow tall. But still reproduce via spores; still need water for fertilization. & Ferns, horsetails, club mosses \\
\rowcolor{rowPu} \textbf{Gymnosperms} (``naked seeds'') & Vascular + \textbf{seeds} (no longer dependent on water for reproduction) + \textbf{pollen} (sperm carried by wind). Seeds are not enclosed in fruit. & Conifers (pines, firs, spruces), cycads, ginkgo \\
\rowcolor{rowA}  \textbf{Angiosperms} (flowering plants) & Vascular + seeds + \textbf{flowers} (attract pollinators) + \textbf{fruit} (protects and disperses seeds). The dominant land plants today. & Roses, oak, grass, corn, apple, sunflower \\
\bottomrule
\end{tabularx}}

\textbf{Two main groups of angiosperms:}
\begin{itemize}[noitemsep]
  \item \textbf{Monocots}: 1 cotyledon (seed leaf), parallel leaf veins, flower parts in 3s, fibrous roots. Examples: grasses, corn, lilies, palms, orchids.
  \item \textbf{Dicots}: 2 cotyledons, branching leaf veins, flower parts in 4s or 5s, taproot. Examples: roses, oaks, beans, sunflowers.
\end{itemize}

\section*{Kingdom Fungi}

\textbf{Eukaryotic, mostly multicellular, heterotrophic} (decomposers --- absorb nutrients
from dead organic matter). Cell walls of \textbf{chitin} (same as insect exoskeletons, NOT
cellulose). About 150{,}000 species described, but estimates suggest there may be 5 million.

\textbf{Key features:}
\begin{itemize}[noitemsep]
  \item Bodies made of \textbf{hyphae} (long thread-like cells) that form a network called a \textbf{mycelium}.
  \item Reproduce primarily by \textbf{spores}; both sexually and asexually.
  \item \textbf{Closer relatives to animals than to plants!} (Both fungi and animals are descended from a common ancestor and share traits like chitin and lacking chloroplasts.)
\end{itemize}

\textbf{Major phyla:}
\begin{itemize}[noitemsep]
  \item \textbf{Zygomycota}: bread molds (\textit{Rhizopus}).
  \item \textbf{Ascomycota} (sac fungi): yeasts, morels, truffles, \textit{Penicillium} (penicillin source).
  \item \textbf{Basidiomycota} (club fungi): mushrooms, puffballs, shelf fungi, rusts, smuts.
  \item \textbf{Glomeromycota}: form mutualistic mycorrhizae (symbiotic associations) with plant roots --- $\sim$80\% of plants depend on these.
\end{itemize}

\textbf{Roles of fungi:}
\begin{itemize}[noitemsep]
  \item \textbf{Decomposers}: recycle dead organic matter --- absolutely essential for the carbon cycle. Without fungi, dead trees would never break down.
  \item \textbf{Symbionts}: lichens (fungus + alga), mycorrhizae (fungus + plant roots).
  \item \textbf{Pathogens}: athlete's foot, ringworm, plant diseases (rusts, smuts), Dutch elm disease.
  \item \textbf{Food and medicine}: edible mushrooms, yeasts (bread, beer), penicillin.
\end{itemize}

\section*{Kingdom Protista}

\textbf{The ``catch-all'' kingdom.} Mostly single-celled eukaryotes that don't fit into
animals, plants, or fungi. Polyphyletic --- protists do not all share a common ancestor that
isn't shared with another kingdom, so modern phylogeny is breaking this kingdom up.

\textbf{Three informal groupings based on lifestyle:}
\begin{itemize}[noitemsep]
  \item \textbf{Animal-like} (\emph{protozoa}): heterotrophic, often motile. Examples: \textit{Amoeba}, \textit{Paramecium}, \textit{Plasmodium} (causes malaria), \textit{Trypanosoma} (sleeping sickness).
  \item \textbf{Plant-like} (\emph{algae}): photosynthetic. Includes \textbf{diatoms} (silica shells; produce $\sim$25\% of Earth's O\textsubscript{2}), \textbf{kelp} (giant brown algae), \textbf{green algae}, \textbf{red algae}, \textbf{euglena}.
  \item \textbf{Fungus-like} (\emph{slime molds, water molds}): absorb nutrients like fungi but move during their life cycle.
\end{itemize}

\section*{Domains Bacteria \& Archaea}

\textbf{Bacteria:} The most diverse and abundant organisms on Earth. Found everywhere ---
soil, water, deep crust, your skin, your gut. (See the Microbiology section for full coverage
of structure, Gram staining, oxygen tolerance, and pathology.)

\textbf{Archaea:} Look like bacteria but biochemically distinct. Many live in extreme environments:

\begin{itemize}[noitemsep]
  \item \textbf{Thermophiles / Hyperthermophiles}: hot springs, hydrothermal vents (some thrive at 100\textdegree C+).
  \item \textbf{Halophiles}: extremely salty environments (Dead Sea, salt evaporation ponds).
  \item \textbf{Methanogens}: anaerobic; produce methane (cow guts, swamps, wetlands, deep ocean sediments).
  \item \textbf{Acidophiles}: extreme acid (pH 0--3, like volcanic vents).
  \item \textbf{Piezophiles (barophiles)}: deep-sea high-pressure environments.
\end{itemize}

\textit{Archaea aren't \emph{only} extremophiles --- they're also abundant in ordinary soil
and ocean. But they're the only known organisms that survive in some of Earth's harshest places.}

\section*{Cladistics \& Phylogenetic Trees}

Modern taxonomy is increasingly based on \textbf{evolutionary relationships}, not just
appearance. \textbf{Cladistics} is the method of grouping organisms by shared
\textbf{evolutionary ancestry}, using DNA, RNA, and protein sequences plus morphology.

\textbf{Key terms:}
\begin{itemize}[noitemsep]
  \item \textbf{Clade}: A group consisting of an ancestor and \emph{all} of its descendants. (e.g., birds + dinosaurs are one clade.)
  \item \textbf{Monophyletic group}: A complete clade. Valid taxonomic group.
  \item \textbf{Paraphyletic group}: An ancestor and \emph{some} of its descendants (excludes some descendants). Often considered invalid in modern taxonomy. Example: ``reptiles'' excluding birds is paraphyletic, because birds evolved from within reptiles.
  \item \textbf{Polyphyletic group}: A group whose members come from \emph{different} ancestors --- usually a mistake based on superficial similarity. Example: ``warm-blooded animals'' (mammals + birds evolved endothermy independently).
  \item \textbf{Homologous traits}: Traits inherited from a common ancestor (human arm, whale flipper, bat wing). Used to build trees.
  \item \textbf{Analogous traits}: Similar functions evolved independently (bird wing vs.\ insect wing). NOT used to build trees --- can be misleading.
\end{itemize}

\textbf{How phylogenetic trees are read:}
\begin{itemize}[noitemsep]
  \item \textbf{Nodes} (branching points) represent the most recent common ancestor of the lineages above.
  \item \textbf{Branch length} sometimes represents time or amount of genetic change.
  \item \textbf{Sister taxa}: two species or groups that share an immediate common ancestor.
  \item Position on the page (top vs.\ bottom) does \textbf{not} indicate ``more evolved'' --- all current species are equally ``modern.''
\end{itemize}

\textbf{Why classification keeps changing:} as DNA sequencing got cheap, biologists found
many traditional groupings (often based on visible features) didn't reflect actual ancestry.
Birds were moved into the dinosaur clade. Whales turned out to be artiodactyls (closely
related to hippos). Fungi turned out to be closer to animals than to plants. Modern taxonomy
is a moving target --- what's in your textbook may be revised within a decade.


\newpage

%% ===========================================================================
%%  PART III – Cheat Sheet
%% ===========================================================================
\purplesections
\addcontentsline{toc}{section}{\textcolor{sectionpurple}{PART III \textemdash\ Quick Reference Cheat Sheet}}
\addcontentsline{toc}{subsection}{Scientific Method}
\addcontentsline{toc}{subsection}{Macromolecules}
\addcontentsline{toc}{subsection}{Prokaryotes vs. Eukaryotes}
\addcontentsline{toc}{subsection}{Must-Know Organelles}
\addcontentsline{toc}{subsection}{SET, Cell Theory \& Key Facts}
\addcontentsline{toc}{subsection}{CHNOPS \& Molecular Biology}
\addcontentsline{toc}{subsection}{Membranes, Transport \& Energy}
\partbanner{part3bg}{sectionpurple}{Part III \quad Quick Reference Cheat Sheet}

\vspace{4pt}
\begin{center}\small\itshape One page. Everything essential. Great for last-minute review.\end{center}
\vspace{4pt}

% ── Scientific Method ────────────────────────────────────────────────────────
\begin{tcolorbox}[colback=part2bg, colframe=sectiongreen,
  title={\bfseries Scientific Method}, arc=3pt, boxrule=1pt]
\begin{itemize}[noitemsep, leftmargin=*]
  \item \textbf{Flow:} Observation $\to$ Hypothesis $\to$ Prediction (``if\ldots then\ldots'') $\to$ Experiment $\to$ Conclusion
  \item \textbf{Theory} = rigorously tested, high predictive value, can still be disproven
  \item \textbf{Law} = describes a pattern (often mathematical); does \emph{not} explain why
  \item \textbf{CRAAP}: Currency, Relevance, Authority, Accuracy, Purpose
\end{itemize}
\end{tcolorbox}

\vspace{4pt}

% ── Macromolecules ───────────────────────────────────────────────────────────
\begin{tcolorbox}[colback=part2bg, colframe=sectiongreen,
  title={\bfseries Macromolecules}, arc=3pt, boxrule=1pt]
\footnotesize\renewcommand{\arraystretch}{1.3}
\begin{tabularx}{\linewidth}{>{\bfseries\raggedright\arraybackslash}X >{\raggedright\arraybackslash}X >{\raggedright\arraybackslash}X >{\raggedright\arraybackslash}X}
\toprule
Molecule & Monomer & Solubility & Key Role \\
\midrule
Carbohydrates & Monosaccharides & Hydrophilic & Quick energy; ID tags \\
Lipids        & Fatty acids / glycerol & Hydrophobic & Long-term energy; membranes; hormones \\
Proteins      & Amino acids (20) & Hydrophilic & Enzymes, antibodies, motors, receptors \\
Nucleic Acids & Nucleotides & Hydrophilic & DNA (storage); RNA (transcription/translation) \\
\bottomrule
\end{tabularx}
\end{tcolorbox}

\vspace{4pt}

% ── Prokaryotes vs Eukaryotes ─────────────────────────────────────────────────
\begin{tcolorbox}[colback=partbg, colframe=sectionblue,
  title={\bfseries Prokaryotes vs.\ Eukaryotes}, arc=3pt, boxrule=1pt]
\begin{multicols}{2}
\textbf{Prokaryotes only:}
\begin{itemize}[noitemsep, leftmargin=*]
  \item 1--10 \textmu m; always unicellular
  \item No membrane-bound organelles
  \item Binary fission; nucleoid (no true nucleus)
  \item Circular DNA; oldest cell type
\end{itemize}
\columnbreak
\textbf{Eukaryotes only:}
\begin{itemize}[noitemsep, leftmargin=*]
  \item 10--100 \textmu m; can be multicellular
  \item Membrane-bound organelles
  \item Mitosis / meiosis; true nucleus
  \item Linear DNA; plants, animals, fungi
\end{itemize}
\end{multicols}
\textbf{Both share:} plasma membrane, ribosomes, cytoplasm, DNA (can have cell walls)
\end{tcolorbox}

\vspace{4pt}

% ── Organelles ───────────────────────────────────────────────────────────────
\begin{tcolorbox}[colback=partbg, colframe=sectionblue,
  title={\bfseries Must-Know Organelles}, arc=3pt, boxrule=1pt]
\footnotesize
\begin{multicols}{2}
\begin{itemize}[noitemsep, leftmargin=*]
  \item \textbf{Nucleus} -- DNA \& gene control
  \item \textbf{Nuclear Pore} -- ribosomes exit here
  \item \textbf{Nuclear Lamina} -- organizes DNA; supports nuclear envelope
  \item \textbf{Rough ER} -- protein synthesis (secreted/membrane proteins)
  \item \textbf{Smooth ER} -- lipids \& detox
  \item \textbf{Golgi} -- sort \& ship proteins/lipids
  \item \textbf{Lysosome} -- waste disposal
  \item \textbf{Peroxisome} -- fatty acid breakdown $\to$ H$_2$O$_2$ $\to$ H$_2$O + CO$_2$
  \item \textbf{Mitochondria} -- ATP via respiration
  \item \textbf{Chloroplast} -- ATP via photosynthesis (plants)
  \item \textbf{Nucleolus} -- makes ribosomes
  \item \textbf{Ribosome} -- assembles proteins
\end{itemize}
\end{multicols}
\end{tcolorbox}

\vspace{4pt}

% ── SET + Cell Theory ─────────────────────────────────────────────────────────
\begin{tcolorbox}[colback=part3bg, colframe=sectionpurple,
  title={\bfseries SET, Cell Theory, \& Key Facts}, arc=3pt, boxrule=1pt]
\begin{itemize}[noitemsep, leftmargin=*]
  \item \textbf{SET}: Mito.\ \& chloroplasts $\leftarrow$ once free prokaryotes. Evidence: circular DNA, double membranes, prokaryotic ribosomes, binary fission, genetic similarity to bacteria.
  \item \textbf{Cell Theory}: (1) All life = cells. (2) Cell = basic unit. (3) Cells from pre-existing cells.
  \item \textbf{Animal cell}: lysosomes/peroxisomes; \emph{no} wall, \emph{no} central vacuole.
  \item \textbf{Plant cell}: chloroplasts, central vacuole, cell wall with plasmodesmata.
  \item \textbf{SA:V}: smaller cell = higher ratio = more efficient. Movement is \textbf{not} a criterion for life.
\end{itemize}
\end{tcolorbox}

\vspace{4pt}

% ── CHNOPS + Molecular Biology ────────────────────────────────────────────────
\begin{tcolorbox}[colback=part4bg, colframe=sectionorange,
  title={\bfseries CHNOPS, Central Dogma \& Mutations}, arc=3pt, boxrule=1pt]
\begin{multicols}{2}
\textbf{Elements by molecule:}
\begin{itemize}[noitemsep, leftmargin=*]
  \item Carbs \& Lipids: C, H, O
  \item Proteins: C, H, N, O, (S)
  \item Nucleic Acids: C, H, N, O, P
\end{itemize}
\textbf{Chargaff's Rules:} \%A=\%T; \%C=\%G
\columnbreak
\textbf{Central Dogma:}\\
DNA $\to$ pre-mRNA $\to$ mRNA $\to$ Protein\\[4pt]
\textbf{mRNA Processing:}
\begin{itemize}[noitemsep, leftmargin=*]
  \item 5' Cap + 3' Poly-A Tail added
  \item Introns removed (spliceosome/snRNPs)
  \item Exons joined = mature mRNA
\end{itemize}
\end{multicols}
\divider
\textbf{Genetic Code:} 64 codons for 20 amino acids. \textbf{Redundant}, \textbf{non-ambiguous}, \textbf{universal}. Start = AUG; Stop = UAA, UAG, UGA.\\[2pt]
\textbf{Mutations:} \textbf{Silent} = same a.a. (wobble); \textbf{Missense} = different a.a.; \textbf{Nonsense} = premature stop codon.
\end{tcolorbox}

\vspace{4pt}

% ── Membranes, Transport & Energy ────────────────────────────────────────────
\begin{tcolorbox}[colback=part5bg, colframe=sectionteal,
  title={\bfseries Membranes, Transport \& Energy}, arc=3pt, boxrule=1pt]
\begin{multicols}{2}
\textbf{Membrane fluidity (cold):}
\begin{itemize}[noitemsep, leftmargin=*]
  \item More \textbf{unsaturated} tails (kinked $\to$ fluid)
  \item More \textbf{cholesterol} (buffers fluidity)
\end{itemize}
\textbf{Transport types:}
\begin{itemize}[noitemsep, leftmargin=*]
  \item \textbf{Simple diffusion} -- small/nonpolar; no protein
  \item \textbf{Facilitated diffusion} -- specific proteins; high$\to$low; no ATP
  \item \textbf{Active transport} -- pumps; low$\to$high; ATP required
  \item \textbf{Endo/Exocytosis} -- vesicles; bulk transport
\end{itemize}
\columnbreak
\textbf{Tonicity:} Hyper = water out; Hypo = water in; Iso = no net flow\\[4pt]
\textbf{Gibbs Free Energy ($\Delta G$):}
\begin{itemize}[noitemsep, leftmargin=*]
  \item $\Delta G < 0$: Exergonic/Catabolic (spontaneous)
  \item $\Delta G > 0$: Endergonic/Anabolic (non-spontaneous)
  \item Equilibrium: $\Delta G = 0$
\end{itemize}
\textbf{Enzymes:} lower $E_a$; do NOT alter $\Delta G$; specific; not consumed. \textbf{ATP:} energy coupling; hydrolysis releases energy; phosphorylation = increased reactivity.
\end{multicols}
\end{tcolorbox}


\vspace{4pt}

% ── Properties of Water ──────────────────────────────────────────────────────
\begin{tcolorbox}[colback=part2bg, colframe=sectiongreen,
  title={\bfseries Properties of Water \& pH}, arc=3pt, boxrule=1pt]
\footnotesize
\begin{multicols}{2}
\begin{itemize}[noitemsep, leftmargin=*]
  \item Polarity $\to$ hydrogen bonds $\to$ all 6 properties
  \item \textbf{Cohesion}: water-to-water (surface tension, xylem)
  \item \textbf{Adhesion}: water-to-polar surface (capillary action)
  \item \textbf{High specific heat}: buffers temperature (4.18 J/g\textdegree C)
\end{itemize}
\columnbreak
\begin{itemize}[noitemsep, leftmargin=*]
  \item \textbf{High heat of vaporization}: evaporative cooling (sweating)
  \item \textbf{Universal solvent}: polar/ionic substances dissolve
  \item \textbf{Ice floats}: liquid denser than solid; protects aquatic life
  \item \textbf{pH}: $<7$ = acid (more H$^+$); $>7$ = base; 7 = neutral; log scale
\end{itemize}
\end{multicols}
\end{tcolorbox}

\vspace{4pt}

% ── Protein Structure ─────────────────────────────────────────────────────────
\begin{tcolorbox}[colback=partbg, colframe=sectionblue,
  title={\bfseries Protein Structure \& DNA Replication}, arc=3pt, boxrule=1pt]
\footnotesize
\begin{multicols}{2}
\textbf{4 levels:}
\begin{itemize}[noitemsep, leftmargin=*]
  \item \textbf{1°}: aa sequence; peptide bonds (covalent)
  \item \textbf{2°}: $\alpha$-helix / $\beta$-sheet; H-bonds (backbone)
  \item \textbf{3°}: overall 3D shape; R-group interactions (disulfide, ionic, hydrophobic)
  \item \textbf{4°}: multiple polypeptides (e.g., hemoglobin, 4 subunits)
  \item Denaturation = lose 3D shape (heat, pH); primary intact
\end{itemize}
\columnbreak
\textbf{DNA Replication:}
\begin{itemize}[noitemsep, leftmargin=*]
  \item Semi-conservative (one old + one new strand)
  \item Helicase unwinds; primase makes RNA primer; DNA pol III synthesizes 5'$\to$3'
  \item Leading strand: continuous; lagging: Okazaki fragments
  \item DNA pol I removes primers; ligase joins fragments
  \item Proofreading by DNA pol III reduces errors
\end{itemize}
\end{multicols}
\end{tcolorbox}

\vspace{4pt}

% ── Cell Cycle & Genetics ─────────────────────────────────────────────────────
\begin{tcolorbox}[colback=part5bg, colframe=sectionteal,
  title={\bfseries Cell Cycle, Mitosis \& Meiosis}, arc=3pt, boxrule=1pt]
\footnotesize
\begin{multicols}{2}
\begin{itemize}[noitemsep, leftmargin=*]
  \item Interphase: G$_1$ (grow) $\to$ S (replicate DNA) $\to$ G$_2$ (prepare)
  \item Mitosis PMAT: Prophase (condense), Metaphase (align), Anaphase (separate chromatids), Telophase (reform nuclei)
  \item Result: 2 identical diploid cells; for growth/repair
  \item \textbf{Cancer} = checkpoint failure (p53 monitors G$_1$/S)
\end{itemize}
\columnbreak
\begin{itemize}[noitemsep, leftmargin=*]
  \item Meiosis I: homologs separate (crossing over in Prophase I!)
  \item Meiosis II: sister chromatids separate (like mitosis)
  \item Result: 4 haploid genetically unique cells (gametes)
  \item Variation sources: crossing over, independent assortment ($2^{23}$), random fertilization
  \item Nondisjunction $\to$ aneuploidy (e.g., Down syndrome = trisomy 21)
\end{itemize}
\end{multicols}
\end{tcolorbox}

\vspace{4pt}

% ── Genetics ──────────────────────────────────────────────────────────────────
\begin{tcolorbox}[colback=part3bg, colframe=sectionpurple,
  title={\bfseries Mendelian \& Non-Mendelian Genetics}, arc=3pt, boxrule=1pt]
\footnotesize
\begin{multicols}{2}
\begin{itemize}[noitemsep, leftmargin=*]
  \item Genotype = alleles present; phenotype = observable trait
  \item Dominant (capital) masks recessive (lowercase)
  \item \textbf{Law of Segregation:} alleles separate into gametes
  \item \textbf{Law of Ind.\ Assortment:} genes on different chromosomes sort independently
  \item Monohybrid $Aa \times Aa$ $\to$ 3:1 ratio
  \item Dihybrid $AaBb \times AaBb$ $\to$ 9:3:3:1 ratio
\end{itemize}
\columnbreak
\begin{itemize}[noitemsep, leftmargin=*]
  \item \textbf{Incomplete dominance:} heterozygote is intermediate (RW = pink)
  \item \textbf{Codominance:} both alleles expressed (AB blood)
  \item \textbf{Polygenic:} many genes $\to$ continuous variation (height)
  \item \textbf{X-linked recessive:} more common in males
  \item Blood type: $I^A$, $I^B$, $i$; O = universal donor; AB = universal recipient
  \item H-W: $p+q=1$; $p^2+2pq+q^2=1$; use to find carrier frequency
\end{itemize}
\end{multicols}
\end{tcolorbox}

\vspace{4pt}

% ── Evolution & Ecology ───────────────────────────────────────────────────────
\begin{tcolorbox}[colback=part2bg, colframe=sectiongreen,
  title={\bfseries Evolution, Classification \& Ecology}, arc=3pt, boxrule=1pt]
\footnotesize
\begin{multicols}{2}
\textbf{Evolution mechanisms (5):}
\begin{itemize}[noitemsep, leftmargin=*]
  \item Natural selection (non-random, fitness-based)
  \item Genetic drift (random; bottleneck/founder)
  \item Gene flow (migration reduces differences)
  \item Mutation (source of new alleles)
  \item Sexual selection (mate choice)
\end{itemize}
\textbf{Selection types:} directional (shift), stabilizing (middle), disruptive (extremes)
\columnbreak
\textbf{Ecology:}
\begin{itemize}[noitemsep, leftmargin=*]
  \item Symbiosis: mutualism (+/+), commensalism (+/0), parasitism (+/$-$)
  \item 10\% rule: only 10\% energy passes per trophic level
  \item Logistic growth $\to$ carrying capacity K
  \item 3 domains: Bacteria, Archaea, Eukarya
  \item Taxonomy: DKPCOFGS (Dear King Philip\ldots)
  \item Allopatric = geographic isolation $\to$ speciation
  \item H-W: no selection, mutation, drift, flow; random mating
\end{itemize}
\end{multicols}
\end{tcolorbox}
\end{document}